\documentclass[11pt]{article}

% ============================================================
% Packages
% ============================================================
\usepackage[a4paper,margin=1in]{geometry}
\usepackage[T1]{fontenc}
\usepackage{lmodern}
\usepackage{microtype}
\usepackage{amsmath,amssymb,amsthm,mathtools}
\usepackage{bm}
\usepackage{graphicx}
\usepackage{float}
\usepackage{booktabs}
\usepackage{enumitem}
\usepackage{tabularx}
\usepackage{array}
\usepackage{hyperref}
\usepackage[nameinlink,noabbrev]{cleveref}

\hypersetup{
  colorlinks=true,
  linkcolor=blue,
  citecolor=blue,
  urlcolor=blue
}

% ============================================================
% Macros / notation
% ============================================================
\newcommand{\dd}{\mathrm d}
\newcommand{\ii}{\mathrm i}
\newcommand{\ee}{\mathrm e}
\newcommand{\etaQ}{\eta_Q}
\newcommand{\qstar}{q_*}
\newcommand{\estar}{e_*}
\newcommand{\qeff}{q_{\mathrm{eff}}}
\newcommand{\eeff}{e_{\mathrm{eff}}}

\newcommand{\R}{\mathbb R}
\newcommand{\X}{\mathbf X}
\newcommand{\x}{\mathbf x}
\newcommand{\Qgeom}{\mathbf Q_{\geom}}
\newcommand{\pvec}{\mathbf p}
\newcommand{\vct}{\mathbf v}
\newcommand{\nvec}{\mathbf n}
\newcommand{\uvec}{\mathbf u}
\newcommand{\grad}{\nabla}
\newcommand{\lap}{\nabla^2}

\newcommand{\OnePN}{\mathrm{1PN}}
\newcommand{\TwoPN}{\mathrm{2PN}}
\newcommand{\ADM}{\mathrm{ADM}}
\newcommand{\DtN}{\mathrm{DtN}}
\newcommand{\EIH}{\mathrm{EIH}}
\newcommand{\Schw}{\mathrm{Schw}}
\newcommand{\eff}{\mathrm{eff}}
\newcommand{\geom}{\mathrm{geom}}
\newcommand{\supp}{\mathrm{support}}
\newcommand{\raw}{\mathrm{raw}}
\newcommand{\loc}{\mathrm{loc}}

\newcommand{\Wproj}{W}
\newcommand{\Zloc}{Z}
\newcommand{\Proj}[1]{\overline{#1}}

\newcommand{\kadd}{\kappa_{\mathrm{add}}}
\newcommand{\kPV}{\kappa_{\mathrm{PV}}}
\newcommand{\krho}{\kappa_{\rho}}
\newcommand{\betaPN}{\beta_{\mathrm{1PN}}}
\newcommand{\Kvec}{K_{\mathrm{vec}}}
\newcommand{\KEOS}{K_{\mathrm{EOS}}}
\newcommand{\alphaSq}{\alpha^2}

\newcommand{\Exact}{\textbf{Exact}}
\newcommand{\Controlled}{\textbf{Controlled reduction}}
\newcommand{\Protocol}{\textbf{Protocol closure}}
\newcommand{\FullMatch}{\textbf{Full assembly}}
\newcommand{\Ansatz}{\textbf{Ansatz-level consistency}}

\newcommand{\claimclass}[1]{\paragraph{Claim class.}\emph{#1.}}
\newcommand{\TODO}[1]{\textbf{[TODO: #1]}}

\newcolumntype{L}{>{\raggedright\arraybackslash}X}

\newtheorem{claim}{Claim}
\newtheorem{remark}{Remark}

% ============================================================
% Title
% ============================================================
\title{A Controlled Derivation of the Full Conservative Second Post-Newtonian Sector\\
from the Unified 4D Toy Model}
\author{Trevor Norris}
\date{\today}

\begin{document}
\maketitle

\begin{abstract}
We present a controlled derivation of the full conservative two-body second
post-Newtonian (\(\TwoPN\)) sector from the unified \(4\)D toy model.  Carrying
forward the exact projection / Poisson-hook / worldtube reduction logic and the
already-fixed Newtonian+\(\OnePN\) ledger, we close the missing one-body
\(\TwoPN\) response slot by a minimal \(\DtN\)-invariant denominator that
preserves the frozen \(\OnePN\) coefficient and reproduces the isotropic
Schwarzschild test-mass target through \(\TwoPN\).  Minimal Bernoulli
exactification and quadratic local mass scaling then supply the self/static
\(\TwoPN\) seed, with the pure-static isotropic gate fixing
\(\lambda_\rho=\tfrac12\).  The remaining comparable-mass mismatch against the
generic-frame \(\ADM\) target collapses to a finite algebraic residual with a
unique solution for the added cross-sector coefficients \(\{q_i,t_i,s_1\}\).
After assembly, the full reduced two-body Lagrangian Legendre-transforms
exactly to the standard generic-frame conservative \(\ADM\) Hamiltonian through
\(\TwoPN\).  The paper therefore gives a full conservative \(\TwoPN\)
derivation within a declared closure hierarchy, while the appendices preserve
the broader throat-response constructions and the remaining dynamic
\(\DtN\)/PDE observables.
\end{abstract}

% ============================================================
\section{Introduction and scope}
\label{sec:intro}

\subsection{Motivation: from full \texorpdfstring{$1$}{1}PN assembly to full conservative \texorpdfstring{$2$}{2}PN assembly}
\label{sec:intro-motivation}

The earlier papers in this series established two pieces of the story that are now taken as the starting point rather than the destination.
First, the unified $4$D toy model supplies the exact projection, Poisson-hook, and worldtube dictionary from which a brane-facing particle mechanics is meant to descend.
Second, the companion $1$PN assembly paper showed that, once that dictionary is frozen and a short list of closures is declared explicitly, the full conservative two-body $1$PN ledger can be recovered in a referee-auditable way.
The purpose of the present paper is to take the next step without relaxing that bookkeeping discipline:
we ask whether the same program closes at $2$PN.

By ``full conservative $2$PN'' we mean the complete near-zone two-body ledger through order $c^{-4}$.
Concretely, the target includes the free-particle sextic terms, the one-body/self corrections tied to the test-mass limit, the comparable-mass quartic velocity sector at order $G/r$, the quadratic velocity sector at order $G^2/r^2$, and the static cubic term at order $G^3/r^3$.
Our goal is not merely to reproduce one orbit-level consequence or one coordinate-specialized formula.
The goal is to derive a reduced two-body Lagrangian whose full $2$PN content Legendre-transforms exactly to the standard generic-frame $\ADM$ Hamiltonian target.

At $2$PN the problem changes character.
At $1$PN, much of the work was to identify and fix the relevant coefficients sector by sector.
At $2$PN, two additional demands become unavoidable.
First, the already-frozen Newtonian+$\OnePN$ ledger must survive unchanged; a new closure is acceptable only if it repairs the missing $2$PN slot without spoiling the settled $1$PN coefficients.
Second, the constructive throat-response record becomes much larger than the proof itself.
The project now contains tensor-wake decompositions, mouth-port reconstructions, support-minus-closure EFT packaging, minimal throat operators, PDE-side $\DtN$ scaffolds, and Family-1 wall reductions.
All of that structure matters, but putting all of it in the main line would make the first theorem harder to read rather than easier to trust.

This paper therefore has a deliberately narrower purpose than the full $2$PN research record and a stronger purpose than a notebook-level target match.
It is narrower because the main text is organized around the smallest proof chain needed for the conservative $2$PN claim.
It is stronger because, within a stated closure hierarchy, that chain now runs from the carried-forward $4$D reduction framework to the full two-body $2$PN target:
\begin{enumerate}[leftmargin=1.5em]
  \item exact carry-forward of the projection / Poisson-hook / worldtube reduction logic and of the frozen Newtonian+$\OnePN$ ledger,
  \item closure of the missing one-body $\TwoPN$ slot by a minimal $\DtN$-invariant conservative denominator that preserves the settled $1$PN coefficient,
  \item construction of the self/static $\TwoPN$ seed from minimal Bernoulli exactification and quadratic local mass scaling,
  \item a unique comparable-mass residual solve against the standard generic-frame $\ADM$ target, and
  \item exact assembly of the full conservative two-body $\TwoPN$ Lagrangian within the declared hierarchy.
\end{enumerate}

A second purpose of the paper is bookkeeping clarity.
In the $4$D framework, projection is not a harmless relabeling; it naturally produces exchange, response, and closure terms.
That makes the separation between exact carry-forward facts, reduced-regime statements, protocol-dependent closures, and final assembled equalities especially important at $2$PN.
The reader should be able to see, line by line, which parts of the derivation are inherited exactly, which parts depend on a declared low-frequency or small-body hierarchy, and which parts represent the minimal new closure input required to finish the conservative proof.
That separation is part of the result.

This paper is also downstream of the corrected Paper~7 / Paper~8 matter--gauge
ontology. Electric-charge sign is the fixed topological puncture orientation
\(\etaQ=\pm1\), the microscopic branch coupling is \(\qstar=\etaQ\estar\) with
\(\estar>0\), and the canonically normalized brane charge is
\(\qeff=\qstar/\sqrt{Z_{\rm int}}=\etaQ\eeff\), where
\(Z_{\rm int}\equiv\int Z(w)\,\dd w\), with
\(\eeff=\estar/\sqrt{Z_{\rm int}}\). Quantized circulation belongs to the
magnetic/vortical sector rather than the electric-charge dictionary, and the
zero-mode Maxwell reduction
\((A_w\approx0,\ \partial_w A_\mu\approx0,\ J^w\approx0,\ F_{\mu w}\approx0)\)
is used here only as a controlled far-field brane limit. Separately, the
scalar Newtonian mass-dressing coefficient historically denoted \(q\) in pre-4D
gravity bookkeeping is not electric charge; in the present downstream notation
it is written \(\krho\), fixed by Newtonian matching to \(\krho=1\).

\subsection{What this paper claims}
\label{sec:intro-claims}

The main claims are the following.
\begin{enumerate}[leftmargin=1.5em]
  \item The exact $4$D projection / Poisson-hook / worldtube framework from the earlier papers can be carried forward unchanged as the reduction backbone of the present $\TwoPN$ derivation.
  \item Within a minimal $\DtN$-invariant conservative closure, the missing one-body $\TwoPN$ response slot is fixed by a denominator that preserves the frozen $\OnePN$ coefficient and exactly reproduces the isotropic Schwarzschild test-mass target through $\TwoPN$.
  \item Minimal Bernoulli exactification together with quadratic local mass scaling yields the self/static $\TwoPN$ seed used in the assembly, with the pure-static isotropic gate selecting $\lambda_\rho=1/2$.
  \item Once the frozen Newtonian+$\OnePN$ block and the one-body/self/static $\TwoPN$ seed are fixed, the remaining comparable-mass residual is solved uniquely in a compact invariant basis, giving an added cross block whose Legendre transform matches the generic-frame $\ADM$ target exactly.
  \item Combining these ingredients yields the full conservative two-body $\TwoPN$ Lagrangian within the stated closure hierarchy, with the status of each ingredient recorded explicitly as exact, controlled, protocol-dependent, or full assembly.
  \item A constructive throat-response interpretation of the new $\TwoPN$ cross content already exists: the carried-forward dipole wake survives from $1$PN, the genuinely new conservative response lives in a small $P_0\oplus P_2$ layer together with a geometry-closure channel, and the appendix program records the corresponding $\DtN$, PDE, and Family-1 realizations.
\end{enumerate}

Stated differently, the point of the paper is not that the $4$D model can be adjusted until a $2$PN-looking polynomial appears.
The point is that, once the derivation is organized sector by sector, the full conservative two-body $2$PN structure can be recovered with a short list of named and auditable ingredients.
The constructive throat-response material then becomes an interpretation and continuation of that proof rather than a substitute for it.

\subsection{What this paper does not claim}
\label{sec:intro-nonclaims}

The negative side of the scope is equally important.

First, this paper does \emph{not} claim an assumption-free theorem that begins with a fully solved moving-throat bulk background and ends automatically with the complete point-particle $\TwoPN$ dynamics.
As at $1$PN, the particle limit used here is a controlled reduction inside a declared closure hierarchy.
The proof in the main text is therefore a reduced-theory derivation, not a finished collective-coordinate theorem for the full bulk PDE.

Second, this paper does \emph{not} claim a first-principles dynamical derivation of every new $\TwoPN$ response datum.
The one-body conservative denominator is fixed here within a minimal $\DtN$-invariant packaging, and the final comparable-mass cross block is fixed here by an exact residual solve once the lower-order ledger and the declared seed are frozen.
But the detailed dynamical origin of the pole scales and of the pure geometry-closure channel belongs to the appendix-side throat program and to later work.
Those quantities should still be read as inner-throat observables, not as already-completed PDE theorems.

Third, this paper is restricted to the conservative sector.
It does not address radiation reaction, dissipative leakage, spin couplings, or genuinely non-adiabatic internal response.
It also does not attempt a full many-body dynamical completion beyond the conservative local-potential consequences already implied by the static scaling and throat-response packaging.

Fourth, the appendices should not be read as optional leftovers, but neither should they be confused with the load-bearing proof of the main theorem.
They preserve the constructive record of how the $2$PN cross content can be packaged in mouth-port, support-minus-closure, $\DtN$, and Family-1 language.
That record matters for the next paper, but the present paper's main claim is the conservative $2$PN proof itself.

These non-claims are not qualifications to be hidden.
They are part of the correct interpretation of the result:
this paper establishes a full conservative two-body $2$PN derivation \emph{within a stated closure hierarchy}.

\subsection{Claim taxonomy}
\label{sec:intro-taxonomy}

To keep the argument referee-safe, we use the same claim taxonomy as in the $\OnePN$ assembly paper.
\begin{description}[leftmargin=1.8em,style=nextline]
  \item[\Exact] identities that follow directly from the declared reduction framework, symmetry bookkeeping, exact algebraic transformations, or exact residual matching once the relevant frozen ledger and closures have been stated.
  \item[\Controlled] steps that require an explicit regime assumption, such as the carried-forward small-body hierarchy, near-zone/low-frequency reduction, quasi-static bookkeeping, or the reduced $\DtN$ packaging used to close the one-body slot.
  \item[\Protocol] coefficients that are fixed only within a declared response closure rather than by a fully solved moving-throat derivation.
  In the present paper, the central example is the minimal conservative packaging used to close the one-body response denominator.
  \item[\FullMatch] the final exact equality between the assembled reduced two-body theory and the standard generic-frame $\ADM$ $\TwoPN$ target once all earlier ingredients are inserted.
\end{description}

This taxonomy is not cosmetic.
At $2$PN it prevents three logically distinct moves from being blurred together:
importing exact $4$D structure,
declaring the smallest new closures needed to finish the conservative ledger,
and performing the exact algebraic solve and final target match once those closures are frozen.

\subsection{Roadmap}
\label{sec:intro-roadmap}

Section~\ref{sec:executive} gives a front-end ledger of the main coefficients, proof steps, and open observables.
Section~\ref{sec:carry-forward} then records the minimal carry-forward dictionary from the earlier papers: the imported projection logic, the frozen Newtonian+$\OnePN$ two-body ledger, and the small closure hierarchy used in the present $\TwoPN$ assembly.
Section~\ref{sec:one-body} closes the missing one-body $\TwoPN$ response slot and compares the resulting denominator directly to the isotropic test-mass target.
Section~\ref{sec:self-static} derives the self/static seed used in the assembly and states the finite-size gate that prevents over-claiming the static sector.
Section~\ref{sec:adm-lift} performs the comparable-mass residual solve against the standard $\ADM$ target, and Section~\ref{sec:full-assembly} assembles the full conservative two-body $\TwoPN$ Lagrangian and states the exact match.
Sections~\ref{sec:status-ledger}--\ref{sec:conclusion} then record the fixed-vs-open status ledger, the verification logic tied to the symbolic harnesses, and the concluding interpretation.

The appendices collect the larger constructive record that is important but not required for the main proof:
one-body $\DtN$ algebra, static-scaling details, the full $\ADM$ residual solve, tensor-wake and mouth-port decompositions, support-minus-closure EFT packaging, minimal throat operators, PDE/$\DtN$ scaffolding, Family-1 soft-wall and traction-balance reductions, geometry-breathing closures, and the final coupled response module.
They are used here to keep the main theorem readable without discarding the broader $2$PN program that motivated it.


% ============================================================
\section{Executive summary of results}
\label{sec:executive}

This paper is the point at which the $4$D program becomes a full conservative second post-Newtonian derivation rather than a partial closure story. The earlier $4$D papers supply the exact projection dictionary, exact balance laws, near-zone Poisson hook, and worldtube reduction needed to define the brane-facing Newtonian sector, while the companion $1$PN assembly paper freezes the complete Newtonian+$\OnePN$ two-body ledger. The present paper then carries those ingredients through the smallest additional chain needed to close the conservative $2$PN problem: a $\DtN$-invariant one-body denominator, a self/static seed from Bernoulli exactification and local mass scaling, and a unique comparable-mass lift of the remaining residual against the standard generic-frame $\ADM$ target.

\paragraph{Headline result.}
Within the stated closure hierarchy, the toy model reproduces the full conservative two-body $\TwoPN$ ledger through order $c^{-4}$. The proof is not a single coefficient fit but a short triangulated chain. The carried Newtonian+$\OnePN$ structure is held fixed, the missing one-body $\TwoPN$ slot is closed in a way that preserves the settled $\OnePN$ coefficient, the pure-static isotropic gate selects $\lambda_\rho=1/2$, and the remaining comparable-mass mismatch then collapses to a finite algebraic residual with a unique solution. The final reduced two-body Lagrangian Legendre-transforms exactly to the standard generic-frame $\ADM$ Hamiltonian through $\TwoPN$.

Table~\ref{tab:executive-ledger} lists the quantities proved in the paper and the logic by which each is obtained.

\begin{table}[H]
  \centering
  \caption{Headline results established in the $\TwoPN$ paper. The table is deliberately front-loaded: it is meant to function as a referee's map of what is carried exactly, what is closed within the declared reduction hierarchy, and what is only claimed after full assembly.}
  \label{tab:executive-ledger}
  \small
  \setlength{\tabcolsep}{4pt}
  \renewcommand{\arraystretch}{1.18}
  \begin{tabularx}{\linewidth}{@{}>{\raggedright\arraybackslash}p{0.23\linewidth} >{\raggedright\arraybackslash}p{0.26\linewidth} >{\raggedright\arraybackslash}p{0.31\linewidth} >{\raggedright\arraybackslash}p{0.14\linewidth}@{}}
    \toprule
    Quantity & Value & Source of derivation & Claim class \\
    \midrule
    Reduction backbone & exact projection / Poisson-hook / worldtube dictionary & carried forward from the earlier $4$D papers & \Exact / \Controlled \\
    Frozen lower-order ledger & full Newtonian+$\OnePN$ two-body block with exact $\EIH$ match & carried forward from the $1$PN assembly paper & \Exact / \FullMatch \\
    One-body denominator & $D_{\eff}(u)=C_s(u)\bigl[1+\tfrac{32768}{3249}(G(u)-1)^2\bigr]$ & minimal $\DtN$-invariant closure preserving the frozen $1$PN slot & \Controlled / \Protocol \\
    Test-mass $\TwoPN$ target & $D_{\eff}(u)=1-4u+8u^2+O(u^3)$ & exact isotropic Schwarzschild match through $\TwoPN$ & \FullMatch \\
    Static scaling coefficient & $\lambda_\rho=\tfrac12$ & pure-static isotropic target gate & \Controlled \\
    Self/static $\TwoPN$ seed & minimal Bernoulli exactification + quadratic local mass scaling & reduced one-body/self/static sector & \Controlled \\
    Comparable-mass residual coefficients & unique solution for $\{q_i,t_i,s_1\}$ & compact invariant-basis solve against the generic-frame $\ADM$ target & \Exact \\
    Full conservative two-body $\TwoPN$ block & exact generic-frame $\ADM$ match after assembly & full two-body reduction and Legendre check & \FullMatch \\
    \bottomrule
  \end{tabularx}
\end{table}

\paragraph{How the ledger closes.}
There are three nested stories in the proof. The first is the one-body slot. Writing the conservative denominator as
\[
D_{\eff}(u)=C_s(u)\left[1+\alpha\,(G-1)+\beta\,(G-1)^2\right],
\]
preservation of the already-frozen $\OnePN$ coefficient forces the linear correction to vanish,
\[
\alpha=0.
\]
The first admissible new term is therefore quadratic in the geometry invariant. With the carried throat slope
\[
G(u)=1+\frac{57}{64}u+O(u^2),
\]
matching the exact isotropic Schwarzschild test-mass target through $\TwoPN$ fixes
\[
D_{\eff}(u)=C_s(u)\left[1+\frac{32768}{3249}(G(u)-1)^2\right]
=1-4u+8u^2+O(u^3).
\]
So the missing one-body conservative $2$PN slot is closed without disturbing the settled $1$PN ledger.

The second story is the self/static seed. Minimal Bernoulli exactification supplies the free-particle sextic term and the self-sector $G/r$ velocity dressing, while quadratic local mass scaling supplies the static cubic hierarchy. The pure-static isotropic gate selects
\[
\lambda_\rho=\frac12,
\]
so the reduced self/static $\TwoPN$ seed entering the assembly is
\begin{align}
L_{2,\text{self+static}}=
&\frac{m_A v_A^6+m_B v_B^6}{16}
+\frac{7Gm_A m_B}{8r}(v_A^4+v_B^4)
+\frac{2G^2 m_A m_B}{r^2}(m_B v_A^2+m_A v_B^2) \notag\\
&+\frac{G^3m_A m_B(m_A^2+m_B^2)}{4r^3},
\end{align}
with the overall $c^{-4}$ factor understood.

The third story is the comparable-mass lift. The key simplification is that for
\[
L=L_0+\epsilon L_1+\epsilon^2 L_2,
\qquad \epsilon\equiv c^{-2},
\]
with quadratic $L_0$, the perturbative Legendre transform through $\TwoPN$ collapses to
\[
H_1=-L_1(v_0),
\qquad
H_2=-L_2(v_0)+\frac12 A_0^T M^{-1}A_0,
\]
where $v_0=p/m$ and $A_0=(\partial L_1/\partial v)|_{v=v_0}$. Because the Newtonian+$\OnePN$ block is already frozen, any genuinely new $\TwoPN$ Lagrangian term enters the Hamiltonian only through the explicit $-L_{2,\text{new}}(v_0)$ piece. The remaining mismatch against the standard generic-frame $\ADM$ target is therefore a finite algebraic residual in a compact invariant basis, and that residual solves uniquely.

\paragraph{Solved comparable-mass coefficients.}
The exact residual solve yields
\[
q_1=-\frac74,\quad
q_2=-\frac14,\quad
q_3=\frac{11}{8},\quad
q_4=\frac14,\quad
q_5=-\frac58,\quad
q_6=\frac32,\quad
q_7=\frac38,
\]
\[
t_1=0,\quad
 t_2=\frac{11}{8},\quad
 t_3=-\frac{15}{4},\quad
 t_4=0,\quad
 t_5=0,\quad
 t_6=\frac{15}{8},
\qquad
s_1=\frac54.
\]
These are the coefficients of the unique added comparable-mass cross block required to upgrade the self/static candidate to the full $\TwoPN$ target.

\paragraph{Normalization warning.}
In the executive formulas below, $L_{2,\text{full}}$ denotes the coefficient of $c^{-4}$ in the total two-body Lagrangian. The Newtonian and $\OnePN$ blocks are carried forward unchanged from the earlier paper; only the new $\TwoPN$ block is displayed explicitly.

\paragraph{Final assembled result.}
Using the shorthand
\[
v_{AB}=\vct_A\!\cdot\!\vct_B,
\qquad
v_{An}=\vct_A\!\cdot\!\nvec,
\qquad
v_{Bn}=\vct_B\!\cdot\!\nvec,
\]
the new $\TwoPN$ block inserted into the carried Newtonian+$\OnePN$ ledger is
\begin{align}
L_{2,\text{full}}=
&\frac{m_A v_A^6+m_B v_B^6}{16}
+\frac{Gm_A m_B}{r}\Big[
\frac78(v_A^4+v_B^4)
-\frac74\,v_{AB}(v_A^2+v_B^2)
-\frac14\,v_{An}v_{Bn}(v_A^2+v_B^2) \notag\\
&\hspace{7em}
+\frac{11}{8}\,v_A^2v_B^2
+\frac14\,v_{AB}^2
-\frac58\,(v_A^2v_{Bn}^2+v_B^2v_{An}^2)
+\frac32\,v_{AB}v_{An}v_{Bn}
+\frac38\,v_{An}^2v_{Bn}^2
\Big] \notag\\
&+\frac{G^2m_A m_B}{r^2}\Big[
\Big(2m_B+\frac{11}{8}m_A\Big)v_A^2
+\Big(2m_A+\frac{11}{8}m_B\Big)v_B^2
-\frac{15}{4}(m_A+m_B)v_{AB} \notag\\
&\hspace{7em}
+\frac{15}{8}(m_A v_{An}^2+m_B v_{Bn}^2)
\Big]
+\frac{G^3m_A m_B}{4r^3}(m_A^2+5m_A m_B+m_B^2).
\label{eq:exec-2pn-block}
\end{align}
The claim of the paper is that this is not merely $\ADM$-like: once the carried and newly solved ingredients are inserted, the Legendre transform of the full reduced two-body theory agrees identically with the standard generic-frame conservative $\ADM$ target through $\TwoPN$.

\paragraph{How to read the claim classes.}
The table mixes four logically different kinds of statements. The carried reduction backbone and the exact residual solve are {\Exact} statements once the relevant ledger and closures have been declared. The one-body denominator and the self/static seed are {\Controlled} statements because they depend on an explicit reduced-regime hierarchy, and the one-body denominator is also a {\Protocol} statement because it is fixed within a declared conservative $\DtN$ packaging rather than by a fully solved moving-throat PDE theorem. The final equality to the generic-frame $\ADM$ target is the central {\FullMatch} statement of the paper.

\paragraph{What is fixed and what remains open.}
Within the scope of the present proof, the values
\[
\lambda_\rho=\frac12,
\qquad
D_{\eff}(u)=C_s(u)\left[1+\frac{32768}{3249}(G(u)-1)^2\right],
\qquad
\{q_i,t_i,s_1\}\ \text{as listed above},
\]
should be read as derived outputs rather than free knobs. The appendix-side constructive program also fixes a nontrivial zero-frequency throat ledger,
\[
R_{1\perp}=\frac72,
\qquad
R_{10}=4,
\qquad
R_0=R_{20}=R_{21}=R_{22}=1,
\]
\[
J=\left(\frac{4}{\sqrt5},\frac54,0,0,0,0\right),
\qquad
\Delta_{\geom}=\frac{281}{80},
\]
where \(J\) is an appendix-side throat-response source-weight vector rather than
an electromagnetic current. This ledger packages the new conservative cross
content into a carried odd dipole backbone, a new even $P_0\oplus P_2$ support
layer, and a pure-$U$ geometry-closure channel. What remains open are the
dynamical pole scales
\[
\Omega_{1\perp},\ \Omega_{10},\ \Omega_0,\ \Omega_{20},\ \Omega_{21},\ \Omega_{22},\ \Omega_g,
\]
and, more broadly, a first-principles moving-throat PDE derivation of the same low-frequency response data. The paper therefore claims a full conservative two-body $2$PN derivation within a stated closure hierarchy, not an assumption-free theorem of the fully solved bulk dynamics.

\paragraph{Appendix map.}
The appendices preserve the larger constructive record that is important for the program but not load-bearing for the proof. They collect the one-body $\DtN$ algebra, the static-scaling and finite-size details, the full $\ADM$ residual solve, tensor-wake and mouth-port decompositions, support-minus-closure EFT packaging, minimal irreducible-channel throat operators, $\DtN$/PDE scaffolding, Family-1 soft-wall and traction-balance reductions, geometry-breathing closures, and the final coupled response module. The main text proves the conservative $2$PN theorem; the appendices document how that theorem is embedded in the broader throat-response program.

% ============================================================
\section{Minimal carry-forward dictionary and frozen \texorpdfstring{$1$}{1}PN ledger}
\label{sec:carry-forward}

This paper does not reopen the derivation of the Newtonian sector or of the full
conservative two-body \(\OnePN\) ledger. Its job is narrower and more rigid: to state,
in one place, exactly what has already been fixed, exactly what notation will be used
below, and exactly which new ingredients a genuine \(\TwoPN\) proof is still allowed to
introduce. The present section therefore plays a bookkeeping role rather than a
constructive one. It freezes the reduction backbone inherited from the earlier \(4\)D
papers, records the already-settled Newtonian+\(\OnePN\) two-body dynamics, and makes
explicit the small closure hierarchy that will be used in the rest of the paper.

The reason for doing this early is simple. At \(2\)PN one can easily blur together
three logically different moves: importing exact or previously-validated lower-order
structure, declaring the minimal new closures needed to repair the missing \(2\)PN
slots, and performing the exact algebraic assembly once those closures are frozen.
The point of the present paper is to keep those moves separate. In particular, the
\(\TwoPN\) calculation is not permitted to retune the already-proved \(\OnePN\) answer.
Any acceptable \(\TwoPN\) completion must reduce to the same Newtonian and \(\OnePN\)
ledger when the \(c^{-4}\) block is switched off.

The notation carried forward here is also downstream of the corrected
matter--gauge stack. Electric charge does not enter the main conservative
\(\TwoPN\) proof; when it is named explicitly it follows the corrected
\((\etaQ,\qstar,\qeff)\) ontology above. The residual-basis coefficients
\(q_i\) used later are simply \(\ADM\)-lift coefficients, and the historical
gravity-side scalar coefficient once written as bare \(q\) is written here as
\(\krho\), fixed by Newtonian matching to \(\krho=1\).

\subsection{What is imported without rederivation}
\label{sec:carry-forward-import}

Only a narrow subset of the earlier \(4\)D program is imported into the present paper,
but that subset is structural and should be kept conceptually visible.

First, we carry forward the distinction between \emph{projection} and \emph{reduction}.
Projection defines what a brane observer measures: it is the operational map from bulk
fields to brane-facing observables. Reduction is a second step, performed only after a
regime has been declared, in which those projected observables are approximated by an
effective particle description. The present \(\TwoPN\) paper uses exactly the same logic
as the \(\OnePN\) assembly paper: projection identities are inherited from the parent
\(4\)D description, whereas the particle mechanics emerges only after the same carried
near-zone and small-body hierarchy has been adopted.

Second, the exact reduction backbone imported from the earlier papers is unchanged.
The bulk-to-brane projection map, the exact projected continuity law, and the exact
longitudinal identity are treated as settled structural inputs. In the quasi-static
near-zone regime, the same Poisson-hook logic is then carried forward to identify the
brane-facing Newtonian potential. Nothing in the present paper modifies that route.
The \(\TwoPN\) calculation begins only \emph{after} those lower-level identities have
already done their work.

Third, we carry forward the coherent-defect / small-body worldtube reduction that turns
a localized object moving in a smooth external field into an effective particle source.
As at \(\OnePN\), this is a controlled reduction rather than an exact theorem of a fully
solved moving-throat bulk configuration. Its purpose here is not to be rederived but to
fix the meaning of the particle variables \(m_A\), \(m_B\), \(\mathbf x_A\), \(\mathbf x_B\),
\(\vct_A\), and \(\vct_B\) that appear throughout the rest of the paper.

Finally, the full conservative two-body \(\OnePN\) result is imported as frozen data.
That means more than importing a few coefficients. It means that the already-assembled
Newtonian+\(\OnePN\) Lagrangian, together with its exact Einstein--Infeld--Hoffmann
match, is treated as part of the fixed backbone of the \(\TwoPN\) derivation. The
present paper is allowed to add the order-\(c^{-4}\) block; it is not allowed to revise
what has already been proven at order \(c^{-2}\).

\claimclass{\Exact{} for the projection dictionary, exact projected continuity, and
exact longitudinal identity inherited from the earlier \(4\)D papers. \Controlled{}
for the quasi-static Poisson hook and the coherent-defect / small-body worldtube
reduction that turns those exact structures into an effective particle mechanics.}

\subsection{Minimal dictionary of variables and invariants}
\label{sec:carry-forward-dictionary}

We work throughout with two bodies \(A\) and \(B\) in a generic frame. Their positions,
velocities, and masses are denoted
\[
\mathbf x_A,\quad \mathbf x_B,
\qquad
\vct_A=\dot{\mathbf x}_A,
\quad
\vct_B=\dot{\mathbf x}_B,
\qquad
m_A,\quad m_B.
\]
The relative separation vector, its magnitude, and the corresponding unit direction are
\begin{equation}
\mathbf r\equiv \mathbf x_A-\mathbf x_B,
\qquad
r\equiv |\mathbf r|,
\qquad
\nvec\equiv \frac{\mathbf r}{r}.
\label{eq:carry-forward-rn}
\end{equation}
All main-text formulas are written in this generic frame; no center-of-mass
specialization is used in the \(\TwoPN\) residual solve.

The velocity invariants that appear repeatedly are
\begin{equation}
v_A^2\equiv \vct_A\!\cdot\!\vct_A,
\qquad
v_B^2\equiv \vct_B\!\cdot\!\vct_B,
\qquad
v_{AB}\equiv \vct_A\!\cdot\!\vct_B,
\label{eq:carry-forward-vinvariants-1}
\end{equation}
with longitudinal projections
\begin{equation}
v_{An}\equiv \vct_A\!\cdot\!\nvec,
\qquad
v_{Bn}\equiv \vct_B\!\cdot\!\nvec.
\label{eq:carry-forward-vinvariants-2}
\end{equation}
These are the only kinematic invariants needed in the main text. Every comparable-mass
\(\TwoPN\) cross monomial used later can be written in terms of
\(v_A^2\), \(v_B^2\), \(v_{AB}\), \(v_{An}\), and \(v_{Bn}\).

It is also convenient to define the local Newtonian pair potentials
\begin{equation}
U_A\equiv \frac{Gm_B}{r},
\qquad
U_B\equiv \frac{Gm_A}{r}.
\label{eq:carry-forward-UAUB}
\end{equation}
In the strict one-body reduction, where a heavy source of mass \(M\) is treated as fixed,
we write
\begin{equation}
U\equiv \frac{GM}{r},
\qquad
u\equiv \frac{U}{c^2},
\label{eq:carry-forward-u}
\end{equation}
so that the one-body denominator problem can be expressed as a weak-field series in the
single dimensionless variable \(u\).

The rest of the paper uses one additional bookkeeping convention repeatedly. The
conservative two-body Lagrangian is expanded as
\begin{equation}
L_{\rm cons}=L_0+\frac{1}{c^2}L_1+\frac{1}{c^4}L_2+O(c^{-6}),
\label{eq:carry-forward-pn-expansion}
\end{equation}
where \(L_0\) is Newtonian, \(L_1\) is the full frozen \(\OnePN\) block, and \(L_2\) is
what the present paper must determine.

\subsection{Frozen Newtonian + \texorpdfstring{$1$}{1}PN ledger}
\label{sec:carry-forward-1pn-ledger}

The Newtonian block is
\begin{equation}
L_0=
\frac12 m_A v_A^2
+\frac12 m_B v_B^2
+\frac{Gm_A m_B}{r}.
\label{eq:carry-forward-L0}
\end{equation}
Nothing in the present paper alters \cref{eq:carry-forward-L0}. All genuinely new
content begins at \(L_2\).

The full carried \(\OnePN\) block is likewise frozen:
\begin{align}
L_1={}&
\frac18\bigl(m_A v_A^4+m_B v_B^4\bigr)
+\frac{Gm_A m_B}{r}
\left[
\frac32\bigl(v_A^2+v_B^2\bigr)
-\frac72\,v_{AB}
-\frac12\,v_{An}v_{Bn}
\right]
\nonumber\\
&
-\frac{G^2m_A m_B(m_A+m_B)}{2r^2}.
\label{eq:carry-forward-L1}
\end{align}
By construction, \cref{eq:carry-forward-L1} is not a new ansatz introduced for the
present paper. It is the already-proved outcome of the \(\OnePN\) assembly and therefore
the fixed lower-order input to every \(\TwoPN\) statement below.

It is useful to record explicitly which lower-order facts are now being treated as
settled.
Newtonian matching fixed the scalar mass-dressing coefficient to
\(\krho=1\) (historically \(q=1\) in earlier gravity-side bookkeeping).
Weak-field optical consistency fixed the barotropic exponent to
\(n=5\).
Wave-supported free kinematics fixed the quartic kinetic coefficient to \(1/8\).
The scalar/self sector fixed the pair coefficient of \(v_A^2+v_B^2\) to \(+3/2\).
The carried vector-wake sector fixed the cross coefficients of \(v_{AB}\) and
\(v_{An}v_{Bn}\) to \(-7/2\) and \(-1/2\), respectively.
The static nonlinear \(\OnePN\) term fixed the coefficient of
\(-G^2m_A m_B(m_A+m_B)/(2r^2)\).
On the response side, the topology ledger fixed
\(\kappa_{\rm add}=1/2\), the declared adiabatic breathing closure fixed
\(\kappa_{\rm PV}=3/2\), and therefore the already-frozen self-sector/precession ledger is
\begin{equation}
\beta_{\OnePN}=1+\frac12+\frac32=3.
\label{eq:carry-forward-beta1pn}
\end{equation}
These values are no longer free parameters. They are carry-forward data.

Equally important is what this means operationally. In the remainder of the paper the
\(\TwoPN\) construction is not allowed to shift any coefficient in
\cref{eq:carry-forward-L1}, nor is it allowed to reopen the lower-order fitting logic
that produced them. The one-body denominator closure must preserve the frozen
\(\OnePN\) slot. The self/static \(\TwoPN\) seed must reduce to the already-fixed static
regression at lower order. And the comparable-mass residual solve is performed with
\(L_0\) and \(L_1\) held exactly fixed from the start.

\claimclass{\Controlled{} for the carried small-body reduction and for the lower-order
fixing of \(\krho=1\) and \(n=5\). \Protocol{} for the already-declared adiabatic response
input \(\kappa_{\rm PV}=3/2\), and therefore for the carried value
\(\beta_{\OnePN}=3\). \FullMatch{} for the statement that
\(L_0+c^{-2}L_1\) is already exactly the standard conservative two-body \(\EIH\)
Lagrangian.}

\subsection{Closure hierarchy used in the present \texorpdfstring{$2$}{2}PN paper}
\label{sec:carry-forward-closure-hierarchy}

With the lower-order ledger frozen, the rest of the paper adds only the smallest set of
new ingredients needed to complete the conservative \(\TwoPN\) proof.

First, the one-body sector introduces a low-frequency conservative
\(\DtN\)-invariant denominator packaging for the missing test-mass \(\TwoPN\) response
slot. Its job is tightly constrained: it must repair the \(\TwoPN\) one-body target
without changing the already-fixed \(\OnePN\) coefficient.

Second, the self sector is extended by a minimal Bernoulli exactification. This is used
only to generate the conservative one-body/self \(\TwoPN\) seed required by the main
assembly. The paper does not claim that this exhausts the full moving-throat response
problem.

Third, the static sector is extended by quadratic local mass scaling. In the main text,
only the load-bearing consequence is needed: the pure-static isotropic gate selects
\begin{equation}
\lambda_\rho=\frac12.
\label{eq:carry-forward-lambda-rho}
\end{equation}
The induced higher-body static hierarchy and related finite-size caveats are recorded in
appendices rather than in the proof line.

Fourth, once the frozen lower-order ledger and the declared one-body/self/static seed are
inserted, the remaining comparable-mass problem is treated as an exact algebraic residual
solve against the standard generic-frame \(\ADM\) target. This is the step that closes
all genuinely new conservative two-body cross content at \(\TwoPN\).

Everything else belongs to constructive interpretation rather than to the minimal proof.
The tensor-wake rebuild, mouth-port decomposition, support-minus-closure EFT packaging,
minimal irreducible-channel throat operators, PDE-side \(\DtN\) scaffolding, Family-1
soft-wall reductions, and the final coupled throat-response module are all preserved in
the appendices. They are important for understanding \emph{why} the solved \(\TwoPN\)
answer admits a compact throat-response interpretation, but they are not required to
state or verify the main conservative theorem.

A useful way to read the rest of the paper is therefore triangular. The bottom layer is
the exact and already-carried Newtonian+\(\OnePN\) backbone. On top of that sits the
smallest new closure stack needed to define the \(\TwoPN\) seed. Once those inputs are
frozen, the final assembly and \(\ADM\) comparison are exact algebra. That is the sense
in which the paper claims a full conservative \(\TwoPN\) derivation within a stated
closure hierarchy, rather than a first-principles solution of the fully dynamical inner
throat.

% ============================================================
\section{One-body \texorpdfstring{$2$}{2}PN closure and the test-mass target}
\label{sec:one-body}

The first genuinely new ingredient at \(\TwoPN\) already appears in the strict
one-body limit.  Once the Newtonian+\(\OnePN\) ledger has been frozen, the problem is
no longer to rediscover the lower-order coefficients but to repair one specific missing
response slot without disturbing anything that is already fixed.  In the present paper
that slot is the order-\(u^2\) correction to the effective one-body velocity denominator,
with
\[
u\equiv \frac{U}{c^2},
\qquad
U\equiv \frac{GM}{r},
\]
for a light body moving in the field of a fixed heavy source of mass \(M\).

This section isolates that one-body problem and closes it in the smallest possible way.
The exact isotropic Schwarzschild test-mass target tells us what the completed one-body
ledger must look like.  The low-frequency mouth / \(\DtN\) data tell us which invariant
building blocks are naturally available on the toy-model side.  The answer is then fixed
by one simple requirement: any new \(\TwoPN\) correction must preserve the already-frozen
\(\OnePN\) coefficient.

\subsection{Exact isotropic Schwarzschild test-mass target through \texorpdfstring{$2$}{2}PN}
\label{sec:one-body-target}

In isotropic coordinates, the exact test-mass Lagrangian for a particle of mass \(m\)
in the Schwarzschild field of a fixed source \(M\) is
\begin{equation}
\frac{L_{\Schw}}{m}
=-c^2\sqrt{\left(\frac{1-u/2}{1+u/2}\right)^2-
(1+u/2)^4\frac{v^2}{c^2}}.
\label{eq:onebody-schwarzschild-exact}
\end{equation}
Expanding \cref{eq:onebody-schwarzschild-exact} through total post-Newtonian order
\(c^{-4}\) gives
\begin{align}
\frac{L_{\Schw}}{m}
={}&-c^2
+\frac12 v^2+U
+\frac{1}{c^2}\left(\frac18 v^4+\frac32 Uv^2-\frac12 U^2\right)
\nonumber\\
&
+\frac{1}{c^4}\left(\frac1{16}v^6+\frac78 Uv^4+2U^2v^2+\frac14 U^3\right)
+O(c^{-6}).
\label{eq:onebody-schwarzschild-2pn}
\end{align}
Equation~\eqref{eq:onebody-schwarzschild-2pn} is the exact one-body target that the
carried toy-model reduction must reproduce within the stated closure hierarchy.

For the present section, one feature of \cref{eq:onebody-schwarzschild-2pn} is the key
point.  In the carried optical/self packaging, the free \(v^6\) term and the \(Uv^4\)
term are already produced by the minimal Bernoulli exactification, while the static
\(U^3\) term belongs to the static sector treated in the next section.  The genuinely
missing one-body response datum is therefore the \(U^2v^2\) coefficient.  The exact
Schwarzschild target requires that coefficient to be
\begin{equation}
2.
\label{eq:onebody-target-u2v2}
\end{equation}

A convenient way to see the resulting denominator target is to write the reduced one-body
kinetic packaging in the form
\begin{equation}
L_{\mathrm{red},1\text{-}\mathrm{body}}
=-m c^2(1-u)
\sqrt{1-\frac{v^2/c^2}{D(u)}},
\qquad
D(u)=1-4u+d_2u^2+O(u^3).
\label{eq:onebody-reduced-denominator}
\end{equation}
Expanding \cref{eq:onebody-reduced-denominator} gives
\begin{align}
\frac{L_{\mathrm{red},1\text{-}\mathrm{body}}}{m}
={}&-c^2+\frac12 v^2+U
+\frac{1}{c^2}\left(\frac18 v^4+\frac32 Uv^2\right)
\nonumber\\
&
+\frac{1}{c^4}\left(\frac1{16}v^6+\frac78 Uv^4+
\left(6-\frac{d_2}{2}\right)U^2v^2\right)
+O(c^{-6}).
\label{eq:onebody-reduced-denominator-expansion}
\end{align}
Matching the exact target coefficient \cref{eq:onebody-target-u2v2} therefore forces
\begin{equation}
d_2=8,
\qquad\text{i.e.}\qquad
D_{\mathrm{target}}(u)=1-4u+8u^2+O(u^3).
\label{eq:onebody-denominator-target}
\end{equation}
This is the precise one-body datum that the \(\DtN\)-side closure must reproduce.

\subsection{Low-frequency \texorpdfstring{$\DtN$}{DtN} invariants and the carried freeze}
\label{sec:one-body-dtn}

The low-frequency monopole mouth operator is written as
\begin{equation}
Z_{00}(\omega;u)=Z_2(u)\,\omega^2+Z_4(u)\,\omega^4+O(\omega^6).
\label{eq:onebody-mouth-operator}
\end{equation}
From it we form the normalized invariant combinations
\begin{equation}
C_s(u)\equiv \frac{Z_4/Z_2^3}{(Z_4/Z_2^3)_0},
\qquad
G(u)\equiv \frac{Z_4/Z_2^2}{(Z_4/Z_2^2)_0}.
\label{eq:onebody-dtn-invariants}
\end{equation}
These are the two dimensionless one-body observables that survive the low-frequency
normalization and therefore provide the natural conservative building blocks for the
present closure.

On the cylinder / Neumann-bottom unit-test branch, the leading coefficients take the form
\begin{equation}
Z_2=-\frac{L}{c_s^2},
\qquad
Z_4=-\frac{L^3}{3c_s^4},
\label{eq:onebody-cylinder-z2z4}
\end{equation}
so the invariants reduce exactly to
\begin{equation}
C_s(u)=\frac{c_s^2(u)}{c_{s0}^2},
\qquad
G(u)=\frac{L(u)}{L_0}.
\label{eq:onebody-cylinder-invariants}
\end{equation}
Under the current conservative freeze, the Bernoulli continuation of the already-fixed
\(n=5\) barotropic sector gives
\begin{equation}
C_s(u)=1-4u,
\label{eq:onebody-Cs}
\end{equation}
while the carried \(\OnePN\) breathing slope implies
\begin{equation}
G(u)=1+g_1u+O(u^2),
\qquad
g_1=\frac{57}{64}.
\label{eq:onebody-G-slope}
\end{equation}
Only the linear coefficient \(g_1\) is needed for the \(\TwoPN\) one-body match in the
main text; higher-order series data are recorded in the appendices.

\subsection{Minimal denominator ansatz that preserves the frozen \texorpdfstring{$1$}{1}PN slot}
\label{sec:one-body-ansatz}

The carried \(\OnePN\) ledger fixes the coefficient of \(u\) in the effective velocity
denominator once and for all.  Therefore any admissible \(\TwoPN\) correction must begin
at order \(u^2\).  Within the minimal conservative \(\DtN\)-invariant packaging, the most
general correction built from the one-body invariants and truncated at quadratic order in
\(G-1\) is
\begin{equation}
D_{\eff}(u)=C_s(u)\left[1+\alpha\,(G-1)+\beta\,(G-1)^2\right].
\label{eq:onebody-general-ansatz}
\end{equation}
This is the smallest ansatz that can both preserve the already-fixed \(\OnePN\) slot and
supply a genuinely new \(\TwoPN\) correction.

\begin{claim}[Minimal one-body \texorpdfstring{$2$}{2}PN closure]
\label{clm:one-body-slot}
Within the conservative ansatz \cref{eq:onebody-general-ansatz}, preservation of the
frozen \(\OnePN\) coefficient forces
\begin{equation}
\alpha=0,
\label{eq:onebody-alpha-zero}
\end{equation}
and the exact isotropic Schwarzschild test-mass target through \(\TwoPN\) uniquely fixes
\begin{equation}
D_{\eff}(u)
=C_s(u)\left[1+\mu\,(G(u)-1)^2\right],
\qquad
\mu=\frac{8}{g_1^2}=\frac{32768}{3249},
\label{eq:onebody-final-closure}
\end{equation}
so that
\begin{equation}
D_{\eff}(u)=1-4u+8u^2+O(u^3).
\label{eq:onebody-final-target}
\end{equation}
\end{claim}

The derivation is immediate.  Insert
\begin{equation}
G(u)=1+g_1u+g_2u^2+O(u^3),
\qquad
C_s(u)=1-4u+O(u^2),
\label{eq:onebody-series-invariants}
\end{equation}
into \cref{eq:onebody-general-ansatz}.  One finds
\begin{equation}
D_{\eff}(u)=1+(-4+\alpha g_1)u+O(u^2).
\label{eq:onebody-linear-coefficient}
\end{equation}
Because the coefficient of \(u\) is already frozen by the carried \(\OnePN\) ledger and
because \(g_1=57/64\neq 0\), the linear correction must vanish, which proves
\cref{eq:onebody-alpha-zero}.  The first admissible new contribution is therefore
quadratic:
\begin{equation}
D_{\eff}(u)=C_s(u)\left[1+\mu\,(G(u)-1)^2\right].
\label{eq:onebody-quadratic-ansatz}
\end{equation}
Using \cref{eq:onebody-G-slope}, this expands to
\begin{equation}
D_{\eff}(u)=1-4u+\mu g_1^2u^2+O(u^3).
\label{eq:onebody-quadratic-series}
\end{equation}
Matching the exact target denominator \cref{eq:onebody-denominator-target} therefore
requires
\begin{equation}
\mu g_1^2=8,
\label{eq:onebody-mu-equation}
\end{equation}
which gives precisely \cref{eq:onebody-final-closure}.

Several features of this result are worth emphasizing.  First, the closure is minimal:
the linear geometry correction is forbidden by the frozen \(\OnePN\) slot, so the first
allowed new term is exactly quadratic in \(G-1\).  Second, the match is exact at the
level required here: once the current slope \(g_1=57/64\) is imported, the resulting
coefficient \(\mu=32768/3249\) reproduces the exact one-body denominator target through
\(\TwoPN\).  Third, the construction is conservative and invariant-based: nothing in the
argument depends on a coordinate-specialized comparable-mass fit or on a re-tuning of the
already-settled lower-order ledger.

\claimclass{\Controlled{} for the low-frequency \(\DtN\) packaging and for the import of
\(C_s(u)=1-4u\) and \(g_1=57/64\) from the carried Bernoulli / breathing ledger.
\Protocol{} for the minimal conservative ansatz
\(D_{\eff}=C_s[1+\alpha(G-1)+\beta(G-1)^2]\).
\FullMatch{} for the exact equality
\(D_{\eff}(u)=1-4u+8u^2+O(u^3)\) with the isotropic Schwarzschild test-mass target
through \(\TwoPN\).}

\subsection{Scope of the one-body result}
\label{sec:one-body-status}

This section closes the one-body missing \(\TwoPN\) response slot, but it is important to
state exactly what that means.

What has been established is the following.  Within the declared conservative
low-frequency \(\DtN\) packaging, the frozen \(\OnePN\) coefficient forces the geometry
correction to begin quadratically in \(G-1\), and the exact isotropic Schwarzschild
target then fixes that quadratic coefficient uniquely.  As a result, the one-body
velocity denominator is no longer an open slot in the \(\TwoPN\) assembly.

What has \emph{not} been claimed is equally important.  This section does not derive the
quadratic packaging \cref{eq:onebody-final-closure} from a fully solved moving-throat PDE
problem.  It also does not by itself determine the static \(U^3\) coefficient or any
comparable-mass \(\TwoPN\) cross term.  Those belong, respectively, to the self/static
sector of the next section and to the exact \(\ADM\) residual solve that follows later.
The appendix records the longer constructive throat-response story---including the
relation to the earlier raw resonance-proxy fit and the PDE-side \(\DtN\) scaffolding---
but none of that extra structure is needed for the proof line established here.

% ============================================================
\section{Minimal self/static \texorpdfstring{$2$}{2}PN sector}
\label{sec:self-static}

The previous section closed the one-body \(\TwoPN\) response slot by fixing the
missing denominator correction needed to reproduce the exact isotropic Schwarzschild
test-mass target.  The present section records the complementary piece of the proof:
the smallest self/static sector that can be carried into the full two-body assembly
without disturbing the frozen Newtonian+\(\OnePN\) ledger.

There are really two ingredients here.  The first is the \emph{self} sector: a minimal
Bernoulli exactification of the one-body worldline that preserves the already-fixed
\(\OnePN\) coefficients while generating the universal \(\TwoPN\) free and optical terms.
The second is the \emph{static} sector: quadratic local mass scaling, treated by the
same pair-counting logic used at lower order, whose two-body coefficient is then fixed
by the pure-static isotropic test-mass gate.  Once those two ingredients are frozen and
the denominator repair from \cref{sec:one-body} is inserted, the resulting
self/static block is ready to be fed into the exact comparable-mass residual solve of
later sections.

\subsection{Minimal Bernoulli exactification of the self sector}
\label{sec:self-static-self}

The carried \(n=5\) Bernoulli continuation fixes the sound-speed factor
\begin{equation}
C_s(u)=1-4u,
\qquad
u\equiv \frac{U}{c^2},
\label{eq:self-static-Cs}
\end{equation}
with \(U=GM/r\) in the strict one-body limit.  If one exactifies only the optical
denominator while keeping the already-frozen linear mass dressing, the minimal one-body
self worldline is
\begin{equation}
L_{\mathrm{self},\min}
=-m c^2(1-u)
\sqrt{1-\frac{v^2/c^2}{1-4u}}.
\label{eq:self-static-self-minimal}
\end{equation}
This is the smallest Bernoulli continuation that leaves the lower-order bookkeeping
intact: it preserves the frozen \(\OnePN\) self coefficient and introduces no new
static term by hand.

Expanding \cref{eq:self-static-self-minimal} through total post-Newtonian order
\(c^{-4}\) gives
\begin{align}
\frac{L_{\mathrm{self},\min}}{m}
={}&-c^2
+\frac12 v^2+U
+\frac{1}{c^2}\left(\frac18 v^4+\frac32 Uv^2\right)
\nonumber\\
&
+\frac{1}{c^4}\left(\frac1{16}v^6+\frac78 Uv^4+6U^2v^2\right)
+O(c^{-6}).
\label{eq:self-static-self-expansion}
\end{align}
Several features of \cref{eq:self-static-self-expansion} are load-bearing.
First, the carried \(\OnePN\) self coefficient remains exactly \(+\frac32\) on the
\(Uv^2\) term, as required by the frozen ledger.  Second, the universal free kinetic
coefficient is fixed to \(1/16\).  Third, the minimal exactification generates the
candidate one-body \(Uv^4\) coefficient \(7/8\).  Finally, before any further response
repair is inserted, the raw self sector places coefficient \(6\) on the \(U^2v^2\) term.

That last coefficient is precisely why the previous section was needed.  Comparing
\cref{eq:self-static-self-expansion} with the exact isotropic Schwarzschild target
\cref{eq:onebody-schwarzschild-2pn}, one sees that the raw minimal self sector already
gets the free \(v^6\) and mixed \(Uv^4\) pieces right, but it overshoots the required
\(U^2v^2\) coefficient by
\begin{equation}
6-2=4.
\label{eq:self-static-u2v2-residual}
\end{equation}
The role of \cref{sec:one-body} was exactly to repair this residual by replacing the raw
optical denominator \(1-4u\) with the \(\DtN\)-completed denominator
\(D_{\mathrm{eff}}(u)=1-4u+8u^2+O(u^3)\).

It is also worth recording what \emph{is not} done here.  One might try to Bernoulli-
exactify the mass prefactor as well, replacing \((1-u)\) by a fully Bernoulli-dressed
factor.  But that more aggressive step immediately spoils the frozen lower-order static
ledger: it generates a \(+\frac32 U^2/c^2\) static term at \(\OnePN\), whereas the carried
pair-counted static sector is already fixed to \(-\frac12 U^2/c^2\).  Thus only the
optical denominator is exactified in the self sector.  This is the sense in which the
present construction is \emph{minimal}.

\claimclass{\Controlled{} for the Bernoulli continuation
\cref{eq:self-static-self-minimal}, which preserves the frozen \(\OnePN\) self ledger
while generating the universal raw \(\TwoPN\) self coefficients
\(1/16\), \(7/8\), and \(6\).}

\subsection{Quadratic local mass scaling and the pure-static gate}
\label{sec:self-static-static}

The static sector is treated by the same local-potential logic used at lower order, but
extended to quadratic order in the local field.  For each body,
\begin{equation}
m_A^{\mathrm{eff}}=m_A\left(1-\frac{U_A}{c^2}+\lambda_\rho\frac{U_A^2}{c^4}\right),
\qquad
m_B^{\mathrm{eff}}=m_B\left(1-\frac{U_B}{c^2}+\lambda_\rho\frac{U_B^2}{c^4}\right),
\label{eq:self-static-meff}
\end{equation}
with the local Newtonian pair potentials
\begin{equation}
U_A=\frac{Gm_B}{r},
\qquad
U_B=\frac{Gm_A}{r}.
\label{eq:self-static-UAUB}
\end{equation}
The pair-counted static contribution is then
\begin{equation}
L_{\mathrm{static}}
=-\frac12\bigl(m_A^{\mathrm{eff}}\Phi_A+m_B^{\mathrm{eff}}\Phi_B\bigr),
\qquad
\Phi_A=-U_A,
\quad
\Phi_B=-U_B.
\label{eq:self-static-paircount}
\end{equation}
Expanding \cref{eq:self-static-paircount} gives, through \(\TwoPN\),
\begin{equation}
L_{\mathrm{static}}
=\frac{Gm_A m_B}{r}
-\frac{G^2m_A m_B(m_A+m_B)}{2c^2r^2}
+\frac{\lambda_\rho G^3m_A m_B(m_A^2+m_B^2)}{2c^4r^3}
+O(c^{-6}).
\label{eq:self-static-static-expansion}
\end{equation}
Thus the two-body static \(\TwoPN\) coefficient is simply \(\lambda_\rho/2\).

The one-body meaning of this coefficient is immediate.  In the strict test-mass limit,
\cref{eq:self-static-static-expansion} contributes
\begin{equation}
\frac{L_{\mathrm{static},1\text{-}\mathrm{body}}}{m}
=U-\frac{1}{2c^2}U^2+\frac{\lambda_\rho}{2c^4}U^3+O(c^{-6}).
\label{eq:self-static-onebody-static}
\end{equation}
The exact isotropic Schwarzschild target of
\cref{eq:onebody-schwarzschild-2pn} requires the \(U^3\) coefficient to be \(+\frac14\).
Therefore the pure-static test-mass gate fixes
\begin{equation}
\lambda_\rho=\frac12.
\label{eq:self-static-lambda-rho}
\end{equation}
This is the unique value compatible with the already-frozen \(\OnePN\) static term and
with the exact isotropic \(\TwoPN\) static coefficient.

\begin{claim}[Minimal static \texorpdfstring{$2$}{2}PN closure]
\label{clm:self-static-lambda-rho}
Within quadratic local mass scaling and the carried pair-counting rule,
matching the pure-static one-body target through \(\TwoPN\) uniquely sets
\begin{equation}
\lambda_\rho=\frac12.
\label{eq:self-static-lambda-rho-claim}
\end{equation}
Equivalently, the two-body static \(\TwoPN\) block is fixed to
\begin{equation}
L_{2,\mathrm{static}}
=\frac{G^3m_A m_B(m_A^2+m_B^2)}{4r^3},
\label{eq:self-static-L2static}
\end{equation}
where, as throughout, \(L_2\) denotes the coefficient of \(c^{-4}\) in the expansion
\(L_{\rm cons}=L_0+c^{-2}L_1+c^{-4}L_2+O(c^{-6})\).
\end{claim}

At the level needed for the main proof, \cref{eq:self-static-lambda-rho} is the only
new static datum that must be carried forward.  The same quadratic local mass scaling
also induces a higher-body cubic-static hierarchy.  Those terms are important for the
broader constructive program and are recorded in the appendices, but they are not needed
for the conservative two-body \(\TwoPN\) theorem of the present paper.

\claimclass{\Controlled{} for quadratic local mass scaling and pair counting.
\FullMatch{} for the pure-static isotropic one-body gate
\(\lambda_\rho=\frac12\).}

\subsection{The frozen self/static seed used in the two-body assembly}
\label{sec:self-static-seed}

It is useful to separate the raw self/static seed from its final one-body repair.
Combining \cref{eq:self-static-self-expansion} with
\cref{eq:self-static-onebody-static} gives the strict one-body self/static candidate
\begin{align}
\frac{L^{\rm raw}_{\mathrm{self}+\mathrm{static},1\text{-}\mathrm{body}}}{m}
={}&-c^2+\frac12 v^2+U
+\frac{1}{c^2}\left(\frac18 v^4+\frac32 Uv^2-\frac12 U^2\right)
\nonumber\\
&
+\frac{1}{c^4}\left(\frac1{16}v^6+\frac78 Uv^4+6U^2v^2+\frac{\lambda_\rho}{2}U^3\right)
+O(c^{-6}).
\label{eq:self-static-raw-onebody}
\end{align}
After imposing \(\lambda_\rho=\frac12\), the pure-static target is matched exactly, and
one is left with the sole residual \(4U^2v^2\) already isolated in
\cref{eq:self-static-u2v2-residual}.  That residual is what the denominator closure of
\cref{sec:one-body} removes.

Accordingly, once the one-body \(\DtN\) repair is inserted, the one-body self/static
seed becomes
\begin{align}
\frac{L_{\mathrm{self}+\mathrm{static},1\text{-}\mathrm{body}}}{m}
={}&-c^2+\frac12 v^2+U
+\frac{1}{c^2}\left(\frac18 v^4+\frac32 Uv^2-\frac12 U^2\right)
\nonumber\\
&
+\frac{1}{c^4}\left(\frac1{16}v^6+\frac78 Uv^4+2U^2v^2+\frac14 U^3\right)
+O(c^{-6}),
\label{eq:self-static-final-onebody}
\end{align}
which is now identical to the exact target
\cref{eq:onebody-schwarzschild-2pn}.

The corresponding generic-frame two-body self/static block is then obtained by the usual
symmetric promotion of the one-body monomials:
\begin{equation}
L_{2,\mathrm{self}+\mathrm{static}}=
\frac{m_A v_A^6+m_B v_B^6}{16}
+\frac{7Gm_A m_B}{8r}(v_A^4+v_B^4)
+\frac{2G^2m_A m_B}{r^2}(m_B v_A^2+m_A v_B^2)
+\frac{G^3m_A m_B(m_A^2+m_B^2)}{4r^3}.
\label{eq:self-static-final-two-body}
\end{equation}
This is the entire self/static input used later in the comparable-mass \(\ADM\) lift.
No quartic cross-velocity monomial appears in \cref{eq:self-static-final-two-body};
there is no term proportional to \(v_{AB}(v_A^2+v_B^2)\), \(v_A^2v_B^2\),
\(v_{AB}^2\), or \(v_{An}v_{Bn}\)-type structures.  Likewise, at order \(G^2/r^2\)
only the direct one-body-descended terms
\((m_B v_A^2+m_A v_B^2)\) are present.  The necessity of the missing comparable-mass
cross structures is therefore not hidden inside the self/static sector; it will emerge
cleanly as a residual in the next-stage \(\ADM\) comparison.

\begin{claim}[Frozen self/static \texorpdfstring{$2$}{2}PN seed]
\label{clm:self-static-seed}
Within the closure hierarchy of the present paper, the full self/static contribution to
\(L_2\) is exactly \cref{eq:self-static-final-two-body}.  In the strict one-body limit
it reduces to \cref{eq:self-static-final-onebody}, which matches the exact isotropic
Schwarzschild test-mass expansion through \(\TwoPN\).
\end{claim}

\claimclass{\Controlled{} for the minimal Bernoulli self exactification and quadratic
local mass scaling.
\FullMatch{} for the statement that the repaired one-body self/static seed matches the
exact isotropic Schwarzschild target through \(\TwoPN\).}

\subsection{Scope and regime of use}
\label{sec:self-static-scope}

Two cautions should be kept visible.

First, \cref{eq:self-static-final-two-body} is not yet the full conservative two-body
\(\TwoPN\) answer.  It is only the self/static seed.  Once inserted into the exact
Legendre-transform comparison, it leaves a nonzero comparable-mass residual.  The whole
point of the next section is to solve that residual in a symmetry-respecting invariant
basis and thereby determine the genuinely new \(\TwoPN\) cross block.

Second, the present sector is derived under the same controlled small-body logic as the
rest of the effective particle reduction.  In particular, universal \(\TwoPN\) terms
must dominate over explicit finite-size corrections.  The appendix derives the sharp
gate
\begin{equation}
\frac{F_{\rm finite\text{-}size}^{(2)}}{F_{\rm universal}^{(2)}}
\sim \left(\frac{a c^2}{GM}\right)^2,
\label{eq:self-static-finite-size-gate}
\end{equation}
so the relevant regime is not merely \(a\ll r\), but the stronger universal-\(\TwoPN\)
condition \(a\ll GM/c^2\).  This caveat does not alter the algebraic two-body proof
once the regime is declared, but it is part of the honest scope statement of the paper.

% ============================================================
\section{Comparable-mass \texorpdfstring{$2$}{2}PN \texorpdfstring{$\ADM$}{ADM} lift}
\label{sec:adm-lift}

The previous two sections fixed everything that can be read off from the one-body
sector and from the minimal self/static continuation of the carried toy-model
ledger.  But that still does not prove the full conservative two-body
\(\TwoPN\) result.  Once the repaired one-body denominator and the static gate
\(\lambda_\rho=1/2\) have been imposed, there can still remain genuinely
comparable-mass terms that vanish in the strict test-mass limit and therefore cannot
be seen by one-body matching alone.

This section isolates exactly those missing terms and solves for them by matching to
the standard generic-frame \(\ADM\) Hamiltonian through \(\TwoPN\).  The key point
is that, once the Newtonian+\(\OnePN\) ledger is frozen, the passage from the
Lagrangian to the Hamiltonian becomes algebraically simple enough that every new
\(\TwoPN\) cross monomial enters linearly.  The comparable-mass lift is therefore
not an open-ended fit: it is an exact residual solve in a compact invariant basis.

\subsection{Perturbative Legendre transform through \texorpdfstring{$2$}{2}PN}
\label{sec:adm-legendre}

Let
\begin{equation}
L=L_0+\epsilon L_1+\epsilon^2 L_2,
\qquad
\epsilon\equiv c^{-2},
\label{eq:adm-epsilon-expansion}
\end{equation}
with \(L_0\) quadratic in the velocities,
\begin{equation}
L_0=\frac12 m_A v_A^2+\frac12 m_B v_B^2+\frac{Gm_A m_B}{r}.
\label{eq:adm-L0-repeat}
\end{equation}
Introduce the six-vector notation
\(v=(\vct_A,\vct_B)\),
\(p=(\mathbf p_A,\mathbf p_B)\),
and the diagonal mass matrix
\begin{equation}
M\equiv \mathrm{diag}(m_A I_3,\,m_B I_3).
\label{eq:adm-mass-matrix}
\end{equation}
Then
\begin{equation}
p=\frac{\partial L}{\partial v}
=Mv+\epsilon\frac{\partial L_1}{\partial v}
+\epsilon^2\frac{\partial L_2}{\partial v}.
\label{eq:adm-momentum-expansion}
\end{equation}
Writing the perturbative inverse as
\begin{equation}
v=v_0+\epsilon v_1+\epsilon^2 v_2+O(\epsilon^3),
\qquad
v_0=M^{-1}p=
\left(\frac{\mathbf p_A}{m_A},\frac{\mathbf p_B}{m_B}\right),
\label{eq:adm-v-expansion}
\end{equation}
and defining
\begin{equation}
A_0\equiv
\left.\frac{\partial L_1}{\partial v}\right|_{v=v_0},
\label{eq:adm-A0}
\end{equation}
the Legendre transform through \(\TwoPN\) collapses to the standard compact form
\begin{equation}
H_1=-L_1(v_0),
\qquad
H_2=-L_2(v_0)+\frac12 A_0^T M^{-1}A_0.
\label{eq:adm-H1H2-formula}
\end{equation}
The derivation is routine and is deferred to
Appendix~\ref{app:adm-details}; the only point needed in the main text is the
consequence.  Once the full carried \(\OnePN\) block is frozen, the quadratic
correction \(\frac12 A_0^T M^{-1}A_0\) is already fixed.  Therefore any
\emph{new} \(\TwoPN\) Lagrangian block contributes to the Hamiltonian only through
\begin{equation}
\Delta H_{2,\mathrm{new}}=-L_{2,\mathrm{new}}(v_0).
\label{eq:adm-new-block-linear}
\end{equation}
This is the algebraic reason the comparable-mass lift is tractable: after the lower-order
ledger is frozen, the remaining \(\TwoPN\) matching problem is linear in the unknown
cross coefficients.

\claimclass{\Exact{} for the perturbative Legendre-transform identity
\cref{eq:adm-H1H2-formula}.}

\subsection{Residual structure with no new cross sector}
\label{sec:adm-residual-structure}

Now insert the full frozen Newtonian+\(\OnePN\) ledger together with the repaired
self/static \(\TwoPN\) seed from \cref{eq:self-static-final-two-body}.  By
construction, the \(\OnePN\) piece still Legendre-transforms exactly to the standard
\(\EIH/\ADM\) target at order \(c^{-2}\).  The only possible mismatch therefore sits
at order \(c^{-4}\).

If one includes \emph{no} new comparable-mass \(\TwoPN\) cross sector beyond the
self/static seed, the Legendre-transformed Hamiltonian does \emph{not} yet equal the
standard generic-frame \(\ADM\) target.  But the mismatch has an important and very
clean form: every missing term is a genuinely comparable-mass structure.  In
particular, the residual contains only three classes of monomials:
\begin{enumerate}
\item quartic cross-velocity terms proportional to \(G/r\),
\item quadratic-velocity comparable-mass terms proportional to \(G^2/r^2\),
\item one static cross mass polynomial proportional to \(G^3 m_A^2 m_B^2/r^3\).
\end{enumerate}
No new free \(v^6\) term appears.  No new one-body \(Uv^4\), \(U^2v^2\), or \(U^3\)
term appears.  Those sectors were already fixed by the minimal self/static seed and by
the one-body denominator closure of \cref{sec:one-body}.  What remains is therefore
purely a comparable-mass lift.

This observation is more than a classification.  It means that the residual vanishes in
the strict test-mass limit with the heavy source held at rest.  So the yet-missing
\(\TwoPN\) content is not hiding in the one-body closure; it is exactly the additional
two-body cross data needed to lift the repaired one-body answer to the full comparable-
mass problem.

\subsection{Compact invariant basis and unique solve}
\label{sec:adm-unique-solve}

To solve the residual, we use the smallest symmetry-respecting generic-frame basis built
from the carried invariants
\(v_A^2\), \(v_B^2\), \(v_{AB}\), \(v_{An}\), and \(v_{Bn}\).  For the
\(G/r\) quartic sector we take
\begin{align}
Q_1&=v_{AB}(v_A^2+v_B^2),
&
Q_2&=v_{An}v_{Bn}(v_A^2+v_B^2),
&
Q_3&=v_A^2v_B^2,
\nonumber\\
Q_4&=v_{AB}^2,
&
Q_5&=v_A^2v_{Bn}^2+v_B^2v_{An}^2,
&
Q_6&=v_{AB}v_{An}v_{Bn},
&
Q_7&=v_{An}^2v_{Bn}^2.
\label{eq:adm-Q-basis}
\end{align}
For the \(G^2/r^2\) quadratic-velocity sector we take
\begin{align}
T_1&=m_B v_A^2+m_A v_B^2,
&
T_2&=m_A v_A^2+m_B v_B^2,
&
T_3&=(m_A+m_B)v_{AB},
\nonumber\\
T_4&=(m_A+m_B)v_{An}v_{Bn},
&
T_5&=m_B v_{An}^2+m_A v_{Bn}^2,
&
T_6&=m_A v_{An}^2+m_B v_{Bn}^2.
\label{eq:adm-T-basis}
\end{align}
The static comparable-mass slot is spanned by the single monomial
\begin{equation}
S_1\equiv \frac{G^3 m_A^2 m_B^2}{r^3}.
\label{eq:adm-S-basis}
\end{equation}
Accordingly, the most general added comparable-mass block in this compact basis is
\begin{equation}
L_{2,\mathrm{cross}}^{\mathrm{ans}}
=\frac{Gm_A m_B}{r}\sum_{i=1}^{7} q_i Q_i
+\frac{G^2m_A m_B}{r^2}\sum_{i=1}^{6} t_i T_i
+s_1 S_1.
\label{eq:adm-cross-ansatz}
\end{equation}
Because of \cref{eq:adm-new-block-linear}, matching the residual Hamiltonian to the
standard generic-frame \(\ADM\) target is now a linear algebra problem for the unknown
coefficients \(q_i\), \(t_i\), and \(s_1\).

\begin{claim}[Unique comparable-mass residual solve]
\label{clm:adm-solve}
In the compact invariant basis
\cref{eq:adm-Q-basis,eq:adm-T-basis,eq:adm-S-basis}, the residual solve is unique and
fixes the added comparable-mass \(\TwoPN\) coefficients to
\begin{align}
q_1&=-\frac74,
& q_2&=-\frac14,
& q_3&=\frac{11}{8},
& q_4&=\frac14,
\nonumber\\
q_5&=-\frac58,
& q_6&=\frac32,
& q_7&=\frac38,
\label{eq:adm-q-solution}
\\[0.5ex]
t_1&=0,
& t_2&=\frac{11}{8},
& t_3&=-\frac{15}{4},
& t_4&=0,
& t_5&=0,
& t_6&=\frac{15}{8},
\label{eq:adm-t-solution}
\\[0.5ex]
s_1&=\frac54.
\label{eq:adm-s-solution}
\end{align}
\end{claim}

Three of the quadratic basis coefficients therefore vanish automatically,
\begin{equation}
t_1=t_4=t_5=0.
\label{eq:adm-vanishing-t}
\end{equation}
This is useful structurally as well as algebraically.  It shows that the missing
comparable-mass information does \emph{not} require every symmetry-allowed monomial.
The exact \(\ADM\) lift lands in a still smaller subspace than the original basis makes
available.

\claimclass{\Exact{} for the residual solve once the frozen lower-order ledger and the
declared self/static \(\TwoPN\) seed are inserted.}

\subsection{Added comparable-mass cross block}
\label{sec:adm-added-cross}

Substituting the solved coefficients into \cref{eq:adm-cross-ansatz} gives the required
added comparable-mass \(\TwoPN\) block,
\begin{align}
L_{2,\mathrm{cross}}=
&\frac{Gm_A m_B}{r}
\Big[
-\frac74\,v_{AB}(v_A^2+v_B^2)
-\frac14\,v_{An}v_{Bn}(v_A^2+v_B^2)
+\frac{11}{8}\,v_A^2v_B^2
+\frac14\,v_{AB}^2 \notag\\
&\hspace{7.25em}
-\frac58\,(v_A^2v_{Bn}^2+v_B^2v_{An}^2)
+\frac32\,v_{AB}v_{An}v_{Bn}
+\frac38\,v_{An}^2v_{Bn}^2
\Big] \notag\\
&+\frac{G^2m_A m_B}{r^2}
\Big[
\frac{11}{8}(m_A v_A^2+m_B v_B^2)
-\frac{15}{4}(m_A+m_B)v_{AB}
+\frac{15}{8}(m_A v_{An}^2+m_B v_{Bn}^2)
\Big]
\notag\\
&+\frac{5G^3m_A^2m_B^2}{4r^3}.
\label{eq:adm-added-cross-block}
\end{align}
As throughout, the overall post-Newtonian factor \(c^{-4}\) is understood.  Equation
\cref{eq:adm-added-cross-block} is one of the central equations of the paper: it is the
unique extra comparable-mass content that must be added to the repaired self/static seed
in order to lift the one-body-complete \(\TwoPN\) answer to the full generic-frame
conservative two-body target.

Several features of \cref{eq:adm-added-cross-block} should be noted immediately.
First, its entire quartic \(G/r\) sector is genuinely cross-coupled; no term reduces to
another copy of a one-body monomial.  Second, in the \(G^2/r^2\) sector, only the
\(T_2\), \(T_3\), and \(T_6\) structures survive the exact solve.  Third, the static
cross term is positive and precisely large enough to upgrade the self/static mass
polynomial from
\((m_A^2+m_B^2)\) to
\((m_A^2+5m_A m_B+m_B^2)\)
once the full \(\TwoPN\) block is assembled in the next section.

\subsection{Sanity checks to report in the main text}
\label{sec:adm-sanity}

Three quick checks should be stated explicitly before moving on to the final assembly.

First, the frozen \(\OnePN\) block still transforms exactly to the standard generic-frame
\(\ADM\) Hamiltonian at order \(c^{-2}\).  The present solve does not repair a hidden
lower-order error; it addresses only a genuine \(\TwoPN\) comparable-mass residual.

Second, the full added cross block \cref{eq:adm-added-cross-block} vanishes in the
strict test-mass limit with the heavy source held at rest.  The quartic \(G/r\) terms
all require cross kinematics and therefore disappear immediately.  The
\(G^2/r^2\) and static pieces are suppressed by at least one additional power of the
light mass, so after dividing by the test mass they vanish as well.  This is exactly as
it should be: the one-body closure had already fixed the entire strict test-mass sector.

Third, the static mass polynomial comes out correctly only after the comparable-mass lift
is included.  The repaired self/static seed contributed
\begin{equation}
\frac{G^3m_A m_B}{4r^3}(m_A^2+m_B^2),
\label{eq:adm-static-seed-masspoly}
\end{equation}
whereas the added cross block contributes
\begin{equation}
\frac{5G^3m_A^2m_B^2}{4r^3}
=\frac{G^3m_A m_B}{4r^3}(5m_A m_B).
\label{eq:adm-static-cross-masspoly}
\end{equation}
Together these give the exact \(\ADM\)-compatible polynomial
\begin{equation}
\frac{G^3m_A m_B}{4r^3}(m_A^2+5m_A m_B+m_B^2).
\label{eq:adm-static-full-masspoly}
\end{equation}
That upgrade is a useful one-line check that the comparable-mass residual solve is not
merely qualitative; it lands on the exact mass structure required by the standard
Hamiltonian target.

The next section simply assembles
\cref{eq:self-static-final-two-body,eq:adm-added-cross-block} into the full two-body
\(\TwoPN\) Lagrangian and records the exact equality to the standard generic-frame
\(\ADM\) result.

% ============================================================
\section{Full two-body \texorpdfstring{$2$}{2}PN Lagrangian and exact \texorpdfstring{$\ADM$}{ADM} match}
\label{sec:full-assembly}

The previous sections isolated the proof into logically distinct parts.  The carry-forward
Newtonian+\(\OnePN\) ledger was frozen.  The one-body missing \(\TwoPN\) response slot was
closed by the denominator repair of \cref{sec:one-body}.  The minimal self/static
\(\TwoPN\) seed was then fixed in \cref{sec:self-static}.  Finally,
\cref{sec:adm-lift} showed that the remaining comparable-mass discrepancy is a purely
algebraic residual and solved it uniquely in a compact invariant basis.

This section simply assembles those ingredients and states the final theorem in the same
style as the \(\OnePN\) EIH-match section.  The result is a complete conservative two-body
\(\TwoPN\) Lagrangian within the stated closure hierarchy, together with exact equality of
its Legendre transform to the standard generic-frame \(\ADM\) Hamiltonian target.

\subsection{Derived two-body Lagrangian through \texorpdfstring{$2$}{2}PN}
\label{sec:full-assembly-lagrangian}

By definition,
\begin{equation}
L_{\rm cons}^{\rm derived}
=
L_0+\frac{1}{c^2}L_1+\frac{1}{c^4}L_{2,\mathrm{full}}+O(c^{-6}),
\label{eq:full-assembly-Lcons}
\end{equation}
with the frozen lower-order blocks
\begin{equation}
L_0=
\frac12 m_A v_A^2
+\frac12 m_B v_B^2
+\frac{Gm_A m_B}{r},
\label{eq:full-assembly-L0}
\end{equation}
\begin{align}
L_1={}&
\frac18\bigl(m_A v_A^4+m_B v_B^4\bigr)
+\frac{Gm_A m_B}{r}
\left[
\frac32\bigl(v_A^2+v_B^2\bigr)
-\frac72\,v_{AB}
-\frac12\,v_{An}v_{Bn}
\right]
\nonumber\\
&
-\frac{G^2m_A m_B(m_A+m_B)}{2r^2},
\label{eq:full-assembly-L1}
\end{align}
and with the full \(\TwoPN\) block obtained by adding the frozen self/static seed
\cref{eq:self-static-final-two-body} to the uniquely solved comparable-mass cross block
\cref{eq:adm-added-cross-block}:
\begin{align}
L_{2,\mathrm{full}}=
&\frac{m_A v_A^6+m_B v_B^6}{16}
+\frac{Gm_A m_B}{r}
\Big[
\frac78\bigl(v_A^4+v_B^4\bigr)
-\frac74\,v_{AB}(v_A^2+v_B^2)
-\frac14\,v_{An}v_{Bn}(v_A^2+v_B^2)
\nonumber\\
&\hspace{7.25em}
+\frac{11}{8}\,v_A^2v_B^2
+\frac14\,v_{AB}^2
-\frac58\,(v_A^2v_{Bn}^2+v_B^2v_{An}^2)
+\frac32\,v_{AB}v_{An}v_{Bn}
+\frac38\,v_{An}^2v_{Bn}^2
\Big]
\nonumber\\
&+\frac{G^2m_A m_B}{r^2}
\Big[
\Bigl(2m_B+\frac{11}{8}m_A\Bigr)v_A^2
+\Bigl(2m_A+\frac{11}{8}m_B\Bigr)v_B^2
-\frac{15}{4}(m_A+m_B)v_{AB}
\nonumber\\
&\hspace{10.5em}
+\frac{15}{8}(m_A v_{An}^2+m_B v_{Bn}^2)
\Big]
+\frac{G^3m_A m_B}{4r^3}(m_A^2+5m_A m_B+m_B^2).
\label{eq:full-assembly-L2full}
\end{align}
Equation~\eqref{eq:full-assembly-L2full} is the final conservative two-body
\(\TwoPN\) Lagrangian block derived in this paper.

Written out in one line, the full conservative two-body Lagrangian through \(\TwoPN\) is
therefore
\begin{align}
L_{\rm cons}^{\rm derived}
={}&
\frac12 m_A v_A^2+\frac12 m_B v_B^2
+\frac{Gm_A m_B}{r}
\nonumber\\
&+\frac{1}{c^2}
\Bigg\{
\frac18\bigl(m_A v_A^4+m_B v_B^4\bigr)
+\frac{Gm_A m_B}{r}
\left[
\frac32\bigl(v_A^2+v_B^2\bigr)
-\frac72\,v_{AB}
-\frac12\,v_{An}v_{Bn}
\right]
\nonumber\\
&\hspace{5.2em}
-\frac{G^2m_A m_B(m_A+m_B)}{2r^2}
\Bigg\}
+\frac{1}{c^4}L_{2,\mathrm{full}}
+O(c^{-6}).
\label{eq:full-assembly-Lcons-expanded}
\end{align}

Several structural features are worth recording explicitly.
First, the free kinetic hierarchy now reads
\(\frac18 m v^4\) at \(\OnePN\) and
\(\frac1{16}m v^6\) at \(\TwoPN\), exactly as required by the carried Bernoulli/self
packaging.  Second, the repaired one-body/self/static data are visible in the
\(\frac78\), \(2\), and \(\frac14\) coefficients that survive in the strict test-mass
reduction of \cref{eq:full-assembly-L2full}.  Third, all genuinely comparable-mass new
content sits in the cross structures that were absent from the self/static seed and were
solved uniquely by the \(\ADM\) residual lift.

A useful bookkeeping remark is that every coefficient in
\cref{eq:full-assembly-L2full} now has a named provenance.  The coefficients
\(1/16\) and \(7/8\) descend from the minimal Bernoulli exactification of the self
sector.  The coefficients \(2\) and \(1/4\) are the repaired one-body
\(U^2v^2\) and pure-static slots after the denominator closure and static gate are
imposed.  The remaining quartic cross-velocity coefficients, the nontrivial
\(G^2/r^2\) comparable-mass coefficients, and the \(+\frac54 m_A^2m_B^2\) static cross
upgrade come from the exact residual solve of \cref{sec:adm-lift}.  So the final answer
is not a single omnibus ansatz; it is an assembly of independently audited sectors.

\subsection{Standard generic-frame \texorpdfstring{$\ADM$}{ADM} target and exact equality}
\label{sec:full-assembly-adm}

Let
\begin{equation}
H_{\rm cons}^{\rm derived}
=
H_0+\frac{1}{c^2}H_1+\frac{1}{c^4}H_2+O(c^{-6})
\label{eq:full-assembly-Hderived}
\end{equation}
be the perturbative Legendre transform of
\cref{eq:full-assembly-Lcons}.  The standard conservative generic-frame two-body
\(\ADM\) Hamiltonian through \(\TwoPN\) is written in the corresponding expansion
\begin{equation}
H_{\rm cons}^{\ADM}
=
H_0^{\ADM}+\frac{1}{c^2}H_{1\mathrm{PN}}^{\ADM}
+\frac{1}{c^4}H_{2\mathrm{PN}}^{\ADM}+O(c^{-6}).
\label{eq:full-assembly-Hadm}
\end{equation}
The Newtonian and \(\OnePN\) equalities were already frozen earlier:
\begin{equation}
H_0=H_0^{\ADM},
\qquad
H_1=H_{1\mathrm{PN}}^{\ADM}.
\label{eq:full-assembly-H01-match}
\end{equation}
Thus the only nontrivial content left is the order-\(c^{-4}\) block.

\begin{claim}[Full conservative two-body \texorpdfstring{$2$}{2}PN assembly]
\label{clm:full-2pn}
After restoring the overall factor \(c^{-4}\), the full assembled two-body
\(\TwoPN\) Lagrangian
\cref{eq:full-assembly-Lcons-expanded} Legendre-transforms exactly to the standard
conservative generic-frame \(\ADM\) Hamiltonian target through \(\TwoPN\).  Equivalently,
\begin{equation}
H_2=H_{2\mathrm{PN}}^{\ADM},
\qquad\text{so that}\qquad
\Delta H_2\equiv H_2-H_{2\mathrm{PN}}^{\ADM}=0.
\label{eq:full-assembly-exact-match}
\end{equation}
\end{claim}

The proof is now short because all of the hard work has already been separated out.
By \cref{eq:adm-H1H2-formula}, once the frozen \(\OnePN\) ledger is fixed, every new
\(\TwoPN\) Lagrangian term enters the Hamiltonian only through its direct value at
\(v_0=p/m\).  The self/static block
\cref{eq:self-static-final-two-body} supplies the repaired one-body content but leaves a
nonzero comparable-mass residual.  In \cref{sec:adm-lift}, that residual was solved
exactly in the compact invariant basis
\cref{eq:adm-Q-basis,eq:adm-T-basis,eq:adm-S-basis}, yielding the unique added cross
block \cref{eq:adm-added-cross-block}.  Substituting that block into the Legendre
transform removes the residual identically, which is precisely the statement of
\cref{eq:full-assembly-exact-match}.

It is worth stating explicitly what kind of equality this is.  The present paper does
not merely show that the assembled two-body Lagrangian has the right test-mass limit or
that its coefficients can be gauge-shifted into a familiar shape.  It shows something
stronger: within the declared closure hierarchy, the concrete assembled Lagrangian
\cref{eq:full-assembly-Lcons-expanded} has a perturbative Legendre transform whose
\(\TwoPN\) block is exactly the standard generic-frame \(\ADM\) target.  So the entire
conservative two-body \(\TwoPN\) coefficient set is closed.

A few short sanity checks should be kept visible here.  First, the frozen
\(\OnePN\) Lagrangian still transforms to the standard \(\ADM\) \(H_{1\mathrm{PN}}\) target
exactly, so the new \(\TwoPN\) completion does not spoil any lower-order result.
Second, the added comparable-mass cross block vanishes in the strict test-mass limit with
body \(B\) held at rest, confirming that no one-body coefficient has been silently
retuned.  Third, the static mass polynomial in the assembled \(\TwoPN\) block is upgraded
from
\((m_A^2+m_B^2)\) in the raw self/static seed to
\((m_A^2+5m_A m_B+m_B^2)\), exactly the mass structure required by the Hamiltonian-side
match.

\claimclass{\FullMatch{} for the exact equality
\cref{eq:full-assembly-exact-match}, with \Exact{} algebra and the previously declared
\Controlled{} / \Protocol{} inputs inherited from the one-body closure and self/static
seed.}

\subsection{Compact statement of the proof chain}
\label{sec:full-assembly-proof-chain}

The logical chain can now be summarized in one paragraph.
The exact projection / Poisson-hook / worldtube backbone supplies the reduced particle
variables and the already-frozen Newtonian+\(\OnePN\) ledger.  The one-body
\(\DtN\)-invariant denominator closure fixes the missing test-mass \(\TwoPN\)
\(U^2v^2\) slot without altering the carried \(\OnePN\) coefficient.  Minimal Bernoulli
exactification plus quadratic local mass scaling then fix the self/static seed, with the
pure-static isotropic gate selecting \(\lambda_\rho=1/2\).  Once those inputs are frozen,
the remaining mismatch with the standard generic-frame \(\ADM\) target is a purely
comparable-mass residual.  That residual is solved uniquely in a compact invariant basis,
and adding the resulting cross block yields the full assembled Lagrangian
\cref{eq:full-assembly-Lcons-expanded}.  The final equality
\cref{eq:full-assembly-exact-match} is therefore not a global fit; it is the end-point of
an auditable sector-by-sector proof chain.

One immediate consequence is worth stating outright.  Once the closure hierarchy of the
present paper is accepted, there are no remaining free dimensionless coefficients in the
full conservative two-body \(\TwoPN\) sector.  What remains open is not the algebraic
assembly, but the deeper constructive explanation of some of its new comparable-mass
coefficients from the inner-throat response side.

\subsection{Short interpretive pointer to the appendices}
\label{sec:full-assembly-interpretive-pointer}

The main proof stops at the exact \(\ADM\) match, because that is the smallest readable
route to the conservative \(\TwoPN\) theorem.  But the assembled cross block is not an
opaque remainder term.  The appendix program shows that it already admits a compact
constructive packaging: the frozen odd dipole wake from \(\OnePN\), a genuinely new even
\(P_0\oplus P_2\) support layer, and a pure-potential geometry-closure term.  Those
appendix developments are important for understanding \emph{why} the solved coefficients
organize so economically, but they are not needed for the main proof that the assembled
Lagrangian matches the standard conservative generic-frame \(\ADM\) target exactly.

% ============================================================
\section{Fixed versus symbolic quantities}
\label{sec:status-ledger}

One of the main purposes of the present paper is to make the coefficient status of the
\(
\TwoPN
\) toy-model program transparent.
Earlier stages of the project necessarily mixed together several logically different
objects:
\begin{enumerate}[leftmargin=1.5em]
  \item quantities already fixed by exact identities, controlled reductions, or exact
  algebraic matching,
  \item quantities fixed only after a specific conservative response hierarchy is
  declared,
  \item appendix-side throat data that are now fixed at zero frequency by the solved
  conservative \(\TwoPN\) cross target, and
  \item quantities that should still remain symbolic because they belong to dynamic
  \(\DtN\)/PDE observables, finite-size structure, or sectors beyond the conservative
  near-zone \(\TwoPN\) problem.
\end{enumerate}

Now that the full conservative two-body \(\TwoPN\) assembly has been completed, it is
useful to record that distinction explicitly.
The point of the present section is not to soften the main theorem, but to sharpen it.
Within the closure hierarchy adopted in this paper, the conservative two-body
\(\TwoPN\) sector contains no remaining free \emph{dimensionless} coefficients.
What remains symbolic belongs either to the declared response packaging itself, to
appendix-side dynamic throat observables, or to physics beyond the scope of the present
conservative derivation.

\subsection{Quantities fixed by the present derivation chain}
\label{sec:status-fixed}

The following quantities should now be regarded as fixed carry-forward values or fixed
outputs of the present \(\TwoPN\) derivation.

\paragraph{Carried lower-order backbone.}
The lower-order ledger remains frozen throughout the paper.
In particular,
\begin{equation}
\krho=1,
\qquad
n=5,
\qquad
\kadd=\frac12,
\label{eq:status-carried-lower-order}
\end{equation}
while the conservative parity-even wake coefficients remain
\begin{equation}
C_\parallel=-\frac72,
\qquad
C_L=-\frac12,
\label{eq:status-carried-wake}
\end{equation}
and the full Newtonian+\(\OnePN\) two-body Lagrangian is already fixed to equal the
standard conservative \(\EIH\) target exactly.
These are no longer tunable inputs.
They are the frozen backbone on which the \(\TwoPN\) calculation is built.

\paragraph{One-body \texorpdfstring{$2$}{2}PN denominator closure.}
The missing test-mass response slot is fixed by the minimal invariant denominator
closure of \cref{sec:one-body}:
\begin{equation}
\alpha=0,
\qquad
\mu=\frac{32768}{3249},
\qquad
D_{\eff}(u)=C_s(u)\left[1+\frac{32768}{3249}(G(u)-1)^2\right].
\label{eq:status-fixed-denominator}
\end{equation}
Equivalently,
\begin{equation}
D_{\eff}(u)=1-4u+8u^2+O(u^3),
\label{eq:status-fixed-denominator-series}
\end{equation}
so the exact isotropic Schwarzschild one-body target is reproduced through
\(\TwoPN\).

\paragraph{Static \texorpdfstring{$2$}{2}PN gate.}
Quadratic local mass scaling plus the pure-static isotropic test-mass gate fix
\begin{equation}
\lambda_\rho=\frac12,
\label{eq:status-fixed-lambda-rho}
\end{equation}
so the two-body static \(\TwoPN\) block is
\begin{equation}
L_{2,\mathrm{static}}
=\frac{G^3m_A m_B(m_A^2+m_B^2)}{4r^3}.
\label{eq:status-fixed-static-block}
\end{equation}
In the one-body reduction this is exactly the statement that the \(U^3\) coefficient is
fixed to \(+\frac14\).

\paragraph{Comparable-mass residual solve.}
Once the repaired one-body/self/static seed is frozen, the remaining comparable-mass
\(\TwoPN\) discrepancy is fixed uniquely by the exact \(\ADM\) residual solve of
\cref{sec:adm-lift}.
The solved coefficients are
\begin{align}
q_1&=-\frac74,
& q_2&=-\frac14,
& q_3&=\frac{11}{8},
& q_4&=\frac14,
\nonumber\\
q_5&=-\frac58,
& q_6&=\frac32,
& q_7&=\frac38,
\label{eq:status-fixed-qs}
\\[0.5ex]
 t_1&=0,
& t_2&=\frac{11}{8},
& t_3&=-\frac{15}{4},
& t_4&=0,
& t_5&=0,
& t_6&=\frac{15}{8},
\label{eq:status-fixed-ts}
\\[0.5ex]
 s_1&=\frac54.
\label{eq:status-fixed-s1}
\end{align}
Thus the full added comparable-mass cross block is no longer an ansatz.
It is a uniquely solved object.

\paragraph{Full conservative two-body \texorpdfstring{$2$}{2}PN sector.}
After assembly, the complete derived two-body Lagrangian
\cref{eq:full-assembly-Lcons-expanded} satisfies
\begin{equation}
H_2=H^{\ADM}_{2\mathrm{PN}},
\qquad
\Delta H_2\equiv H_2-H^{\ADM}_{2\mathrm{PN}}=0.
\label{eq:status-fixed-full-match}
\end{equation}
So, within the declared closure hierarchy, the entire conservative two-body
\(\TwoPN\) coefficient ledger is fixed.
There are no remaining free dimensionless coefficients in the main theorem.

\subsection{Quantities fixed only within the declared response hierarchy}
\label{sec:status-protocol}

A smaller set of quantities is also fixed in the present paper, but only after one
explicit conservative response hierarchy is declared.
These should be treated as \emph{protocol-fixed}, not as assumption-free theorems of the
fully dynamical moving-throat bulk problem.

\paragraph{Carried breathing/pressure-volume response.}
The \(\OnePN\) closure already fixed
\begin{equation}
\kPV=\frac32,
\qquad
\beta_{\OnePN}=\krho+\kadd+\kPV=3,
\label{eq:status-protocol-kpv-beta}
\end{equation}
within the adiabatic one-degree-of-freedom response protocol.
These values are carried unchanged into the present \(\TwoPN\) paper.
Likewise, the same closure fixes the breathing slope
\begin{equation}
g_1=\frac{57}{64},
\label{eq:status-protocol-g1}
\end{equation}
which is the only geometry-slope datum needed in the main-text one-body
\(\TwoPN\) denominator match.

\paragraph{Minimal conservative one-body packaging.}
The specific invariant ansatz
\begin{equation}
D_{\eff}(u)=C_s(u)\bigl[1+\alpha(G-1)+\beta(G-1)^2\bigr]
\label{eq:status-protocol-deff-ansatz}
\end{equation}
is itself part of the declared low-frequency conservative packaging.
Within that ansatz, preserving the frozen \(\OnePN\) slot forces \(\alpha=0\), and the
exact isotropic one-body target then fixes the quadratic coefficient uniquely.
But the ansatz is still a declared closure move rather than a direct theorem of the full
moving-throat PDE.

\paragraph{Minimal self/static continuation.}
The self-sector Bernoulli exactification and the quadratic local mass-scaling rule are
also part of the stated closure hierarchy.
They are the smallest conservative choices that preserve the frozen lower-order ledger
while generating the universal \(\TwoPN\) one-body/self/static seed.
The exact \(\ADM\) match proves that they are sufficient for the conservative two-body
\(\TwoPN\) theorem, but it does not by itself prove that no richer inner-throat response
exists beyond this minimal packaging.

The cleanest way to summarize the status is therefore the following.
The main conservative coefficient ledger is fixed, but some of the response-side
\emph{explanatory} machinery remains protocol-dependent until the inner-throat
\(\DtN\)/PDE problem is solved directly.

\subsection{Appendix-side throat data already fixed at zero frequency}
\label{sec:status-throat-fixed}

Although the main proof stops at the exact \(\ADM\) match, the appendix program already
shows that the solved \(\TwoPN\) cross block admits a compact throat-response packaging.
Within that packaging, several zero-frequency throat data are no longer free.

\paragraph{Odd-channel static residues.}
The conservative odd dipole sector fixes
\begin{equation}
R_{1\perp}=\frac72,
\qquad
R_{10}=4.
\label{eq:status-fixed-odd-residues}
\end{equation}
These are the zero-frequency transverse and longitudinal dipole residues carried by the
minimal dressed odd channel.

\paragraph{Even support residues.}
The conservative even support sector fixes the canonical zero-frequency support residues
to
\begin{equation}
R_0=R_{20}=R_{21}=R_{22}=1.
\label{eq:status-fixed-even-residues}
\end{equation}
Equivalently, the even support metric is canonically normalized at zero frequency in the
basis used by the appendices.

\paragraph{Scalar source vector and geometry-closure deficit.}
The solved even \(P_0\oplus P_2\) operator fixes the source vector
\begin{equation}
J=\left(\frac{4}{\sqrt5},\frac54,0,0,0,0\right),
\label{eq:status-fixed-J}
\end{equation}
so that
\begin{equation}
J\cdot J=\frac{381}{80}.
\label{eq:status-fixed-JdotJ}
\end{equation}
Here \(J\) is the appendix-side throat-response source-weight vector, not a
Maxwell current.
Because the exact solved static cross coefficient is only \(5/4\), the direct closure
term is fixed to the deficit
\begin{equation}
\Delta_{\geom}=\frac{281}{80}.
\label{eq:status-fixed-deltageom}
\end{equation}
Thus the appendix-side zero-frequency even operator is no longer phenomenological: both
its support/source data and its required pure-\(U\) closure subtraction are fixed by the
solved conservative \(\TwoPN\) cross target.

These quantities are not needed to state the main theorem, but they matter for the
constructive interpretation of the result.
They show that the full conservative \(\TwoPN\) cross sector can already be packaged as
carried odd dipole wake data, a minimal even support/source layer, and one direct
geometry-closure term.

\subsection{Open \texorpdfstring{$\DtN$}{DtN}/PDE observables}
\label{sec:status-open}

What remains open on the inner-throat side is now sharply constrained.
The conservative \(\TwoPN\) algebra fixes the zero-frequency data, but it does not yet
fix the dynamic pole scales of the minimal axisymmetric low-frequency \(\DtN\)
completion.
Those open observables are
\begin{equation}
\Omega_{1\perp},\quad
\Omega_{10},\quad
\Omega_0,\quad
\Omega_{20},\quad
\Omega_{21},\quad
\Omega_{22},\quad
\Omega_g.
\label{eq:status-open-poles}
\end{equation}
Equivalently, in a low-frequency Taylor expansion one may write
\begin{equation}
Y_0(\omega)=1+\chi_0\omega^2+\cdots,
\qquad
Y_2(\omega)=1+\chi_2\omega^2+\cdots,
\label{eq:status-open-chi-even}
\end{equation}
\begin{equation}
Y_{1\perp}(\omega)=\frac72+\chi_{1\perp}\omega^2+\cdots,
\qquad
Y_{10}(\omega)=4+\chi_{10}\omega^2+\cdots,
\label{eq:status-open-chi-odd}
\end{equation}
and the coefficients \(\chi\) are precisely the first genuinely new inner-throat dynamic
observables beyond the conservative zero-frequency \(\TwoPN\) match.

A near-spherical support simplification is still available if desired,
\begin{equation}
\Omega_{20}=\Omega_{21}=\Omega_{22},
\label{eq:status-open-nearspherical}
\end{equation}
but this is optional and is not required anywhere in the main proof.
The important point is that the remaining PDE problem is now narrow and explicit:
one no longer needs to guess the conservative zero-frequency operator, but only to derive
its dynamical completion.

\subsection{Other quantities that remain symbolic}
\label{sec:status-symbolic}

Beyond the open pole data above, several broader structures should still be treated as
symbolic or unsolved.

\paragraph{Absolute dimensional normalizations.}
The present paper fixes dimensionless conservative coefficients, but it does not by itself
fix the absolute dimensional normalization of the parent medium.
In particular, in
\begin{equation}
P(\rho)=\KEOS\rho^n,
\qquad n=5,
\label{eq:status-symbolic-keos}
\end{equation}
the exponent \(n\) is fixed, but the dimensional constant \(\KEOS\) is not.
Likewise, \(G\) and \(c\) enter the reduced description as measured effective constants;
the present paper fixes how they appear in the dimensionless ledger, not their numerical
values from first principles.

\paragraph{Projection/localization profiles and finite-size moments.}
The exact projected identities used by the reduction do not require explicit closed forms
for the detailed projection and localization profiles,
\begin{equation}
\Wproj(w),
\qquad
\Zloc(w),
\label{eq:status-symbolic-profiles}
\end{equation}
and the \(\TwoPN\) proof likewise does not fix explicit finite-size or higher-multipole
data such as worldtube quadrupoles.  These remain profile-dependent and belong to a more
detailed finite-size theory.

\paragraph{General dynamic response beyond the conservative low-frequency closure.}
Outside the specific conservative low-frequency packaging adopted here, the internal
response should remain symbolic.  In particular, one should expect more general data such
as damping, phase lag, compliance, multi-mode mixing, or driven/open-system couplings,
none of which is fixed by the present conservative near-zone \(\TwoPN\) derivation.
The same caution applies to the detailed Family-1 soft-wall / endcap inertia completion:
those constructions strongly constrain the dynamic interpretation, but they are not part
of the main algebraic theorem.

\paragraph{Beyond-conservative-\texorpdfstring{$2$}{2}PN sectors.}
The present paper does not solve radiation reaction, dissipative leakage, spin couplings,
higher post-Newtonian orders, or strong-field non-perturbative completion.
Likewise, the paper does not claim a first-principles solution of the fully dynamical
moving-throat bulk PDE for arbitrary motion.
The worldtube/particle reduction remains a controlled effective description rather than a
complete direct bulk solution.

\subsection{Compact status ledger}
\label{sec:status-ledger-table}

For convenience, \cref{tab:status-ledger-2pn} summarizes the status of the main
quantities entering the present \(\TwoPN\) paper.

\begin{table}[t]
  \centering
  \caption{Status ledger for the main quantities entering the conservative two-body
  \(\TwoPN\) derivation.}
  \label{tab:status-ledger-2pn}
  \small
  \renewcommand{\arraystretch}{1.15}
  \begin{tabularx}{\linewidth}{@{}L L L@{}}
    \toprule
    Quantity & Status & Comments \\
    \midrule
    \(\krho=1\), \(n=5\) & fixed carry-forward & Newtonian and optical lower-order backbone; historical gravity-side \(q=1\) is now written \(\krho=1\) \\
    \(\kadd=1/2\), \(C_\parallel=-7/2\), \(C_L=-1/2\) & fixed carry-forward & frozen topology/wake data from the \(\OnePN\) program \\
    \(\kPV=3/2\), \(\beta_{\OnePN}=3\) & protocol-fixed carry-forward & adiabatic one-DOF breathing closure \\
    \(g_1=57/64\) & protocol-fixed carry-forward & breathing slope used in the one-body \(\TwoPN\) closure \\
    \(D_{\eff}(u)=C_s[1+\tfrac{32768}{3249}(G-1)^2]\) & fixed within hierarchy & exact one-body denominator repair through \(\TwoPN\) \\
    \(\lambda_\rho=1/2\) & fixed & pure-static isotropic gate \\
    \(q_i,t_i,s_1\) of \cref{eq:status-fixed-qs,eq:status-fixed-ts,eq:status-fixed-s1} & fixed & unique comparable-mass residual solve \\
    full assembled \(L_{\rm cons}^{\rm derived}\) & fixed & exact generic-frame conservative \(\ADM\) match through \(\TwoPN\) \\
    \(R_{1\perp}=7/2\), \(R_{10}=4\) & appendix-side fixed & odd zero-frequency throat residues \\
    \(R_0=R_{20}=R_{21}=R_{22}=1\) & appendix-side fixed & even support residues \\
    \(J=(4/\sqrt5,5/4,0,0,0,0)\) & appendix-side fixed & scalar source vector \\
    \(\Delta_{\geom}=281/80\) & appendix-side fixed & pure-\(U\) geometry-closure deficit \\
    \(\Omega_{1\perp},\Omega_{10},\Omega_0,\Omega_{20},\Omega_{21},\Omega_{22},\Omega_g\) & symbolic/open & dynamic \(\DtN\)/PDE pole scales \\
    \(\KEOS\), \(\Wproj(w)\), \(\Zloc(w)\) & symbolic & dimensional normalization and profile data \\
    finite-size moments / higher multipoles & symbolic & beyond the universal monopole \(\TwoPN\) limit \\
    dissipative, radiative, open-system sectors & symbolic & outside the conservative near-zone theorem \\
    \bottomrule
  \end{tabularx}
\end{table}

\subsection{Interpretation}
\label{sec:status-interpretation}

The cleanest way to summarize the paper is therefore the following.

\begin{quote}
Within the declared closure hierarchy, the conservative two-body \(\TwoPN\) sector of
the toy model is fully fixed and matches the standard generic-frame conservative
\(\ADM\) target exactly.  What remains symbolic is no longer the main coefficient ledger
itself, but absolute normalization, generalized dynamic throat response, finite-size
structure, open-system or dissipative physics, and the fully solved moving-throat bulk
PDE completion.
\end{quote}

That is a stronger statement than was available at \(\OnePN\).
At \(\OnePN\), the main open items still included the comparable-mass \(\TwoPN\)
cross content and its throat-side constructive packaging.
After the present paper, the remaining problem has narrowed to a much more specific one:
to derive from the inner-throat dynamics the pole structure and dynamic completion of a
zero-frequency operator whose conservative \(\TwoPN\) coefficient ledger is already known.

% ============================================================
\section{Verification and reproducibility ledger}
\label{sec:verification}

The derivation in this paper is analytic rather than computational.  The accompanying
Wolfram Language (WL) files and symbolic derivation notes are therefore not presented as
black-box generators of the final result.  Their role is narrower and more auditable:
they act as a referee archive for the algebraic identities, carried regressions,
coefficient solves, and end-to-end assembly statements used in the main text.

At \(\TwoPN\), the verification archive is split into two layers.
The first layer consists of compact executable harnesses, whose job is to check the
frozen lower-order ledger, the one-body closure, the self/static seed, the test-mass
target, and the finite-size gate in explicit PASS/FAIL form.
The second layer consists of symbolic derivation ledgers, whose job is to record the
comparable-mass \(\ADM\) residual solve and the appendix-side throat-response
reconstruction in a way that can be rechecked coefficient by coefficient.
This split reflects the actual structure of the project archive: the compact self/static
and one-body checks are already bundled as executable scripts, whereas the heavier
comparable-mass and throat-response modules are currently documented as exact symbolic
solves rather than as a single released WL harness.

This distinction matters for interpretation.
The archive verifies that the coefficient ledger derived in the previous sections is
internally consistent \emph{within the stated closure hierarchy}.  It does not replace
any analytic step, and it does not upgrade the paper into an assumption-free theorem of
a fully solved moving-throat bulk dynamics.  Its purpose is reproducibility and
auditability: every load-bearing statement of the conservative \(\TwoPN\) assembly can be
checked either by a named executable test or by an explicit symbolic residual recorded in
the derivation ledger.

\subsection{Structure of the verification archive}
\label{sec:verification-archive}

The verification archive used for the present paper is summarized in
\cref{tab:verification-2pn}.  It has three layers: the inherited \(\OnePN\) freeze, the
compact executable \(\TwoPN\) harness, and the symbolic \(\TwoPN\) derivation ledger.

\begin{table}[t]
  \centering
  \caption{Verification archive used for the conservative two-body \(\TwoPN\) paper.}
  \label{tab:verification-2pn}
  \small
  \renewcommand{\arraystretch}{1.15}
  \begin{tabularx}{\linewidth}{@{}L L L@{}}
    \toprule
    Artifact & Role in the derivation & Reported summary \\
    \midrule
    \nolinkurl{4d_gravity_1pn_master_harness.wl}\\
    \nolinkurl{vector_wake_rebuild.wl}\\
    \nolinkurl{4d_full_1pn_derivation_with_vectorwake.wl}
    &
    Inherited \(\OnePN\) referee suite.  Freezes the lower-order backbone used unchanged
    in the present paper: Newtonian mass-dressing coefficient \(\krho=1\), optical matching \(n=5\),
    \(\kadd=1/2\), the adiabatic \(\kPV=3/2\) closure, \(\beta_{\OnePN}=3\), the
    parity-even wake coefficients, and the exact conservative EIH equality.
    &
    Reference transcripts report \texttt{PASS count = 60}, \texttt{FAIL count = 0}
    for the broad master harness; the wake module recovers
    \(a_H=0\), \(a_L=\pm\sqrt3/2\), \(K_{\rm vec}=2/\pi^2\),
    \(C_{\parallel}=-7/2\), \(C_L=-1/2\); and the full \(\OnePN\) assembly script
    reports \texttt{Passes: 37}, \texttt{Fails: 0}, \texttt{Skips: 0}.
    \\[4pt]

    \nolinkurl{4d_gravity_2pn_master_harness.wl}
    &
    Compact executable \(\TwoPN\) referee harness.  Re-runs the carried \(\OnePN\)
    regressions; verifies the exact \(n=5\) Bernoulli continuation; checks the minimal
    self-sector exactification, the quadratic local-mass-scaling static hierarchy, the
    generic cubic adiabatic-elimination template, the self/static test-mass reduction,
    the isotropic Schwarzschild target coefficients, and the universal finite-size gate.
    &
    Reference transcript reports
    \texttt{Passes: 42}, \texttt{Fails: 0}, \texttt{Skips: 0},
    with an overall all-pass summary.
    \\[4pt]

    \nolinkurl{4d_2pn_full_notes.md}
    &
    Symbolic \(\TwoPN\) derivation ledger for the exact comparable-mass \(\ADM\) lift
    and for the appendix-side constructive throat-response program.
    Records the perturbative Legendre identity used for the lift, the compact residual
    basis, the unique solved coefficients \(q_i\), \(t_i\), \(s_1\), the zero-residual
    full \(\ADM\) match, and the exact zero-frequency throat-module reconstruction.
    &
    Reports an exact generic-frame \(\ADM\) match for the assembled conservative
    \(\TwoPN\) block and records that the final low-frequency throat module has a
    symbolic residual that vanishes identically against the solved added cross target.
    This layer is presently a symbolic derivation ledger rather than a single released
    PASS/FAIL WL harness.
    \\
    \bottomrule
  \end{tabularx}
\end{table}

The division of labor among these artifacts mirrors the logic of the paper itself.
The inherited \(\OnePN\) suite freezes the lower-order ledger that the present paper is
not allowed to retune.
The compact \(\TwoPN\) harness checks the one-body, self, static, and test-mass parts of
the proof in explicit executable form.
The comparable-mass \(\ADM\) lift and the appendix-side throat-response modules are then
made reproducible by writing their residual basis and solved coefficient set explicitly,
so that the final match can be checked coefficient by coefficient rather than trusted as
an opaque notebook output.

\subsection{What is actually being checked}
\label{sec:verification-checked}

The archive is designed around the same claim taxonomy used in the main text, and the
checks naturally fall into five groups.

\paragraph{(i) Frozen lower-order regressions.}
Before any \(\TwoPN\) step is accepted, the carried lower-order ledger is rechecked.
The executable archive verifies, among other items,
\begin{equation}
\krho=1,
\qquad
n=5,
\qquad
\kadd=\frac12,
\qquad
\kPV=\frac32,
\qquad
\beta_{\OnePN}=3,
\label{eq:verification-frozen-ledger}
\end{equation}
together with the one-body scalar/optical \(\OnePN\) self coefficient and the pair-counted
static \(\OnePN\) coefficient.  This matters because the present \(\TwoPN\) paper is not
permitted to improve its target by silently shifting any part of the already-solved
Newtonian+\(\OnePN\) backbone.

\paragraph{(ii) One-body and self/static \texorpdfstring{$2$}{2}PN checks.}
The compact \(\TwoPN\) harness verifies the exact \(n=5\) Bernoulli continuation,
including the optical-index expansion
\begin{equation}
N(U)=1+2\frac{U}{c^2}+6\frac{U^2}{c^4}+O(c^{-6}),
\label{eq:verification-bernoulli-index}
\end{equation}
and then checks the minimal self-sector exactification,
\begin{equation}
\frac{L_{\mathrm{self},\min}}{m}
=-c^2+\frac12 v^2+U
+\frac{1}{c^2}\left(\frac18 v^4+\frac32 Uv^2\right)
+\frac{1}{c^4}\left(\frac1{16}v^6+\frac78 Uv^4+6U^2v^2\right)+O(c^{-6}),
\label{eq:verification-self-minimal}
\end{equation}
including the specific coefficients \(1/16\), \(7/8\), and \(6\).
It also verifies the cautionary check that a fully Bernoulli-exactified mass prefactor
would generate a forbidden \(+3/2\) static coefficient already at \(\OnePN\).

\paragraph{(iii) Static hierarchy and response-template checks.}
The same harness verifies that quadratic local mass scaling yields
\begin{equation}
L_{2,\mathrm{static}}
=\frac{\lambda_\rho G^3m_A m_B(m_A^2+m_B^2)}{2r^3},
\label{eq:verification-static-general}
\end{equation}
checks selected unit-normalized three-body and four-body cubic-static monomials, and
verifies the generic one-fast-mode adiabatic-elimination template through cubic order.
These checks do not close the full throat dynamics, but they certify that the exact algebra
used to package the self/static seed and the one-degree-of-freedom response correction is
internally consistent.

\paragraph{(iv) Test-mass target and regime checks.}
The compact executable harness confirms that the two-body self/static reduction agrees
exactly with the direct one-body construction, that the isotropic Schwarzschild target has
its required static and mixed coefficients, and that the pure-static isotropic gate fixes
\begin{equation}
\lambda_\rho=\frac12.
\label{eq:verification-lambda-rho}
\end{equation}
It then checks that the remaining self/static discrepancy is exactly
\begin{equation}
\Delta L_{\mathrm{self+static}\rightarrow\mathrm{iso}}
=4U^2v^2,
\label{eq:verification-residual-4u2v2}
\end{equation}
which is precisely the missing one-body slot repaired in \cref{sec:one-body}.
Finally, it verifies the universal finite-size gate
\begin{equation}
\frac{F_{\rm finite\text{-}size}^{(2)}}{F_{\rm universal}^{(2)}}
=\left(\frac{a c^2}{GM}\right)^2,
\label{eq:verification-finite-size-gate}
\end{equation}
so the regime of validity of the universal conservative \(\TwoPN\) force is made explicit.

\paragraph{(v) Comparable-mass \texorpdfstring{$\ADM$}{ADM} lift and appendix-side constructive checks.}
The exact comparable-mass lift is verified in a different style.
The symbolic derivation ledger records the perturbative Legendre identity
\begin{equation}
H_1=-L_1(v_0),
\qquad
H_2=-L_2(v_0)+\frac12 A_0^T M^{-1}A_0,
\label{eq:verification-legendre-identity}
\end{equation}
with \(v_0=p/m\) and \(A_0=(\partial L_1/\partial v)_{v=v_0}\), then writes the residual
in an explicit invariant basis and solves it uniquely for
\(q_i\), \(t_i\), and \(s_1\).
Because the basis and coefficients are given in closed form, the full Hamiltonian residual
can be rechecked term by term.
Likewise, the appendix-side constructive program records an explicit final low-frequency
throat module whose symbolic residual against the solved added conservative
\(\TwoPN\) cross target vanishes identically.
These are not presently packaged as one compact released PASS/FAIL script, but they are
still reproducible because the residual basis, coefficient ledger, and zero-residual
statements are written out explicitly.

\subsection{Reference run summaries}
\label{sec:verification-summaries}

For the executable layers of the archive, the PASS/FAIL status can be read off directly
from the reference transcripts.
The inherited \(\OnePN\) master harness and end-to-end assembly script report
\begin{equation}
\texttt{PASS count = 60},
\qquad
\texttt{FAIL count = 0},
\qquad
\texttt{Passes: 37},
\qquad
\texttt{Fails: 0},
\qquad
\texttt{Skips: 0},
\label{eq:verification-1pn-counts}
\end{equation}
while the constructive wake module reports the real-parameter solution
\begin{equation}
a_H=0,
\qquad
a_L=\pm\frac{\sqrt3}{2},
\qquad
K_{\rm vec}=\frac{2}{\pi^2},
\qquad
C_{\parallel}=-\frac72,
\qquad
C_L=-\frac12.
\label{eq:verification-1pn-wake}
\end{equation}

For the compact executable \(\TwoPN\) layer, the reference transcript reports
\begin{equation}
\texttt{Passes: 42},
\qquad
\texttt{Fails: 0},
\qquad
\texttt{Skips: 0},
\label{eq:verification-2pn-counts}
\end{equation}
and terminates with an explicit all-pass summary.
In other words, every executable check currently bundled with the \(\TwoPN\) project
snapshot closes cleanly.

The symbolic derivation ledger does not report a single aggregate pass count, because it is
organized as an exact coefficient solve rather than as a compact transcripted harness.
Its reproducibility metadata is instead the sequence of exact outputs written in the notes:
the unique comparable-mass coefficients \((q_i,t_i,s_1)\), the full assembled
\(L_{2,\mathrm{full}}\), the exact generic-frame conservative \(\ADM\) equality
\begin{equation}
H_2=H^{\ADM}_{2\mathrm{PN}},
\label{eq:verification-adm-zero}
\end{equation}
and the final appendix-side statement that the low-frequency throat module has a symbolic
residual that vanishes identically against the solved added cross target.
These should be read as symbolic reproducibility metadata rather than as substitutes for
the derivation itself.

\subsection{What the verification archive does \emph{not} verify}
\label{sec:verification-scope}

It is equally important to state the scope limits of the archive.

First, neither the executable harnesses nor the symbolic derivation notes solve the fully
dynamical moving-throat bulk PDE for arbitrary motion.
The paper remains a conservative near-zone \(\TwoPN\) derivation within a declared closure
hierarchy, built on the same controlled worldtube and response reductions used elsewhere
in the project.

Second, the compact executable \(\TwoPN\) harness is intentionally narrower than the full
paper.
It covers the frozen \(\OnePN\) regressions, the Bernoulli continuation, the self/static
seed, the isotropic one-body target, and the finite-size gate.
The exact comparable-mass \(\ADM\) lift and the appendix-side throat-response packaging
are presently documented as symbolic ledgers rather than as one public all-in-one WL
script.
This is a statement about how the archive is currently organized, not about any missing
analytic step in the paper itself.

Third, the archive verifies only the conservative sector addressed in the present paper.
It does not claim to cover dissipative leakage, radiation reaction, spin couplings,
higher-PN orders, strong-field completion, or non-adiabatic internal dynamics beyond the
declared response closures.

In short, the archive verifies the paper \emph{as stated}, not a stronger claim than the
paper itself makes.

\subsection{Reproducibility conventions}
\label{sec:verification-conventions}

Four reproducibility conventions are worth recording explicitly.

\paragraph{Named executable checks.}
Every nontrivial statement in the compact harnesses is given a human-readable PASS/FAIL
name.  This makes it possible to compare the paper's prose directly with the verification
transcript.

\paragraph{Local assumptions.}
The executable checks do not hide their regime assumptions globally.
Quasi-staticity, positivity domains, small-body ordering, and the declared response
closures are attached to the relevant symbolic reductions, so the reader can see which
equalities are exact and which are conditional.

\paragraph{Explicit residual bases.}
For the comparable-mass lift and the appendix-side throat modules, reproducibility comes
from explicit residual bases and explicit solved coefficient ledgers.
The point is not to hide a large notebook behind a single ``zero'' output, but to expose
the exact polynomial basis on which the residual is being solved.

\paragraph{Modular audit trail.}
The full verification story is intentionally modular.
Lower-order freeze, one-body closure, self/static seed, comparable-mass lift, and
constructive throat packaging are kept separate so that failure in one sector cannot be
masked by compensating changes in another.
This is especially important at \(\TwoPN\), where it would otherwise be easy to blur
together one-body repairs, static bookkeeping, and genuinely new comparable-mass cross
content.

For readability, the main text reports only the archive structure, the all-pass summaries,
and the zero-residual assembly statements.
A longer PASS/FAIL transcript or a future consolidated \(\TwoPN\) WL harness can be added
later as a repository note or supplementary appendix, but the present paper already
contains enough metadata for a referee to audit the derivation sector by sector.

\paragraph{Bottom line.}
The purpose of the verification archive is not to ask the reader to trust a computation
instead of a derivation.
It is to provide an auditable backstop showing that the carried and newly introduced
coefficients close exactly under the stated assumptions:
\begin{equation}
D_{\eff}(u)=1-4u+8u^2+O(u^3),
\qquad
\lambda_\rho=\frac12,
\qquad
H_2=H^{\ADM}_{2\mathrm{PN}}.
\label{eq:verification-bottom-line-2pn}
\end{equation}
The appendix-side constructive program then goes further and shows that the solved added
cross sector can already be packaged into a minimal low-frequency throat-response module
with zero symbolic residual.

\claimclass{Documentary / reproducibility note.  This section does not add a new dynamical
proof.  It records how the inherited \(\OnePN\) suite, the compact executable \(\TwoPN\)
harness, and the symbolic \(\TwoPN\) derivation ledger jointly verify the exact
identities, controlled reductions, protocol closures, and full-match assembly statements
used elsewhere in the paper.}

% ============================================================
\section{Conclusions}
\label{sec:conclusion}

\subsection{What has been proven}
\label{sec:conclusion-proven}

The main result of this paper is that the \(4\)D toy model now supports a full
\emph{conservative two-body \(2\)PN derivation} within a clearly stated closure
hierarchy.
The point of the paper is not merely that the model can be tuned to resemble the
standard generic-frame \(\ADM\) Hamiltonian through order \(c^{-4}\).
Rather, once the derivation is organized by sector and the already-proved
Newtonian+\(\OnePN\) ledger is frozen, every coefficient in the conservative
two-body \(\TwoPN\) block can be traced to a named ingredient of the parent
\(4\)D program or to one exact residual solve whose basis is fully explicit.

The proof chain established here is:
\begin{enumerate}[leftmargin=1.5em]
  \item exact projection identities, the quasi-static Poisson hook, and the
  controlled coherent-defect / small-body worldtube reduction inherited from the
  earlier \(4\)D papers;
  \item the frozen Newtonian+\(\OnePN\) ledger, including the exact EIH match, which
  the present paper is not allowed to retune;
  \item a minimal one-body \(\DtN\)-invariant denominator closure,
  \begin{equation}
  D_{\eff}(u)=C_s(u)\left[1+\frac{32768}{3249}(G(u)-1)^2\right]
  =1-4u+8u^2+O(u^3),
  \label{eq:conclusion-deff}
  \end{equation}
  which closes the missing test-mass \(\TwoPN\) response slot while preserving the
  already-frozen \(\OnePN\) coefficient;
  \item a minimal self/static \(\TwoPN\) seed consisting of Bernoulli exactification in
  the self sector together with quadratic local mass scaling in the static sector,
  the latter being fixed by the pure-static isotropic gate to
  \begin{equation}
  \lambda_\rho=\frac12;
  \label{eq:conclusion-lambda-rho}
  \end{equation}
  \item an exact comparable-mass \(\ADM\) lift, which solves the remaining
  quartic-velocity, mixed \(G^2/r^2\), and static-cross residual in a compact,
  symmetry-respecting invariant basis; and
  \item final assembly, for which the full derived conservative two-body Lagrangian
  satisfies
  \begin{equation}
  L_{\rm cons}^{\rm derived}
  =L_0+\frac{1}{c^2}L_1+\frac{1}{c^4}L_{2,\mathrm{full}}+O(c^{-6}),
  \qquad
  H_2=H_{2\mathrm{PN}}^{\ADM}.
  \label{eq:conclusion-main-match}
  \end{equation}
\end{enumerate}

That last equality is the paper's central \FullMatch statement.
Within the stated closure hierarchy, the conservative two-body \(\TwoPN\) coefficient
ledger is no longer open.
The assembled answer is not just correct in the strict test-mass limit and not just
compatible with a change of gauge; its perturbative Legendre transform is exactly the
standard generic-frame conservative \(\ADM\) target through \(\TwoPN\).

Conceptually, the paper does three things at once.
First, it upgrades the earlier program from a full conservative \(\OnePN\) derivation to
full conservative \(\TwoPN\) order.
Second, it keeps that upgrade readable by separating the proof line from the larger
constructive throat-response story: the main text proves the \(\TwoPN\) theorem, while
the appendices record the tensor-wake, mouth-port, support-minus-closure, PDE-side
\(\DtN\), and Family-1 developments that explain how the new comparable-mass sector may
be understood from the inner-throat side.
Third, it sharpens the bookkeeping of what is fixed, what is protocol-fixed, and what
should still remain symbolic.

In that sharpened sense, the paper closes more than one algebraic gap.
It shows that the missing one-body \(\TwoPN\) response slot can be repaired without
spoiling the frozen \(\OnePN\) ledger; it shows that the pure-static one-body target
selects a unique quadratic local-mass-scaling coefficient; and it shows that the
remaining comparable-mass discrepancy is not mysterious but a finite linear residual
solve with a unique answer.
Once these points are separated cleanly, the final conservative \(\TwoPN\) assembly is
surprisingly short.

As at \(\OnePN\), the status labels matter.
Some steps are \Exact identities or exact algebraic residual solves.
Some are \Controlled reductions, such as the Poisson hook and the worldtube particle
limit.
Some are \Protocol closures, such as the low-frequency invariant one-body packaging and
minimal self/static response hierarchy.
The final equality \cref{eq:conclusion-main-match} is then a \FullMatch statement once
those previously declared ingredients are inserted.

That is the correct strength of the result.
This paper does \emph{not} claim an assumption-free theorem derived from a fully solved
moving-throat bulk background for arbitrary motion.
It establishes a full conservative two-body \(\TwoPN\) derivation \emph{within a stated
closure hierarchy}.
Within that hierarchy, however, there are no remaining free dimensionless coefficients in
the main conservative two-body \(\TwoPN\) ledger.

\subsection{What remains open}
\label{sec:conclusion-open}

The scope of the result should be kept equally explicit on the negative side.

First, the paper does not solve the fully dynamical moving-throat PDE from first
principles and then derive the particle limit as an automatic theorem.
The worldtube / coherent-defect reduction remains a controlled effective description.
Likewise, the one-body denominator closure and the minimal self/static continuation are
not claimed as universal facts of arbitrary throat dynamics; they are the minimal
conservative closures used in the present paper and are justified precisely by the way
they preserve the frozen lower-order ledger and close the exact \(\TwoPN\) target.

Second, the paper does not claim that the appendix-side throat-response reconstruction is
already dynamically complete.
What the appendices show is that the solved conservative \(\TwoPN\) cross block admits a
compact zero-frequency packaging in terms of the carried odd dipole wake, a new even
\(P_0\oplus P_2\) support layer, and a pure-potential geometry-closure term.
At that level, important static data are already fixed.
But the dynamic pole scales remain genuine PDE/\(\DtN\) observables.
A convenient representative list is
\begin{equation}
\Omega_{1\perp},\quad
\Omega_{10},\quad
\Omega_0,\quad
\Omega_{20},\quad
\Omega_{21},\quad
\Omega_{22},\quad
\Omega_g.
\label{eq:conclusion-open-poles}
\end{equation}
These are the first genuinely new inner-throat dynamic observables after the
conservative zero-frequency \(\TwoPN\) coefficient ledger has been fixed.

Third, the present paper is restricted to the conservative near-zone \(\TwoPN\) sector.
It does not address dissipative leakage, radiation reaction, spin couplings, strong-field
nonperturbative completion, or the general finite-size / higher-multipole problem beyond
the universal monopole regime.
The finite-size gate derived in the main text and appendices should therefore be read as
part of the honest scope statement rather than as a hidden footnote: the universal
conservative \(\TwoPN\) ledger applies in the declared small-body regime, not as a claim
about arbitrary extended objects.

These non-claims do not weaken the main theorem.
They tell us exactly where the genuine remaining physics now sits.
Before the present paper, it was still unclear whether the conservative two-body
\(\TwoPN\) ledger could be closed cleanly at all.
After the present paper, that algebraic question is settled.
What remains open is the dynamical throat explanation of already-solved conservative data,
not the existence of the data themselves.

\subsection{Interface to the next paper}
\label{sec:conclusion-next}

This gives the project a much cleaner handoff than before.
The next true job is no longer to continue coefficient matching at conservative
zero-frequency \(\TwoPN\) order.
That work is finished.
The next job is to derive the appendix-side throat data directly from the inner-throat
PDE / \(\DtN\) problem and to turn the present static and low-frequency packaging into a
fully dynamical response module.

Put differently, the conservative zero-frequency target has now been promoted from a
heuristic guide to a precise construction target.
The isotropic cavity has become a useful no-go unit test rather than a vague model.
The minimal axisymmetric channel structure is known.
The zero-frequency support/source ledger is known.
And even the pure-geometry closure deficit is known.
So the next paper can focus on a narrower and more physical question:
can the Family-1 throat dynamics, with its wall and endcap structure treated honestly,
derive the open pole data and dynamic response functions whose static limits are already
fixed here?

That is a qualitatively better situation than at the end of the \(\OnePN\) paper.
At \(\OnePN\), one still had to discover the new comparable-mass \(\TwoPN\) cross block
and determine whether it admitted any constructive throat-response interpretation.
After the present paper, the conservative \(\TwoPN\) algebra is complete and the appendix
program already shows that the answer is yes: the new cross content does admit a compact
throat-response packaging.
The remaining task is therefore to derive its dynamics, not to guess its static form.

The most natural immediate follow-ups are then clear.
One is a dedicated dynamical throat-response paper, centered on the pole structure,
soft-wall / surface inertia completion, and the low-frequency \(\DtN\) operator.
Another is a more systematic finite-size and higher-multipole extension of the worldtube
reduction.
A third is the move beyond conservative near-zone physics toward dissipation, radiation,
and higher post-Newtonian order.
But those are now properly sequenced follow-ups rather than unfinished business inside the
present paper.

What has been shown here is that the \(4\)D toy model can reproduce the full conservative
two-body \(\TwoPN\) ledger in a way that is sector-by-sector transparent, algebraically
exact at final assembly, and explicit about its remaining scope conditions.
That is the sense in which the project has moved from a successful \(\OnePN\) derivation
to a genuine conservative \(\TwoPN\) proof.

\claimclass{\FullMatch as a summary statement of the paper's central result, with the
important scope condition that the derivation is complete within the stated closure
hierarchy rather than as an assumption-free theorem of the fully solved parent \(4\)D
dynamics.}

% ============================================================
\appendix

\section{Notation and conventions}
\label{app:notation}

This appendix records the notation and bookkeeping conventions used throughout the
paper.  It is intentionally shorter than the corresponding appendix in the
\(\OnePN\) paper.  The goal here is only to collect the symbols and conventions that
actually enter the conservative two-body \(\TwoPN\) derivation and the appendices that
support it.

The main text is organized by derivation sector rather than by symbol class, so the same
quantities reappear in several different roles: as exact carry-forward data, as
low-frequency response invariants, as coefficients in the generic-frame two-body
Lagrangian, and as appendix-side throat-response parameters.  The purpose of the present
appendix is to keep those uses unambiguous and to state explicitly which equalities are
literal definitions, which are reduction conventions, and which are claim-class labels.

\begin{table}[t]
  \centering
  \caption{Core notation used repeatedly in the main text and appendices.}
  \label{tab:notation-core-2pn}
  \small
  \renewcommand{\arraystretch}{1.15}
  \begin{tabularx}{\linewidth}{@{}L L@{}}
    \toprule
    Symbol & Meaning \\
    \midrule
    \(X^M=(t,x,y,z,w)=(t,\x,w)\) & parent spacetime coordinates \\
    \(\X=(x,y,z,w)\in\R^4\) & bulk spatial position \\
    \(\x=(x,y,z)\in\R^3\) & brane spatial position \\
    \(m_A,m_B\) & two-body masses \\
    \(\mathbf x_A,\mathbf x_B\) & body positions in a generic frame \\
    \(\vct_A,\vct_B\) & coordinate-time velocities, \(\vct_A=\dot{\mathbf x}_A\), \(\vct_B=\dot{\mathbf x}_B\) \\
    \(\mathbf r=\mathbf x_A-\mathbf x_B\), \(r=|\mathbf r|\), \(\nvec=\mathbf r/r\) & pair separation vector, magnitude, and unit direction \\
    \(v_A^2, v_B^2, v_{AB}, v_{An}, v_{Bn}\) & standard velocity invariants used in the generic-frame \(\TwoPN\) basis \\
    \(U_A=Gm_B/r\), \(U_B=Gm_A/r\) & local Newtonian pair potentials \\
    \(U=GM/r\), \(u=U/c^2\) & one-body Newtonian potential and its dimensionless PN variable \\
    \(L_0,L_1,L_2\) & Newtonian, \(\OnePN\), and \(\TwoPN\) Lagrangian blocks in
    \(L_{\rm cons}=L_0+c^{-2}L_1+c^{-4}L_2+O(c^{-6})\) \\
    \(H_0,H_1,H_2\) & corresponding Hamiltonian blocks after perturbative Legendre transform \\
    \(\ADM\), \(\EIH\), \(\Schw\), \(\DtN\) & standard abbreviations used throughout the paper \\
    \(C_s(u),G(u)\) & normalized low-frequency one-body \(\DtN\) invariants \\
    \(q_i,t_i,s_1\) & solved \(\ADM\)-lift coefficients of the added comparable-mass \(\TwoPN\) cross block; unrelated to electric charge \\
    \(R_{1\perp},R_{10},R_0,R_{20},R_{21},R_{22},J,\Delta_{\geom}\) & appendix-side zero-frequency throat-response data; \(J\) here is a throat-response source-weight vector, not a Maxwell current \\
    \bottomrule
  \end{tabularx}
\end{table}

\subsection{Coordinates, pair variables, and PN bookkeeping}
\label{app:notation-bookkeeping}

We work in a parent spacetime with coordinates
\begin{equation}
X^M=(t,x,y,z,w)\equiv (t,\x,w),
\qquad
\X=(x,y,z,w)\in\R^4,
\qquad
\x=(x,y,z)\in\R^3.
\label{eq:notation-2pn-coordinates}
\end{equation}
Here \(\X\) denotes the full spatial position in the parent \(4\)-space, whereas
\(\x\) denotes the ordinary brane spatial position.  Except where appendix-side
carry-forward formulas explicitly refer back to the parent construction, the main paper is
written entirely in terms of brane-observed effective particle variables.

Body labels \(A,B,C,\dots\) are never summed implicitly.  By contrast, ordinary spatial
Cartesian components are summed in the usual Euclidean way when dot products are written in
index-free form.  An overdot always denotes differentiation with respect to coordinate time
\(t\):
\begin{equation}
\vct_A\equiv \dot{\mathbf x}_A,
\qquad
\vct_B\equiv \dot{\mathbf x}_B.
\label{eq:notation-2pn-dots}
\end{equation}
All dot products such as \(\vct_A\!\cdot\!\vct_B\) and
\(\nvec\!\cdot\!\vct_A\) are ordinary Euclidean \(3\)-dimensional brane dot products.

For the two-body problem we define
\begin{equation}
\mathbf r\equiv \mathbf x_A-\mathbf x_B,
\qquad
r\equiv |\mathbf r|,
\qquad
\nvec\equiv \frac{\mathbf r}{r},
\label{eq:notation-2pn-rn}
\end{equation}
and use the invariant velocity combinations
\begin{equation}
v_A^2\equiv \vct_A\!\cdot\!\vct_A,
\qquad
v_B^2\equiv \vct_B\!\cdot\!\vct_B,
\qquad
v_{AB}\equiv \vct_A\!\cdot\!\vct_B,
\label{eq:notation-2pn-velocity-invariants-1}
\end{equation}
\begin{equation}
v_{An}\equiv \vct_A\!\cdot\!\nvec,
\qquad
v_{Bn}\equiv \vct_B\!\cdot\!\nvec.
\label{eq:notation-2pn-velocity-invariants-2}
\end{equation}
These are the only kinematic invariants needed for the generic-frame conservative
\(\TwoPN\) Lagrangian basis used in the main text.

The local Newtonian pair potentials are written
\begin{equation}
U_A\equiv \frac{Gm_B}{r},
\qquad
U_B\equiv \frac{Gm_A}{r}.
\label{eq:notation-2pn-local-potentials}
\end{equation}
In the strict one-body reduction, where body \(B\) is promoted to a fixed heavy source,
we write
\begin{equation}
m_A\to m,
\qquad
m_B\to M,
\qquad
\vct_B=0,
\qquad
U\equiv \frac{GM}{r},
\qquad
u\equiv \frac{U}{c^2}.
\label{eq:notation-2pn-one-body}
\end{equation}
The dimensionless variable \(u\) is the natural weak-field bookkeeping variable in the
one-body \(\DtN\) closure and isotropic Schwarzschild comparisons.

We keep both \(G\) and \(c\) explicit throughout the paper; no geometric-unit convention
\((G=c=1)\) is adopted.  The post-Newtonian counting is standard:
\begin{equation}
\frac{v^2}{c^2}\sim \frac{U}{c^2}\sim O(c^{-2}).
\label{eq:notation-2pn-counting}
\end{equation}
Accordingly, the conservative two-body Lagrangian is expanded as
\begin{equation}
L_{\rm cons}=L_0+\frac{1}{c^2}L_1+\frac{1}{c^4}L_2+O(c^{-6}),
\label{eq:notation-2pn-lagrangian-expansion}
\end{equation}
where \(L_0\) is the Newtonian block, \(L_1\) the full frozen conservative \(\OnePN\)
block, and \(L_2\) the full conservative \(\TwoPN\) block derived in this paper.
Likewise,
\begin{equation}
H_{\rm cons}=H_0+\frac{1}{c^2}H_1+\frac{1}{c^4}H_2+O(c^{-6})
\label{eq:notation-2pn-hamiltonian-expansion}
\end{equation}
denotes the perturbatively Legendre-transformed Hamiltonian.

When canonical variables are needed in the \(\ADM\) comparison, we write
\(\pvec_A\) and \(\pvec_B\) for the momenta conjugate to
\(\mathbf x_A\) and \(\mathbf x_B\).  As in the main text, the perturbative Legendre
transform is performed in the generic frame rather than in a center-of-mass-specialized
basis.

The paper uses the standard Lagrangian sign convention
\begin{equation}
L=T-V.
\label{eq:notation-2pn-LTV}
\end{equation}
Because the Newtonian pair potential energy is
\begin{equation}
V_{AB}=-\frac{Gm_A m_B}{r},
\label{eq:notation-2pn-VAB}
\end{equation}
the Newtonian interaction term appears in the Lagrangian as
\begin{equation}
L_{\rm int}^{(AB)}=+\frac{Gm_A m_B}{r}.
\label{eq:notation-2pn-Lint}
\end{equation}
This positive Newtonian pair term is used uniformly throughout the paper and its
appendices.

Two further bookkeeping conventions are worth making explicit.
First, unless stated otherwise, equalities between displayed Lagrangians in the main text
are literal equalities in the chosen generic-frame representative, not merely equivalences
modulo an unspecified contact transformation.  Second, when the paper states an “exact
\(\ADM\) match,” the precise claim is that the perturbative Legendre transform of the
assembled Lagrangian satisfies
\begin{equation}
H_2=H_{2\mathrm{PN}}^{\ADM}
\label{eq:notation-2pn-exact-adm}
\end{equation}
in the standard generic-frame conservative target basis.

Finally, the appendix-side throat-response sections use the labels
\begin{equation}
R_{1\perp},\ R_{10},\ R_0,\ R_{20},\ R_{21},\ R_{22},\ J,\ \Delta_{\geom},
\label{eq:notation-2pn-throat-symbols}
\end{equation}
for zero-frequency odd and even mouth-response data, source weights, and the
pure-potential geometry-closure deficit, and
\begin{equation}
\Omega_{1\perp},\ \Omega_{10},\ \Omega_0,\ \Omega_{20},\ \Omega_{21},\ \Omega_{22},\ \Omega_g
\label{eq:notation-2pn-pole-symbols}
\end{equation}
for the corresponding dynamical pole scales.  These quantities are used only in the
appendix-side constructive program and are not needed to read the main \(\TwoPN\) proof.
In particular, the symbol \(J\) in this appendix notation is a throat-response
source-weight vector rather than a Maxwell current.

\subsection{Claim-class labels}
\label{app:notation-claims}

The paper uses four claim-class labels to distinguish exact statements from controlled
reductions, declared closures, and final end-to-end matches.  These labels are part of the
scientific bookkeeping of the paper, not merely an expository convenience.

\begin{description}[leftmargin=1.8em,style=nextline]
  \item[\Exact]
  Statements that follow directly from definitions, exact algebra, exact projection
  identities already established in the parent theory, or exact coefficient solves once a
  residual basis has been declared.  In the present paper, the perturbative Legendre
  algebra and the final residual solve for the added comparable-mass block are typical
  examples.

  \item[\Controlled reduction]
  Statements that require an explicit regime assumption or reduction hierarchy, such as the
  quasi-static Poisson hook, the coherent-defect / small-body worldtube reduction, weak-field
  truncation, or the universal small-body regime used in the self/static \(\TwoPN\)
  sector.  These are not arbitrary guesses, but they are not exact theorems of the fully
  dynamical parent bulk problem either.

  \item[\Protocol closure]
  Statements that depend on an explicitly declared internal-response or operator closure.
  In the present paper, this includes the adiabatic breathing-response input inherited from
  the \(\OnePN\) program and the minimal low-frequency invariant packaging used to close
  the one-body missing \(\TwoPN\) denominator slot.

  \item[\textbf{Full assembly}]
  End-to-end equalities that become exact once all previously declared exact ingredients,
  controlled reductions, and protocol closures are inserted.  The primary example in the
  present paper is the final statement that the assembled conservative two-body
  \(\TwoPN\) Lagrangian Legendre-transforms exactly to the standard generic-frame
  conservative \(\ADM\) Hamiltonian target.
\end{description}

These labels should be read cumulatively rather than competitively.
A final \FullMatch statement may rely on earlier \Exact algebra, on explicitly named
\Controlled reductions, and on one or more declared \Protocol closures.
The point of the labels is to keep that dependence visible instead of hiding it.

\paragraph{Practical reading guide.}
Three conventions are especially useful to keep in mind while reading the paper:
\begin{enumerate}[leftmargin=1.5em]
  \item the generic-frame two-body calculation is written entirely in terms of the
  invariants \(v_A^2\), \(v_B^2\), \(v_{AB}\), \(v_{An}\), and \(v_{Bn}\);
  \item the one-body variable \(u=U/c^2\) is reserved for the strict test-mass reduction;
  and
  \item “exact \(\ADM\) match” always refers to equality of the Legendre-transformed
  Hamiltonian block, not merely to similarity of the Lagrangian coefficient pattern.
\end{enumerate}
These three points account for most of the otherwise avoidable notational ambiguity in the
main derivation.

% ============================================================
\section{Carry-forward formulas from the earlier 4D and \texorpdfstring{$1$}{1}PN papers}
\label{app:carry-forward-formulas}

This appendix collects the specific formulas imported from the earlier \(4\)D and
full conservative \(\OnePN\) papers and used without rederivation in the present
\(\TwoPN\) work.  It is not a replacement for those earlier derivations.  Its job is
narrower: to place in one location the exact projection identities, controlled
near-zone reductions, and frozen \(\OnePN\) coefficient formulas that the main text
assumes as already established.

Only the load-bearing formulas are repeated here.  In particular, we do \emph{not}
reprint the full parent action, the full projected momentum appendix, or the longer
intermediate derivations that led to the wake and breathing-response closures.  The
present paper uses those earlier results only through the formulas recorded below.

\subsection{Exact projection formulas carried from the parent \texorpdfstring{$4$}{4}D paper}
\label{app:carry-forward-4d-projection}

The parent \(4\)D paper defines brane observables by a normalized projection kernel
\(W(w)\):
\begin{equation}
\mathcal P_W[Q](t,\mathbf x)
\equiv
\int_{-\infty}^{\infty} W(w)\,Q(t,\mathbf x,w)\,dw,
\qquad
\int_{-\infty}^{\infty}W(w)\,dw=1.
\label{eq:app-cf-projection}
\end{equation}
This is a measurement prescription, not an approximation.  Projection and reduction are
kept distinct throughout the program.

For the bulk density \(\rho\) and bulk spatial current \(j^i\), the brane-facing density,
current, and velocity are
\begin{equation}
\rho_{\rm br}(t,\mathbf x)\equiv \mathcal P_W[\rho],
\qquad
\mathbf J_{\rm br}(t,\mathbf x)
\equiv
\int_{-\infty}^{\infty}W(w)\,\mathbf j_{xyz}(t,\mathbf x,w)\,dw,
\qquad
\mathbf v_{\rm br}\equiv \frac{\mathbf J_{\rm br}}{\rho_{\rm br}},
\label{eq:app-cf-brane-rho-J-v}
\end{equation}
where \(\mathbf j_{xyz}=(j^x,j^y,j^z)\).  The point to remember is that
\(\mathbf v_{\rm br}\) is a ratio of projected quantities, not the projection of the bulk
velocity.

Starting from exact bulk continuity,
\begin{equation}
\partial_t\rho+\partial_i j^i=0,
\qquad i\in\{x,y,z,w\},
\label{eq:app-cf-bulk-continuity}
\end{equation}
projection gives the exact open-system brane continuity law
\begin{equation}
\partial_t\rho_{\rm br}+\nabla_3\!\cdot\!\mathbf J_{\rm br}=S_{\rm leak},
\label{eq:app-cf-projected-continuity}
\end{equation}
with leakage source
\begin{equation}
S_{\rm leak}(\mathbf x,t)
=
-\left[W(w)j^w(\mathbf x,w,t)\right]_{-\infty}^{+\infty}
+\int_{-\infty}^{\infty}W'(w)\,j^w(\mathbf x,w,t)\,dw.
\label{eq:app-cf-leakage}
\end{equation}
Under fast decay of \(Wj^w\) at infinity this simplifies to
\begin{equation}
S_{\rm leak}(\mathbf x,t)
=
\int_{-\infty}^{\infty}W'(w)\,j^w(\mathbf x,w,t)\,dw.
\label{eq:app-cf-leakage-simplified}
\end{equation}

Next decompose the brane velocity by the standard Helmholtz split,
\begin{equation}
\mathbf v_{\rm br}=\nabla_3\phi_{\rm br}+\mathbf v_T,
\qquad
\nabla_3\!\cdot\!\mathbf v_T=0.
\label{eq:app-cf-helmholtz}
\end{equation}
Then \cref{eq:app-cf-projected-continuity} becomes the exact longitudinal identity
\begin{equation}
\rho_{\rm br}\,\nabla_3^2\phi_{\rm br}
=
S_{\rm leak}
-\partial_t\rho_{\rm br}
-(\nabla_3\rho_{\rm br})\!\cdot\!(\nabla_3\phi_{\rm br}+\mathbf v_T).
\label{eq:app-cf-longitudinal-identity}
\end{equation}
This identity is one of the core exact carry-forward formulas from the parent \(4\)D
paper.  It is not yet a Poisson equation; the Poisson equation appears only after an
explicit near-zone reduction.

\subsection{Controlled Poisson hook and Newtonian point-particle limit}
\label{app:carry-forward-poisson-worldtube}

In the quasi-static near-zone regime, with approximately constant projected density,
subleading advection, and slow leakage/time variation, \cref{eq:app-cf-longitudinal-identity}
reduces to the Poisson hook
\begin{equation}
\nabla_3^2\phi_{\rm br}
\simeq
\frac{S_{\rm eff}}{\rho_0}.
\label{eq:app-cf-poisson-hook}
\end{equation}
For a localized effective source,
\begin{equation}
S_{\rm eff}(\mathbf x)\sim \mathcal S\,\delta^{(3)}(\mathbf x),
\label{eq:app-cf-localized-source}
\end{equation}
this gives the standard inverse-square longitudinal profile
\begin{equation}
\phi_{\rm br}(\mathbf x)
\sim
-\frac{\mathcal S}{4\pi\rho_0}\frac{1}{r},
\qquad
\nabla_3\phi_{\rm br}(\mathbf x)
\sim
\frac{\mathcal S}{4\pi\rho_0}\frac{\hat{\mathbf r}}{r^2}.
\label{eq:app-cf-inverse-square}
\end{equation}
The measured Newtonian potential is then defined by a normalization of this scalar
channel,
\begin{equation}
\Phi_N(\mathbf x,t)=\lambda_\Phi\,\phi_{\rm br}(\mathbf x,t).
\label{eq:app-cf-phiN-def}
\end{equation}
For an isolated compact source \(B\), the near-zone exterior field is carried forward as
\begin{equation}
\Phi_B(\mathbf x,t)
=
-\frac{Gm_B}{|\mathbf x-\mathbf X_B(t)|}
+O\!\left(\frac{Gm_B a_B^2}{|\mathbf x-\mathbf X_B|^3}\right).
\label{eq:app-cf-isolated-source-potential}
\end{equation}

The coherent-defect / small-body reduction then turns a compact worldtube into an
effective particle source.  For a defect with center-of-mass position
\(\mathbf X_{\rm cm}(t)\), characteristic size \(a\), and external field smooth across the
support, the leading equation of motion is Newtonian,
\begin{equation}
M\ddot{\mathbf X}_{\rm cm}
=
-M\nabla_3\Phi_N(\mathbf X_{\rm cm},t)
+O\!\left(Q\,\nabla^3\Phi_N\right),
\label{eq:app-cf-worldtube-force}
\end{equation}
with the first finite-size correction quadrupolar and suppressed by
\begin{equation}
\frac{|\mathbf F_{\rm quad}|}{|\mathbf F_{\rm mono}|}
=O\!\left(\frac{a^2}{r^2}\right).
\label{eq:app-cf-worldtube-quadrupole-suppression}
\end{equation}
The corresponding one-body Newtonian point-particle Lagrangian is
\begin{equation}
L_N
=
\frac12 M\dot{\mathbf X}_{\rm cm}^{\,2}-M\Phi_N(\mathbf X_{\rm cm},t),
\label{eq:app-cf-LN-one-body}
\end{equation}
and pair counting for two compact defects gives the frozen Newtonian two-body block
\begin{equation}
L_0
=
\frac12 m_Av_A^2+\frac12 m_Bv_B^2+\frac{Gm_A m_B}{r_{AB}}.
\label{eq:app-cf-L0-two-body}
\end{equation}
This is the only Newtonian particle mechanics imported by the present \(\TwoPN\) paper.

\subsection{Frozen scalar/optical \texorpdfstring{$1$}{1}PN worldline formulas}
\label{app:carry-forward-scalar-optical}

The \(\OnePN\) paper fixed the self-sector and universal free-particle coefficients by the
reduced scalar/optical worldline ansatz
\begin{equation}
L_{\rm sc}
=
-M_0\left(1+\krho\,\frac{\Phi_N}{c^2}\right)c^2
\sqrt{1-
\frac{v^2/c^2}{1+(n-1)\Phi_N/c^2}}.
\label{eq:app-cf-scalar-optical-ansatz}
\end{equation}
Its weak-field expansion is
\begin{align}
L_{\rm sc}
={}&-M_0c^2-\krho M_0\Phi_N+\frac12M_0v^2+\frac18M_0\frac{v^4}{c^2}
\nonumber\\
&+M_0\left(\frac{\krho}{2}-\frac{n-1}{2}\right)\frac{\Phi_N v^2}{c^2}+O(c^{-4}).
\label{eq:app-cf-scalar-optical-expansion}
\end{align}
Newtonian matching fixes
\begin{equation}
\kappa_\rho=1,
\label{eq:app-cf-krho}
\end{equation}
where \(\krho\) is the historical gravity-side scalar mass-dressing coefficient
rather than an electric charge.
Weak-field optical consistency for the polytropic equation of state
\begin{equation}
P(\rho)=K_{\rm EOS}\rho^n
\label{eq:app-cf-polytrope}
\end{equation}
fixes
\begin{equation}
n=5.
\label{eq:app-cf-n-five}
\end{equation}
With these values, the free-particle part of
\cref{eq:app-cf-scalar-optical-expansion} becomes
\begin{equation}
L_{\rm free}
=
-M_0c^2+\frac12M_0v^2+\frac18M_0\frac{v^4}{c^2}+O(c^{-4}),
\label{eq:app-cf-free-quartic}
\end{equation}
so the frozen \(\OnePN\) Lagrangian coefficient of \(v^4\) is \(1/8\).

The mixed coefficient in \cref{eq:app-cf-scalar-optical-expansion} is
\begin{equation}
\mathcal C_{\rm self}=\frac{\krho}{2}-\frac{n-1}{2}.
\label{eq:app-cf-Cself-def}
\end{equation}
Using \cref{eq:app-cf-krho,eq:app-cf-n-five} gives
\begin{equation}
\mathcal C_{\rm self}=\frac12-2=-\frac32.
\label{eq:app-cf-Cself-value}
\end{equation}
Since \(\Phi_B(\mathbf X_A)=-Gm_B/r_{AB}\), the symmetric two-body self term is therefore
\begin{equation}
L_{\rm self}^{(AB)}
=
\frac{Gm_A m_B}{c^2r_{AB}}\,\frac32\,(v_A^2+v_B^2).
\label{eq:app-cf-self-pair-term}
\end{equation}
This \(+3/2\) coefficient is one of the central frozen inputs to the present \(\TwoPN\)
paper.

\subsection{Frozen static and response ledgers at \texorpdfstring{$1$}{1}PN}
\label{app:carry-forward-static-response}

The static nonlinear \(\OnePN\) term comes from local mass scaling and exact pair
counting.  The effective mass of body \(A\) is written
\begin{equation}
m_A^{\rm eff}
=
m_A\left(1+\kappa_\rho\frac{\Phi_A^{\rm loc}}{c^2}\right)+O(c^{-4}),
\qquad
\Phi_A^{\rm loc}
=
-\sum_{C\neq A}\frac{Gm_C}{r_{AC}},
\label{eq:app-cf-local-mass-scaling}
\end{equation}
with the pair-counted potential
\begin{equation}
L_{\rm pot}
=
-\frac12\sum_A m_A^{\rm eff}\,\Phi_A^{\rm loc}.
\label{eq:app-cf-pair-counted-potential}
\end{equation}
Expanding \cref{eq:app-cf-pair-counted-potential} gives the compact many-body static
formula
\begin{equation}
L_{\rm stat}
=
-\frac{\kappa_\rho G^2}{2c^2}
\sum_A m_A\left(\sum_{C\neq A}\frac{m_C}{r_{AC}}\right)^2.
\label{eq:app-cf-static-many-body}
\end{equation}
For two bodies, with \(\kappa_\rho=1\), this reduces to
\begin{equation}
L_{\rm stat}^{(AB)}
=
-\frac{G^2m_A m_B(m_A+m_B)}{2c^2r_{AB}^2}.
\label{eq:app-cf-static-two-body}
\end{equation}
Thus the frozen static \(\OnePN\) coefficient is exactly \(-1/2\).

The \(\OnePN\) paper also fixed the self-sector ledger in the familiar split
\begin{equation}
\beta_{\OnePN}=\kappa_\rho+\kappa_{\rm add}+\kappa_{\rm PV}.
\label{eq:app-cf-beta-ledger}
\end{equation}
The topology calculation for a \(w\)-uniform translating throat gives
\begin{equation}
\kappa_{\rm add}=\frac12,
\label{eq:app-cf-kappa-add}
\end{equation}
while the counterfactual compact \(4\)D bubble would give
\begin{equation}
\kappa_{\rm add}^{(B^4)}=\frac13.
\label{eq:app-cf-kappa-add-bubble}
\end{equation}
This distinction is carried forward as a geometry/topology discriminator rather than as a
tunable inertial parameter.

Within the declared adiabatic one-degree-of-freedom breathing closure, the response side
is fixed to
\begin{equation}
\kappa_{\rm PV}=\frac32,
\qquad
E_w:E_f:E_{\rm PV}=11:2:5,
\qquad
\frac{d\ln a_*}{d\ln\rho}=-\frac{57}{64}.
\label{eq:app-cf-kappaPV-breathing}
\end{equation}
Therefore the full frozen \(\OnePN\) self ledger is
\begin{equation}
\beta_{\OnePN}
=
\kappa_\rho+\kappa_{\rm add}+\kappa_{\rm PV}
=
1+\frac12+\frac32=3.
\label{eq:app-cf-beta-value}
\end{equation}
The present \(\TwoPN\) paper imports all of
\cref{eq:app-cf-kappa-add,eq:app-cf-kappaPV-breathing,eq:app-cf-beta-value}
without reopening their derivation.

\subsection{Frozen wake formulas and the parity-even cross sector}
\label{app:carry-forward-wake}

The parity-even velocity cross sector is imported from the constructive isotropic wake
calculation of the \(\OnePN\) paper.  For two bodies,
\begin{equation}
L_{\rm vec}^{(AB)}
=
\frac{Gm_A m_B}{c^2r_{AB}}
\left[
C_\parallel\,\mathbf v_A\!\cdot\!\mathbf v_B
+
C_L(\mathbf v_A\!\cdot\!\hat{\mathbf n}_{AB})(\mathbf v_B\!\cdot\!\hat{\mathbf n}_{AB})
\right].
\label{eq:app-cf-cross-sector-template}
\end{equation}
The isotropic Fourier-space wake basis gives
\begin{equation}
C_\parallel=K_{\rm vec}\pi^2(-1+a_H^2-\alpha^2),
\qquad
C_L=K_{\rm vec}\pi^2(-1+a_H^2+\alpha^2),
\label{eq:app-cf-cross-coeff-template}
\end{equation}
where \(\alpha^2\equiv a_L^2\) is the longitudinal wake fraction.

The thermodynamic wake closure imported from the \(\OnePN\) paper is
\begin{equation}
\alpha^2
=
1-\frac{U}{P}
=
1-\frac{1}{n-1}.
\label{eq:app-cf-alpha-thermo}
\end{equation}
Using the already-frozen value \(n=5\) gives
\begin{equation}
\alpha^2=\frac34.
\label{eq:app-cf-alpha-value}
\end{equation}
The parity-even real solution then collapses to
\begin{equation}
a_H=0,
\qquad
K_{\rm vec}=\frac{2}{\pi^2},
\label{eq:app-cf-ah-kvec}
\end{equation}
which yields the exact EIH cross coefficients
\begin{equation}
C_\parallel=-\frac72,
\qquad
C_L=-\frac12.
\label{eq:app-cf-Cpar-CL}
\end{equation}
Thus the frozen two-body velocity cross term carried into the present paper is
\begin{equation}
L_{\rm cross}^{(AB)}
=
\frac{Gm_A m_B}{c^2r_{AB}}
\left[
-\frac72\,\mathbf v_A\!\cdot\!\mathbf v_B
-\frac12(\mathbf v_A\!\cdot\!\hat{\mathbf n}_{AB})(\mathbf v_B\!\cdot\!\hat{\mathbf n}_{AB})
\right].
\label{eq:app-cf-cross-two-body}
\end{equation}
No part of the present \(\TwoPN\) paper is allowed to retune these \(\OnePN\) cross
coefficients.

\subsection{Full carried conservative \texorpdfstring{$1$}{1}PN assembly}
\label{app:carry-forward-full-1pn}

Combining the Newtonian sector, free quartic term, self term, cross term, and static term
recorded above gives the full conservative two-body \(\OnePN\) Lagrangian imported from
the earlier paper:
\begin{align}
L_{\OnePN}^{\rm carried}
={}&
\frac12m_Av_A^2+\frac12m_Bv_B^2
+\frac{m_Av_A^4+m_Bv_B^4}{8c^2}
+\frac{Gm_A m_B}{r_{AB}}
\nonumber\\
&+\frac{Gm_A m_B}{c^2r_{AB}}
\left[
\frac32(v_A^2+v_B^2)
-\frac72\,\mathbf v_A\!\cdot\!\mathbf v_B
-\frac12(\mathbf v_A\!\cdot\!\hat{\mathbf n}_{AB})(\mathbf v_B\!\cdot\!\hat{\mathbf n}_{AB})
\right]
\nonumber\\
&-\frac{G^2m_A m_B(m_A+m_B)}{2c^2r_{AB}^2}.
\label{eq:app-cf-full-1pn-lagrangian}
\end{align}
The exact equality carried forward from the \(\OnePN\) paper is
\begin{equation}
L_{\OnePN}^{\rm carried}=L_{\OnePN}^{\rm EIH}.
\label{eq:app-cf-exact-eih}
\end{equation}
This is one of the main frozen inputs to the present \(\TwoPN\) assembly: the current
paper adds only the order-\(c^{-4}\) block and is not allowed to revise any coefficient
in \cref{eq:app-cf-full-1pn-lagrangian}.

For later reference, the one-body reduction of
\cref{eq:app-cf-full-1pn-lagrangian} is
\begin{equation}
\frac{L_{\rm test}}{m}
=
\frac12v^2+\frac{\mu}{r}
+\frac{1}{c^2}
\left[
\frac18v^4+\frac32\frac{\mu}{r}v^2-\frac12\frac{\mu^2}{r^2}
\right],
\qquad
\mu\equiv GM,
\label{eq:app-cf-test-mass-1pn}
\end{equation}
which can be repackaged in the standard orbit-level form
\begin{equation}
\mathcal L_{\rm orb}
=
\frac12\bigl(1+\sigma(r)\bigr)(\dot r^2+r^2\dot\phi^2)+\frac{\mu}{r},
\qquad
\sigma(r)=\frac{\beta_{\OnePN}\mu}{c^2r}.
\label{eq:app-cf-orbit-form}
\end{equation}
Using \(\beta_{\OnePN}=3\) then gives the carried perihelion result
\begin{equation}
\Delta\phi_{\rm model}
=
\frac{2\pi\beta_{\OnePN}\mu}{c^2a(1-e^2)}
=
\frac{6\pi\mu}{c^2a(1-e^2)}
=\Delta\phi_{\rm GR}.
\label{eq:app-cf-perihelion}
\end{equation}
This is not used directly in the \(\TwoPN\) coefficient solve, but it remains part of the
frozen lower-order status ledger inherited by the present paper.

\claimclass{\Exact{} for the projection formulas
\cref{eq:app-cf-projection,eq:app-cf-projected-continuity,eq:app-cf-longitudinal-identity},
for pair counting in
\cref{eq:app-cf-static-many-body},
for the algebraic assembly
\cref{eq:app-cf-full-1pn-lagrangian},
and for the exact equality
\cref{eq:app-cf-exact-eih} once the imported coefficients are inserted.
\Controlled{} for the Poisson hook, the small-body worldtube reduction, the local
mass-scaling ansatz, the topology result \(\kappa_{\rm add}=1/2\), and the wake closure
leading to
\cref{eq:app-cf-alpha-value,eq:app-cf-ah-kvec,eq:app-cf-Cpar-CL}.
\Protocol{} for the adiabatic breathing-response input
\cref{eq:app-cf-kappaPV-breathing} and therefore for the carried value
\(\beta_{\OnePN}=3\).
\FullMatch{} for the exact EIH equality
\cref{eq:app-cf-exact-eih} and the carried perihelion formula
\cref{eq:app-cf-perihelion}.}

% ============================================================
\section{One-body \texorpdfstring{$\DtN$}{DtN} closure details}
\label{app:one-body-dtn}

This appendix records the detailed algebra behind the one-body \(\TwoPN\)
closure used in \cref{sec:one-body}.  The main text gave only the minimal proof
line: the low-frequency monopole mouth operator supplies two normalized invariants,
the exact isotropic Schwarzschild test-mass target fixes the required effective
velocity denominator, and preservation of the frozen \(\OnePN\) slot then leaves a
unique minimal quadratic geometry correction.

The purpose of the present appendix is to unpack that chain in a more constructive
way.  We first write the simplest unit-test mouth operator explicitly, then show why
its normalized invariant combinations isolate sound-speed and geometry information,
then derive the target denominator directly from the isotropic Schwarzschild
one-body expansion, and finally show how the earlier raw resonance-proxy fit
factorizes into a clean \emph{raw DtN piece} times a \emph{port correction}.
Nothing in this appendix is needed to state the main theorem, but it makes the
one-body closure much less opaque.

\subsection{Low-frequency monopole mouth operator on the cylinder unit-test branch}
\label{app:one-body-dtn-cylinder}

The simplest conservative mouth model retained in the \(\TwoPN\) notes is the
monopole branch of a cylinder-like throat with one axial coordinate \(z\in[0,L]\),
a mouth at \(z=0\), and a Neumann bottom at \(z=L\).  This is a unit-test branch,
not the final moving-throat PDE.  Its value is that the low-frequency Dirichlet-to-
Neumann map can be written in closed form and already exposes the two geometric
quantities that matter for the one-body closure: the effective sound speed \(c_s\)
and the throat depth \(L\).

Let \(\psi(z,\omega)\) satisfy the one-dimensional Helmholtz equation
\begin{equation}
\psi''+k^2\psi=0,
\qquad
k\equiv \frac{\omega}{c_s},
\label{eq:app-dtn-helmholtz}
\end{equation}
with bottom condition
\begin{equation}
\psi'(L,\omega)=0
\label{eq:app-dtn-neumann-bottom}
\end{equation}
and mouth datum
\begin{equation}
\psi(0,\omega)=\psi_m(\omega).
\label{eq:app-dtn-mouth-datum}
\end{equation}
The unique regular solution is
\begin{equation}
\psi(z,\omega)=
\psi_m(\omega)
\frac{\cos\!\bigl(k(L-z)\bigr)}{\cos(kL)}.
\label{eq:app-dtn-cylinder-solution}
\end{equation}
Taking the outward normal derivative at the mouth gives
\begin{equation}
-\partial_z\psi(0,\omega)
=
-k\tan(kL)\,\psi_m(\omega).
\label{eq:app-dtn-mouth-derivative}
\end{equation}
Thus the monopole mouth operator is
\begin{equation}
Z_{00}(\omega;u)=-k\tan(kL)
=-\frac{\omega}{c_s}\tan\!\left(\frac{\omega L}{c_s}\right).
\label{eq:app-dtn-cylinder-exact}
\end{equation}
Expanding at low frequency yields
\begin{equation}
Z_{00}(\omega;u)=Z_2(u)\,\omega^2+Z_4(u)\,\omega^4+O(\omega^6),
\label{eq:app-dtn-Z-expansion}
\end{equation}
with
\begin{equation}
Z_2(u)=-\frac{L(u)}{c_s^2(u)},
\qquad
Z_4(u)=-\frac{L^3(u)}{3c_s^4(u)}.
\label{eq:app-dtn-Z2Z4}
\end{equation}
The absence of an \(\omega^0\) term is automatic in this conservative unit-test
branch: once the static monopole potential is accounted for separately, the retained
mouth operator begins at even time-derivative order.

The main text uses the normalized invariant combinations
\begin{equation}
C_s(u)
\equiv
\frac{Z_4/Z_2^3}{(Z_4/Z_2^3)_0},
\qquad
G(u)
\equiv
\frac{Z_4/Z_2^2}{(Z_4/Z_2^2)_0},
\label{eq:app-dtn-invariants}
\end{equation}
where the subscript \(0\) means evaluation at \(u=0\).  In the cylinder branch,
these become exactly
\begin{equation}
C_s(u)=\frac{c_s^2(u)}{c_{s0}^2},
\qquad
G(u)=\frac{L(u)}{L_0}.
\label{eq:app-dtn-cylinder-invariants}
\end{equation}
This is why the choice \cref{eq:app-dtn-invariants} is so useful: all irrelevant
normalizations and numerical factors cancel, leaving one invariant that tracks the
sound-speed sector and one invariant that tracks the geometry / depth sector.

Equivalently, the low-frequency coefficients themselves can be reconstructed from the
invariants as
\begin{equation}
\frac{Z_2(u)}{Z_2(0)}=\frac{G(u)}{C_s(u)},
\qquad
\frac{Z_4(u)}{Z_4(0)}=\frac{G^3(u)}{C_s^2(u)}.
\label{eq:app-dtn-Z-from-invariants}
\end{equation}
This inversion is often the most convenient route for checking explicit series.

\subsection{Current conservative freeze and explicit series data}
\label{app:one-body-dtn-current-series}

The one-body \(\TwoPN\) closure is not performed in a vacuum; it inherits specific
lower-order data from the frozen \(\OnePN\) program.

First, the exact Bernoulli continuation of the already-fixed \(n=5\) barotropic
sector gives the sound-speed factor
\begin{equation}
\frac{c_s^2(u)}{c^2}=1-4u,
\qquad
u\equiv \frac{U}{c^2},
\qquad
U=\frac{GM}{r}.
\label{eq:app-dtn-cs-bernoulli}
\end{equation}
In the present normalized language this is simply
\begin{equation}
C_s(u)=1-4u.
\label{eq:app-dtn-Cs-current}
\end{equation}
Second, the carried low-frequency throat closure writes the geometric depth factor as
\begin{equation}
G(u)=\frac{L(u)}{L_0}=a(u),
\label{eq:app-dtn-G-a}
\end{equation}
with the already-fixed \(\OnePN\) breathing slope
\begin{equation}
g_1\equiv \left.\frac{dG}{du}\right|_{u=0}=\frac{57}{64}.
\label{eq:app-dtn-g1}
\end{equation}
At the current closure level, the explicit series used in the notes is
\begin{equation}
G(u)=1+\frac{57}{64}u+\frac{298821}{131072}u^2+O(u^3).
\label{eq:app-dtn-G-current-series}
\end{equation}
Substituting \cref{eq:app-dtn-Cs-current,eq:app-dtn-G-current-series} into
\cref{eq:app-dtn-Z-from-invariants} yields
\begin{equation}
\frac{Z_2(u)}{Z_2(0)}
=1+\frac{313}{64}u+\frac{2862917}{131072}u^2+O(u^3),
\label{eq:app-dtn-Z2-series}
\end{equation}
\begin{equation}
\frac{Z_4(u)}{Z_4(0)}
=1+\frac{683}{64}u+\frac{10301487}{131072}u^2+O(u^3).
\label{eq:app-dtn-Z4-series}
\end{equation}
Conversely, inserting these series back into \cref{eq:app-dtn-invariants} reproduces
\begin{equation}
C_s(u)=1-4u,
\qquad
G(u)=1+\frac{57}{64}u+\frac{298821}{131072}u^2+O(u^3),
\label{eq:app-dtn-invariants-current-series}
\end{equation}
so the normalized invariant extraction is self-consistent.

Two features of \cref{eq:app-dtn-invariants-current-series} matter for the rest of the
appendix.  First, the sound-speed invariant is already completely fixed through the
current order.  Second, only the \emph{linear} geometry slope \(g_1=57/64\) will enter
into the actual \(\TwoPN\) denominator match; the coefficient of \(u^2\) in
\(G(u)\) does not affect the conservative one-body target through \(c^{-4}\).

\subsection{The exact isotropic Schwarzschild target as a denominator problem}
\label{app:one-body-dtn-target}

The exact isotropic Schwarzschild test-mass Lagrangian is
\begin{equation}
\frac{L_{\mathrm{Schw}}}{m}
=-c^2\sqrt{\left(\frac{1-u/2}{1+u/2}\right)^2-(1+u/2)^4\frac{v^2}{c^2}}.
\label{eq:app-dtn-schwarzschild-exact}
\end{equation}
Expanding through total post-Newtonian order \(c^{-4}\) gives
\begin{align}
\frac{L_{\mathrm{Schw}}}{m}
={}&-c^2
+\frac12v^2+U
+\frac{1}{c^2}\left(\frac18v^4+\frac32Uv^2-\frac12U^2\right)
\nonumber\\
&
+\frac{1}{c^4}\left(\frac1{16}v^6+\frac78Uv^4+2U^2v^2+\frac14U^3\right)
+O(c^{-6}).
\label{eq:app-dtn-schwarzschild-2pn}
\end{align}
In the carried toy-model assembly, the \(v^6\) and \(Uv^4\) terms are already produced
by the minimal Bernoulli exactification of the self sector, while the pure-static
\(U^3\) term is handled separately by quadratic local mass scaling.  So the genuinely
missing one-body response datum is the coefficient of \(U^2v^2\), which must be
\begin{equation}
2.
\label{eq:app-dtn-target-u2v2}
\end{equation}

To isolate that datum, write the reduced one-body worldline in the denominator form
\begin{equation}
L_{\mathrm{red}}
=-mc^2(1-u)\sqrt{1-\frac{v^2/c^2}{D(u)}},
\qquad
D(u)=1-4u+d_2u^2+O(u^3).
\label{eq:app-dtn-reduced-worldline}
\end{equation}
Expanding \cref{eq:app-dtn-reduced-worldline} gives
\begin{align}
\frac{L_{\mathrm{red}}}{m}
={}&-c^2+\frac12v^2+U
+\frac{1}{c^2}\left(\frac18v^4+\frac32Uv^2\right)
\nonumber\\
&
+\frac{1}{c^4}\left(
\frac1{16}v^6+\frac78Uv^4+
\left(6-\frac{d_2}{2}\right)U^2v^2
\right)
+O(c^{-6}).
\label{eq:app-dtn-reduced-worldline-expanded}
\end{align}
Comparing \cref{eq:app-dtn-reduced-worldline-expanded} with the exact target
\cref{eq:app-dtn-schwarzschild-2pn} shows that the effective denominator must satisfy
\begin{equation}
6-\frac{d_2}{2}=2,
\qquad\text{hence}\qquad
d_2=8.
\label{eq:app-dtn-d2-solve}
\end{equation}
So the exact isotropic one-body target can be restated as the denominator condition
\begin{equation}
D_{\mathrm{target}}(u)=1-4u+8u^2+O(u^3).
\label{eq:app-dtn-denominator-target}
\end{equation}
This repackaging is the reason the low-frequency mouth invariants are useful at all: the
entire one-body closure problem reduces to producing the correct quadratic coefficient in
\cref{eq:app-dtn-denominator-target} without changing the already-frozen linear term.

\subsection{Minimal invariant ansatz and uniqueness of the quadratic geometry correction}
\label{app:one-body-dtn-uniqueness}

Within the current conservative hierarchy, the sound-speed invariant \(C_s(u)\) is
already frozen by the Bernoulli continuation, so any genuinely new one-body response
information must enter through the geometry invariant \(G(u)\).  The smallest analytic
invariant ansatz is therefore
\begin{equation}
D_{\mathrm{eff}}(u)=C_s(u)\,F\bigl(G(u)\bigr),
\qquad
F(1)=1,
\label{eq:app-dtn-general-F}
\end{equation}
with \(F\) expanded near \(G=1\).  Through \(\TwoPN\), this is equivalent to
\begin{equation}
D_{\mathrm{eff}}(u)=C_s(u)\left[1+\alpha\,(G-1)+\beta\,(G-1)^2\right]+O(u^3).
\label{eq:app-dtn-general-ansatz}
\end{equation}
Write
\begin{equation}
G(u)=1+g_1u+g_2u^2+O(u^3),
\qquad
C_s(u)=1-4u+O(u^2).
\label{eq:app-dtn-G-generic-series}
\end{equation}
Then
\begin{equation}
D_{\mathrm{eff}}(u)=1+(-4+\alpha g_1)u+\bigl(\beta g_1^2+\cdots\bigr)u^2+O(u^3).
\label{eq:app-dtn-ansatz-expanded}
\end{equation}
Because the \(u\)-coefficient is already frozen by the exact \(\OnePN\) ledger,
preservation of the lower-order slot requires
\begin{equation}
\alpha g_1=0.
\label{eq:app-dtn-alpha-condition}
\end{equation}
Since the carried breathing slope is nonzero,
\begin{equation}
g_1=\frac{57}{64}\neq 0,
\label{eq:app-dtn-g1-nonzero}
\end{equation}
we must have
\begin{equation}
\alpha=0.
\label{eq:app-dtn-alpha-zero}
\end{equation}
So the first admissible correction is quadratic:
\begin{equation}
D_{\mathrm{eff}}(u)=C_s(u)\left[1+\mu\,(G(u)-1)^2\right].
\label{eq:app-dtn-quadratic-ansatz}
\end{equation}
Expanding \cref{eq:app-dtn-quadratic-ansatz} gives
\begin{equation}
D_{\mathrm{eff}}(u)=1-4u+\mu g_1^2u^2+O(u^3).
\label{eq:app-dtn-quadratic-expanded}
\end{equation}
Matching the exact denominator target \cref{eq:app-dtn-denominator-target} therefore
forces
\begin{equation}
\mu g_1^2=8.
\label{eq:app-dtn-mu-equation}
\end{equation}
Using \cref{eq:app-dtn-g1} gives
\begin{equation}
\mu=\frac{8}{g_1^2}=\frac{32768}{3249}\approx 10.0855647892.
\label{eq:app-dtn-mu-value}
\end{equation}
Hence the final one-body closure used in the main text is
\begin{equation}
D_{\mathrm{eff}}(u)
=C_s(u)\left[1+\frac{32768}{3249}\,(G(u)-1)^2\right],
\label{eq:app-dtn-final-closure}
\end{equation}
which indeed expands to
\begin{equation}
D_{\mathrm{eff}}(u)=1-4u+8u^2+O(u^3).
\label{eq:app-dtn-final-series}
\end{equation}

Three remarks are worth making explicit.
First, the uniqueness here is a \emph{through-\(\TwoPN\)} uniqueness within the stated
analytic invariant hierarchy.  Nothing stronger is claimed.
Second, only the linear breathing slope \(g_1\) matters for the coefficient match;
the higher geometric series coefficients of \(G\) do not enter until higher PN order.
Third, the closure is automatically compatible with the frozen \(\OnePN\) ledger because
that compatibility is exactly what forced \(\alpha=0\).

\begin{claim}[One-body invariant closure through \texorpdfstring{$2$}{2}PN]
\label{clm:app-dtn-one-body-closure}
Within the conservative analytic ansatz \cref{eq:app-dtn-general-ansatz}, the frozen
\(\OnePN\) denominator coefficient forces the geometry correction to begin quadratically
in \(G-1\), and the exact isotropic Schwarzschild test-mass target then fixes the
quadratic coefficient uniquely to
\begin{equation}
\mu=\frac{32768}{3249}.
\label{eq:app-dtn-claim-mu}
\end{equation}
Equivalently, the completed one-body denominator is exactly
\begin{equation}
D_{\mathrm{eff}}(u)=1-4u+8u^2+O(u^3)
\label{eq:app-dtn-claim-target}
\end{equation}
through conservative \(\TwoPN\) order.
\end{claim}

\subsection{Relation to the earlier raw resonance-proxy fit}
\label{app:one-body-dtn-raw-proxy}

Before the invariant packaging \cref{eq:app-dtn-final-closure} was written in its final
form, the same low-frequency DtN data had been summarized by a more primitive “raw
resonance proxy,” namely
\begin{equation}
D_{\mathrm{raw}}(u)
\equiv
\frac{[Z_2/Z_4](u)}{[Z_2/Z_4](0)}.
\label{eq:app-dtn-draw-def}
\end{equation}
Using \cref{eq:app-dtn-Z2Z4,eq:app-dtn-cylinder-invariants}, this can be written exactly
as
\begin{equation}
D_{\mathrm{raw}}(u)=\frac{C_s(u)}{G(u)^2}.
\label{eq:app-dtn-draw-csg}
\end{equation}
With the current series data,
\begin{equation}
D_{\mathrm{raw}}(u)
=1-\frac{185}{32}u+\frac{324075}{65536}u^2+O(u^3).
\label{eq:app-dtn-draw-series}
\end{equation}
The problem with \cref{eq:app-dtn-draw-series} is immediate: it does \emph{not}
preserve the already-frozen \(\OnePN\) coefficient, because the linear term is
\(-185/32\) instead of \(-4\).

The final closure shows exactly how that raw proxy must be corrected.  Define the
multiplicative port factor
\begin{equation}
P_{\mathrm{port}}(u)
\equiv
G(u)^2\left[1+\mu\,(G(u)-1)^2\right].
\label{eq:app-dtn-port-factor-def}
\end{equation}
Then
\begin{equation}
D_{\mathrm{eff}}(u)=D_{\mathrm{raw}}(u)\,P_{\mathrm{port}}(u).
\label{eq:app-dtn-factorization}
\end{equation}
With \(\mu=32768/3249\), the current series is
\begin{equation}
P_{\mathrm{port}}(u)
=1+\frac{57}{32}u+\frac{875093}{65536}u^2+O(u^3).
\label{eq:app-dtn-port-factor-series}
\end{equation}
Multiplying \cref{eq:app-dtn-draw-series,eq:app-dtn-port-factor-series} gives back the
exact target series
\begin{equation}
D_{\mathrm{eff}}(u)=1-4u+8u^2+O(u^3).
\label{eq:app-dtn-factorized-target}
\end{equation}
So the earlier port fit was not an arbitrary extra correction.  It can be read more
cleanly as a factorization statement: the ratio \(Z_2/Z_4\) supplies a raw DtN proxy,
while the missing geometry-port dressing is exactly the factor required to preserve the
frozen \(\OnePN\) slot and hit the \(\TwoPN\) one-body target.

This factorization is also useful conceptually.  It shows that the one-body closure is
not inventing a new denominator from scratch; it is reorganizing the same low-frequency
mouth data into a product of a raw response ratio and a compensating geometry-port
factor.  The raw piece by itself is too naive, but it already contains the right
low-frequency ingredients.

\subsection{Scope of the one-body closure}
\label{app:one-body-dtn-scope}

The one-body closure derived above is deliberately narrow.
It closes the missing conservative test-mass \(\TwoPN\) response slot and nothing more.
Several things are fixed, and several things remain open.

What is fixed is the following.
The sound-speed invariant is carried as \(C_s(u)=1-4u\).
The geometry slope is carried as \(g_1=57/64\).
Within the declared conservative analytic invariant hierarchy, these data uniquely fix
\(\mu=32768/3249\) and therefore fix the denominator target
\(D_{\mathrm{eff}}(u)=1-4u+8u^2+O(u^3)\).
As a result, the one-body \(U^2v^2\) slot is no longer open in the main \(\TwoPN\)
assembly.

What remains open is the deeper derivation of this closure from the fully dynamical
inner-throat PDE and its frequency-dependent \(\DtN\) operator.
In particular, the present appendix does not derive the analytic structure of
\(F(G)\) beyond the minimal quadratic term needed at \(\TwoPN\), nor does it determine
the dynamic pole scales of the appendix-side mouth-response operator.  Those are genuine
inner-throat observables and belong to the later PDE / \(\DtN\) program rather than to
the conservative algebraic proof completed in the main text.

\claimclass{\Controlled{} for the cylinder / Neumann-bottom unit-test branch and the
import of the Bernoulli and breathing series data.
\Protocol{} for the analytic invariant ansatz
\(D_{\mathrm{eff}}=C_s[1+\alpha(G-1)+\beta(G-1)^2]\).
\Exact{} for the algebraic elimination of \(\alpha\), for the determination of
\(\mu=32768/3249\), and for the factorization
\(D_{\mathrm{eff}}=D_{\mathrm{raw}}P_{\mathrm{port}}\) once the closure hierarchy is declared.
\FullMatch{} for the statement that the resulting denominator reproduces the exact
isotropic Schwarzschild one-body target through conservative \(\TwoPN\) order.}


% ============================================================
\section{Bernoulli exactification, quadratic local mass scaling, and the static hierarchy}
\label{app:static-scaling}

This appendix collects the detailed self/static formulas used in
\cref{sec:self-static}.  The main text kept only the proof line: exact
\(n=5\) Bernoulli continuation fixes the sound-speed sector, minimal optical
exactification supplies the raw self contribution through \(\TwoPN\), quadratic local
mass scaling fixes the static hierarchy, and the pure-static isotropic test-mass gate
selects
\(\lambda_\rho=1/2\).
The purpose of the present appendix is to write those statements in a form that can be
checked term by term.

There are four points to keep separate.
First, the barotropic \(n=5\) Bernoulli continuation is exact within the reduced matter
model and already determines the optical index through \(\TwoPN\).
Second, only the optical denominator is Bernoulli-exactified in the self sector; a naïve
full exactification of the mass prefactor would spoil the already-frozen \(\OnePN\)
static ledger.
Third, quadratic local mass scaling generates a compact many-body cubic-static hierarchy,
of which the two-body term is selected by the isotropic one-body target and the higher-body
terms are predictions of the same bookkeeping rule.
Fourth, any interpretation of these terms as universal \(1/r^4\)-order structure must
respect the worldtube finite-size gate derived at the end of the appendix.

\subsection{Exact \texorpdfstring{$n=5$}{n=5} continuation of the barotropic sector}
\label{app:static-bernoulli}

The barotropic sector carried from the earlier papers is
\begin{equation}
P(\rho)=K_{\rm EOS}\rho^5,
\qquad
U_{\rm int}(\rho)=\frac{K_{\rm EOS}}{4}\rho^5,
\label{eq:app-static-eos}
\end{equation}
so the enthalpy variable is
\begin{equation}
h(\rho)\equiv \frac{dU_{\rm int}}{d\rho}=\frac{5K_{\rm EOS}}{4}\rho^4,
\label{eq:app-static-enthalpy}
\end{equation}
and the sound speed satisfies
\begin{equation}
c_s^2(\rho)=\frac{1}{m_{\rm mat}}\frac{dP}{d\rho}
=\frac{5K_{\rm EOS}}{m_{\rm mat}}\rho^4.
\label{eq:app-static-cs2-def}
\end{equation}
In the reduced matter units used in the harness, one sets \(m_{\rm mat}=1\) and fixes the
background normalization by demanding
\begin{equation}
c_s^2(\rho_0)=c^2.
\label{eq:app-static-background-cs}
\end{equation}
This gives
\begin{equation}
K_{\rm EOS}=\frac{c^2}{5\rho_0^4},
\qquad
h_0\equiv h(\rho_0)=\frac{c^2}{4}.
\label{eq:app-static-KEOS-h0}
\end{equation}

The Bernoulli continuation is then exact:
\begin{equation}
h(\rho)=h_0+\Phi,
\label{eq:app-static-bernoulli-eq}
\end{equation}
with \(\Phi\) the local Newtonian potential.
Using \cref{eq:app-static-enthalpy,eq:app-static-KEOS-h0}, one solves
\cref{eq:app-static-bernoulli-eq} directly for the density profile,
\begin{equation}
\rho_{\rm Bern}(\Phi)=\rho_0\left(1+\frac{4\Phi}{c^2}\right)^{1/4}.
\label{eq:app-static-rho-bernoulli}
\end{equation}
Substituting \cref{eq:app-static-rho-bernoulli} into \cref{eq:app-static-cs2-def} gives
\begin{equation}
\frac{c_s^2(\Phi)}{c^2}=1+\frac{4\Phi}{c^2}.
\label{eq:app-static-cs2-phi}
\end{equation}
Equivalently, in the attractive one-body notation \(\Phi=-U\) with
\(u\equiv U/c^2\),
\begin{equation}
\frac{c_s^2}{c^2}=1-4u.
\label{eq:app-static-cs2-u}
\end{equation}
This is the exact Bernoulli input used throughout the \(\TwoPN\) one-body and self-sector
analysis.

The corresponding optical index is
\begin{equation}
N(\Phi)\equiv \frac{c}{c_s(\Phi)}=
\left(1+\frac{4\Phi}{c^2}\right)^{-1/2},
\label{eq:app-static-optical-index}
\end{equation}
so in \(U=-\Phi\) variables,
\begin{equation}
N(U)=\left(1-\frac{4U}{c^2}\right)^{-1/2}
=1+2\frac{U}{c^2}+6\frac{U^2}{c^4}+O(c^{-6}).
\label{eq:app-static-optical-index-series}
\end{equation}
The coefficients in \cref{eq:app-static-optical-index-series} are exactly the ones checked
by the compact \(\TwoPN\) harness and are the source of the carried
\(\frac{7}{8}Uv^4\) and raw \(6U^2v^2\) terms in the self sector.

With this optical continuation in hand, the minimal self exactification used in the main
text is
\begin{equation}
L_{\rm self,min}
=-mc^2\left(1+\frac{\Phi}{c^2}\right)
\sqrt{1-\frac{v^2}{c^2\left(1+4\Phi/c^2\right)}}.
\label{eq:app-static-Lself-min}
\end{equation}
Expanding \cref{eq:app-static-Lself-min} through total \(\TwoPN\) order gives
\begin{align}
\frac{L_{\rm self,min}}{m}
={}&-c^2+\frac12v^2-\Phi
+\frac{1}{c^2}\left(\frac18v^4-\frac32\Phi v^2\right)
\nonumber\\
&+\frac{1}{c^4}\left(\frac1{16}v^6-\frac78\Phi v^4+6\Phi^2v^2\right)
+O(c^{-6}).
\label{eq:app-static-Lself-min-phi}
\end{align}
Writing \(\Phi=-U\) converts this to the more familiar attractive-field form
\begin{align}
\frac{L_{\rm self,min}}{m}
={}&-c^2+\frac12v^2+U
+\frac{1}{c^2}\left(\frac18v^4+\frac32Uv^2\right)
\nonumber\\
&+\frac{1}{c^4}\left(\frac1{16}v^6+\frac78Uv^4+6U^2v^2\right)
+O(c^{-6}).
\label{eq:app-static-Lself-min-U}
\end{align}
Thus the exact \(n=5\) Bernoulli continuation plus minimal optical exactification already
fixes the raw self-sector coefficients
\begin{equation}
\frac1{16},\qquad \frac78,\qquad 6.
\label{eq:app-static-self-coeffs}
\end{equation}

The important obstruction appears if one tries to Bernoulli-exactify the mass prefactor
as well.  The naïve aggressive continuation would be
\begin{equation}
L_{\rm self,aggr}
=-mc^2\left(1+\frac{4\Phi}{c^2}\right)^{1/4}
\sqrt{1-\frac{v^2}{c^2\left(1+4\Phi/c^2\right)}}.
\label{eq:app-static-Lself-aggr}
\end{equation}
Its \(\OnePN\) expansion is
\begin{equation}
\frac{L_{\rm self,aggr}}{m}
=-c^2+\frac12v^2-\Phi
+\frac{1}{c^2}\left(\frac18v^4-\frac32\Phi v^2+\frac32\Phi^2\right)
+O(c^{-4}).
\label{eq:app-static-Lself-aggr-1pn}
\end{equation}
Since \(\Phi^2=U^2\), this would generate a
\(+\frac32U^2/c^2\) static term at \(\OnePN\), which is incompatible with the already-frozen
pair-counted static ledger
\(-\frac12U^2/c^2\).
That is why the present paper exactifies only the optical denominator and leaves the mass
prefactor at its frozen linear form.

\claimclass{\Exact{} for the Bernoulli continuation
\cref{eq:app-static-rho-bernoulli,eq:app-static-cs2-phi,eq:app-static-optical-index}.
\Controlled{} for the minimal self-sector continuation
\cref{eq:app-static-Lself-min}, whose role is to preserve the frozen \(\OnePN\) self
ledger while generating the raw universal \(\TwoPN\) self coefficients.
\Protocol{} for the decision not to Bernoulli-exactify the mass prefactor beyond linear
order because doing so would violate the already-frozen \(\OnePN\) static regression.}

\subsection{Two-body static sector and \texorpdfstring{$\lambda_\rho=1/2$}{lambda\_rho=1/2}}
\label{app:static-two-body}

The static sector is extended independently of the optical/self continuation by
quadratic local mass scaling.  For each body,
\begin{equation}
m_A^{\rm eff}
=m_A\left(1+\frac{\Phi_A^{\rm loc}}{c^2}
+\lambda_\rho\frac{(\Phi_A^{\rm loc})^2}{c^4}\right),
\qquad
m_B^{\rm eff}
=m_B\left(1+\frac{\Phi_B^{\rm loc}}{c^2}
+\lambda_\rho\frac{(\Phi_B^{\rm loc})^2}{c^4}\right),
\label{eq:app-static-meff}
\end{equation}
where
\begin{equation}
\Phi_A^{\rm loc}=-\frac{Gm_B}{r},
\qquad
\Phi_B^{\rm loc}=-\frac{Gm_A}{r}
\label{eq:app-static-local-potentials-two-body}
\end{equation}
for the two-body system.
The pair-counted potential is
\begin{equation}
L_{\rm pot}
=-\frac12\left(m_A^{\rm eff}\Phi_A^{\rm loc}+m_B^{\rm eff}\Phi_B^{\rm loc}\right).
\label{eq:app-static-Lpot-pair}
\end{equation}
Expanding \cref{eq:app-static-Lpot-pair} gives the complete two-body static sector through
\(\TwoPN\):
\begin{equation}
L_{\rm static}^{(AB)}
=\frac{Gm_A m_B}{r}
-\frac{G^2m_A m_B(m_A+m_B)}{2c^2r^2}
+\frac{\lambda_\rho G^3m_A m_B(m_A^2+m_B^2)}{2c^4r^3}
+O(c^{-6}).
\label{eq:app-static-two-body-static}
\end{equation}
Thus the quadratic local-mass rule fixes the two-body \(\TwoPN\) static coefficient to
\begin{equation}
\frac{\lambda_\rho}{2}.
\label{eq:app-static-two-body-lambda-half}
\end{equation}

In the strict test-mass reduction, \cref{eq:app-static-two-body-static} becomes
\begin{equation}
\frac{L_{\rm static,1body}}{m}
=U-\frac12\frac{U^2}{c^2}+\frac{\lambda_\rho}{2}\frac{U^3}{c^4}+O(c^{-6}).
\label{eq:app-static-one-body-static}
\end{equation}
The exact isotropic Schwarzschild target requires the pure-static coefficient of
\(U^3/c^4\) to be \(+\frac14\), so the isotropic gate immediately fixes
\begin{equation}
\lambda_\rho=\frac12.
\label{eq:app-static-lambda-rho-half}
\end{equation}
Once \cref{eq:app-static-lambda-rho-half} is imposed, the two-body static block becomes
\begin{equation}
L_{2,\rm static}^{(AB)}
=\frac{G^3m_A m_B(m_A^2+m_B^2)}{4r^3},
\label{eq:app-static-two-body-static-fixed}
\end{equation}
where, as throughout the paper, \(L_2\) denotes the coefficient of \(c^{-4}\) in the
post-Newtonian expansion.

Combining the static contribution \cref{eq:app-static-one-body-static} with the minimal
self exactification \cref{eq:app-static-Lself-min-U} gives the raw one-body self/static
candidate
\begin{align}
\frac{L_{\rm self+static}^{\rm raw}}{m}
={}&-c^2+\frac12v^2+U
+\frac{1}{c^2}\left(\frac18v^4+\frac32Uv^2-\frac12U^2\right)
\nonumber\\
&+\frac{1}{c^4}\left(\frac1{16}v^6+\frac78Uv^4+6U^2v^2+\frac{\lambda_\rho}{2}U^3\right)
+O(c^{-6}).
\label{eq:app-static-self-static-raw}
\end{align}
Fixing \(\lambda_\rho=1/2\) matches the pure-static isotropic coefficient exactly and leaves
\begin{equation}
\Delta L_{\rm raw\to iso}
=\frac{4U^2v^2}{c^4}
\label{eq:app-static-residual-4u2v2}
\end{equation}
as the sole remaining one-body discrepancy.  This is precisely the missing test-mass slot
closed by the \(\DtN\)-invariant denominator repair in \cref{sec:one-body} and
\cref{app:one-body-dtn}.

\claimclass{\Exact{} for the pair-counted two-body static expansion
\cref{eq:app-static-two-body-static}. \Controlled{} for the quadratic local mass-scaling
rule. \FullMatch{} for the pure-static isotropic selection
\cref{eq:app-static-lambda-rho-half} and for the statement that, after this selection,
the unrepaired one-body discrepancy sits entirely in the \(U^2v^2\) sector.}

\subsection{Selected three-body and four-body static terms}
\label{app:static-higher-body}

The same quadratic local mass-scaling rule generates a compact many-body static hierarchy.
For an arbitrary collection of bodies with local potentials
\begin{equation}
\Phi_A^{\rm loc}=-\sum_{C\neq A}\frac{Gm_C}{r_{AC}},
\label{eq:app-static-local-potential-general}
\end{equation}
the \(\TwoPN\) static sector is simply
\begin{equation}
L_{2,\rm static}^{(N)}
=-\frac{\lambda_\rho}{2}\sum_A m_A\bigl(\Phi_A^{\rm loc}\bigr)^3.
\label{eq:app-static-general-many-body}
\end{equation}
Equation~\eqref{eq:app-static-general-many-body} is the entire higher-body static
hierarchy in compact form.  The point of the present subsection is only to record the
first representative monomials that were checked explicitly in the harness.

For three bodies \(A,B,C\), one has
\begin{equation}
\Phi_A^{\rm loc}=-G\left(\frac{m_B}{r_{AB}}+\frac{m_C}{r_{AC}}\right),
\qquad
\Phi_B^{\rm loc}=-G\left(\frac{m_A}{r_{AB}}+\frac{m_C}{r_{BC}}\right),
\qquad
\Phi_C^{\rm loc}=-G\left(\frac{m_A}{r_{AC}}+\frac{m_B}{r_{BC}}\right).
\label{eq:app-static-three-body-phi}
\end{equation}
Expanding \cref{eq:app-static-general-many-body} gives selected cubic-static terms of the
form
\begin{equation}
L_{2,\rm static}^{(ABC)}\supset
\frac{3\lambda_\rho G^3}{2}
\left[
\frac{m_A m_B^2 m_C}{r_{AB}^2 r_{AC}}
+\frac{m_A m_B m_C^2}{r_{AB}r_{AC}^2}
+\text{cyclic relabelings}
\right],
\label{eq:app-static-three-body-selected}
\end{equation}
where the overall \(c^{-4}\) factor is understood as part of \(L_2\).
The corresponding unit-normalized coefficients are exactly unity, in agreement with the
compact harness checks.

Just as important is what does \emph{not} appear.  The coefficient of the pure triangle
triple product
\begin{equation}
\frac{1}{r_{AB}r_{AC}r_{BC}}
\label{eq:app-static-pure-triangle}
\end{equation}
vanishes identically in this hierarchy.  Equivalently,
\begin{equation}
\Bigl[\frac{1}{r_{AB}r_{AC}r_{BC}}\Bigr]L_{2,\rm static}^{(ABC)}=0.
\label{eq:app-static-pure-triangle-zero}
\end{equation}
This is not accidental.  Every term in
\cref{eq:app-static-general-many-body} is anchored on one body, so each cubic monomial is
a product of three pair potentials that all emanate from the same anchor body.  A pure
triangle product involving three different edges and no common anchor is therefore absent.

The first genuine four-body monomial appears when one body couples simultaneously to three
others.  For four bodies \(A,B,C,D\), the harness explicitly checks the coefficient of
\begin{equation}
\frac{1}{r_{AB}r_{AC}r_{AD}},
\label{eq:app-static-four-body-monomial}
\end{equation}
which is
\begin{equation}
L_{2,\rm static}^{(ABCD)}\supset
\frac{3\lambda_\rho G^3 m_A m_B m_C m_D}{r_{AB}r_{AC}r_{AD}}
+\text{relabelings of the anchor body}.
\label{eq:app-static-four-body-selected}
\end{equation}
Again, the unit-normalized coefficient is exactly one.  Thus the first higher-body static
prediction of the quadratic local mass-scaling rule is not a pure triangle term, but a
common-anchor cubic product.

After imposing the isotropic gate \(\lambda_\rho=1/2\), the selected coefficients in
\cref{eq:app-static-three-body-selected,eq:app-static-four-body-selected} are simply
halved relative to the formal \(\lambda_\rho\)-normalized hierarchy.  The pattern itself,
however, is fixed already by \cref{eq:app-static-general-many-body} and does not depend on
the test-mass matching step.

\claimclass{\Exact{} for the compact many-body formula
\cref{eq:app-static-general-many-body} once quadratic local mass scaling is declared.
\Controlled{} for the decision to interpret this hierarchy as the static extension of the
same reduced local-potential bookkeeping used at lower order.}

\subsection{Worldtube finite-size gate}
\label{app:static-finite-size}

The static hierarchy above is algebraically clean, but any claim that it represents a
universal \(1/r^4\)-order force law must be interpreted with care.  The same small-body
worldtube logic used throughout the project implies an explicit finite-size gate.

Consider a compact body of characteristic radius \(a\) moving in the field of a central
mass \(M\) at separation \(r\).  The Newtonian force scale is
\begin{equation}
F_N\sim \frac{GMm}{r^2}.
\label{eq:app-static-FN}
\end{equation}
The leading finite-size correction from a quadrupolar worldtube distortion is then of
order
\begin{equation}
F_{\rm quad}\sim F_N\left(\frac{a}{r}\right)^2.
\label{eq:app-static-Fquad}
\end{equation}
By contrast, a universal conservative \(\TwoPN\) correction scales as
\begin{equation}
F_{\rm universal}^{(2)}\sim F_N\left(\frac{GM}{rc^2}\right)^2.
\label{eq:app-static-F2PN}
\end{equation}
Therefore the ratio of finite-size to universal-\(\TwoPN\) effects is
\begin{equation}
\frac{F_{\rm quad}}{F_{\rm universal}^{(2)}}
\sim
\frac{(a/r)^2}{\left(GM/(rc^2)\right)^2}
=
\left(\frac{ac^2}{GM}\right)^2.
\label{eq:app-static-finite-size-ratio}
\end{equation}
The striking feature of \cref{eq:app-static-finite-size-ratio} is that the separation
\(r\) cancels: once one compares against a universal \(\TwoPN\) claim, the relevant gate
is no longer merely geometric smallness relative to the orbit.

The required regime is therefore
\begin{equation}
a\ll \frac{GM}{c^2},
\label{eq:app-static-finite-size-gate}
\end{equation}
not just \(a\ll r\).
This is the honest interpretation bound for any statement that the static hierarchy of the
present appendix yields a universal conservative \(1/r^4\) sector.  The algebraic
coefficient ledger remains correct once the reduction is declared, but the physical scope
of universality is narrower than a casual ``small body'' phrase might suggest.

In practice this means the static hierarchy should be read in two layers.  As an
algebraic output of quadratic local mass scaling, it is exact within the declared reduced
bookkeeping.  As a literal universal force law for finite extended bodies, it requires the
stronger gate \cref{eq:app-static-finite-size-gate}.  This distinction is why the main
text states the finite-size caveat explicitly even though it plays no role in the final
comparable-mass \(\ADM\) match.

\claimclass{\Controlled{} for the worldtube estimate
\cref{eq:app-static-finite-size-ratio,eq:app-static-finite-size-gate}.  This is not a new
coefficient solve; it is the regime statement needed to interpret the static hierarchy
honestly.}


% ============================================================
\section{Perturbative Legendre transform and \texorpdfstring{$\ADM$}{ADM} residual-solve details}
\label{app:adm-details}

This appendix records the algebra behind the main-text \(\TwoPN\) lift of
\cref{sec:adm-lift}.  There are two separate points to keep visible.
First, once the Newtonian+\(\OnePN\) ledger is frozen, the perturbative Legendre
transform simplifies enough that any \emph{new} \(\TwoPN\) Lagrangian block enters the
Hamiltonian linearly.  Second, after the repaired self/static seed is inserted, the
remaining mismatch with the standard generic-frame conservative \(\ADM\) target lives in a
small, symmetry-respecting comparable-mass basis.  The exact \(\TwoPN\) lift is therefore
an ordinary linear coefficient solve rather than an open-ended fit.

\subsection{Legendre-transform bookkeeping through \texorpdfstring{$2$}{2}PN}
\label{app:adm-legendre-details}

Write the conservative two-body Lagrangian in the bookkeeping form
\begin{equation}
L=L_0+\epsilon L_1+\epsilon^2 L_2,
\qquad
\epsilon\equiv c^{-2},
\label{eq:app-adm-epsilon-expansion}
\end{equation}
with quadratic Newtonian block
\begin{equation}
L_0=\frac12 v^T M v+V_0,
\qquad
V_0\equiv \frac{Gm_A m_B}{r},
\qquad
M\equiv \mathrm{diag}(m_A I_3,\,m_B I_3).
\label{eq:app-adm-L0-matrix}
\end{equation}
Here
\begin{equation}
v\equiv (\vct_A,\vct_B),
\qquad
p\equiv (\pvec_A,\pvec_B),
\label{eq:app-adm-vp-def}
\end{equation}
with six-vector dot products understood in the obvious block-Euclidean sense.
Because \(V_0\) is velocity-independent, the canonical momenta are
\begin{equation}
p=\frac{\partial L}{\partial v}
=Mv+\epsilon\frac{\partial L_1}{\partial v}
+\epsilon^2\frac{\partial L_2}{\partial v}.
\label{eq:app-adm-momentum-expansion}
\end{equation}
We seek the inverse relation in the form
\begin{equation}
v=v_0+\epsilon v_1+\epsilon^2 v_2+O(\epsilon^3).
\label{eq:app-adm-v-series}
\end{equation}
At leading order,
\begin{equation}
p=Mv_0,
\qquad\text{so that}\qquad
v_0=M^{-1}p=
\left(\frac{\pvec_A}{m_A},\frac{\pvec_B}{m_B}\right).
\label{eq:app-adm-v0}
\end{equation}
Now define the frozen \(\OnePN\) gradient, Hessian, and the raw \(\TwoPN\) gradient at
\(v=v_0\):
\begin{equation}
A_0\equiv \left.\frac{\partial L_1}{\partial v}\right|_{v=v_0},
\qquad
B_0\equiv \left.\frac{\partial^2 L_1}{\partial v\,\partial v}\right|_{v=v_0},
\qquad
C_0\equiv \left.\frac{\partial L_2}{\partial v}\right|_{v=v_0}.
\label{eq:app-adm-A0B0C0}
\end{equation}
Expanding \cref{eq:app-adm-momentum-expansion} order by order gives
\begin{equation}
Mv_1+A_0=0,
\qquad
Mv_2+B_0v_1+C_0=0.
\label{eq:app-adm-v1v2-eqs}
\end{equation}
Hence
\begin{equation}
v_1=-M^{-1}A_0,
\qquad
v_2=-M^{-1}(B_0v_1+C_0).
\label{eq:app-adm-v1v2}
\end{equation}

The Hamiltonian is
\begin{equation}
H=p\cdot v-L.
\label{eq:app-adm-H-def}
\end{equation}
Expanding the quadratic Newtonian part gives
\begin{align}
p\cdot v-\frac12 v^T M v
={}&
\frac12 v_0^T M v_0
-\frac{\epsilon^2}{2}\,v_1^T M v_1
+O(\epsilon^3),
\label{eq:app-adm-quadratic-piece}
\end{align}
while Taylor expansion of the higher-order blocks around \(v_0\) gives
\begin{equation}
L_1(v)=L_1(v_0)+\epsilon A_0\cdot v_1+O(\epsilon^2),
\qquad
L_2(v)=L_2(v_0)+O(\epsilon).
\label{eq:app-adm-L1L2-taylor}
\end{equation}
Substituting into \cref{eq:app-adm-H-def} yields
\begin{align}
H
={}&
\left(\frac12 v_0^T M v_0-V_0\right)
+\epsilon\bigl[-L_1(v_0)\bigr]
\nonumber\\
&+\epsilon^2\left[-L_2(v_0)-A_0\cdot v_1-\frac12 v_1^T M v_1\right]
+O(\epsilon^3).
\label{eq:app-adm-H-preliminary}
\end{align}
Using \cref{eq:app-adm-v1v2}, one has
\begin{equation}
A_0\cdot v_1=-A_0^T M^{-1}A_0,
\qquad
v_1^T M v_1=A_0^T M^{-1}A_0,
\label{eq:app-adm-v1-contractions}
\end{equation}
so the \(\epsilon^2\) block collapses to
\begin{equation}
H_2=-L_2(v_0)+\frac12 A_0^T M^{-1}A_0.
\label{eq:app-adm-H2-derived}
\end{equation}
Thus the full perturbative Legendre transform through \(\TwoPN\) is
\begin{equation}
H_0=\frac12 v_0^T M v_0-V_0,
\qquad
H_1=-L_1(v_0),
\qquad
H_2=-L_2(v_0)+\frac12 A_0^T M^{-1}A_0.
\label{eq:app-adm-H0H1H2}
\end{equation}
This is the appendix derivation of the compact main-text identity
\cref{eq:adm-H1H2-formula}.

The key structural point is that \(A_0\) depends only on the already-frozen
\(\OnePN\) block.  Therefore the quadratic correction
\(\frac12 A_0^T M^{-1}A_0\) is fixed once and for all before any new \(\TwoPN\)
comparable-mass information is introduced.  In particular, for any 
\emph{new} \(\TwoPN\) Lagrangian block \(L_{2,\mathrm{new}}\),
\begin{equation}
\Delta H_{2,\mathrm{new}}=-L_{2,\mathrm{new}}(v_0).
\label{eq:app-adm-linear-new-block}
\end{equation}
This is the algebraic reason the lift in \cref{sec:adm-lift} is linear.

For completeness, the explicit frozen \(\OnePN\) gradient entering
\cref{eq:app-adm-H2-derived} is
\begin{align}
A_{0,A}
={}&
\frac{p_A^2}{2m_A^2}\,\pvec_A
+\frac{Gm_A m_B}{r}
\left[
\frac{3}{m_A}\,\pvec_A
-\frac{7}{2m_B}\,\pvec_B
-\frac{\pvec_B\!\cdot\!\nvec}{2m_B}\,\nvec
\right],
\label{eq:app-adm-A0A}
\\
A_{0,B}
={}&
\frac{p_B^2}{2m_B^2}\,\pvec_B
+\frac{Gm_A m_B}{r}
\left[
\frac{3}{m_B}\,\pvec_B
-\frac{7}{2m_A}\,\pvec_A
-\frac{\pvec_A\!\cdot\!\nvec}{2m_A}\,\nvec
\right],
\label{eq:app-adm-A0B}
\end{align}
where \(A_0=(A_{0,A},A_{0,B})\).  We never need to expand
\(\frac12 A_0^T M^{-1}A_0\) by hand in the paper, but it is useful to record
\cref{eq:app-adm-A0A,eq:app-adm-A0B} explicitly because they show that the entire
quadratic correction is fixed by the already-carried \(\OnePN\) coefficients.

\subsection{Residual basis and linear solve}
\label{app:adm-basis}

The repaired self/static \(\TwoPN\) seed carried into the lift is
\begin{equation}
L_{2,\mathrm{self+static}}
=
\frac{m_A v_A^6+m_B v_B^6}{16}
+\frac{7Gm_A m_B}{8r}(v_A^4+v_B^4)
+\frac{2G^2m_A m_B}{r^2}(m_B v_A^2+m_A v_B^2)
+\frac{G^3m_A m_B(m_A^2+m_B^2)}{4r^3},
\label{eq:app-adm-self-static-seed}
\end{equation}
with the overall factor \(c^{-4}\) understood.  As explained in the main text, this
seed is already exact in the strict one-body sector, so any remaining mismatch with the
standard generic-frame conservative \(\ADM\) target must vanish in the test-mass limit.
That observation fixes the shape of the admissible comparable-mass ansatz.

For convenience, we restate the compact generic-frame basis used in
\cref{sec:adm-lift}.  In velocity variables, the quartic \(G/r\) sector is spanned by
\begin{align}
Q_1&=v_{AB}(v_A^2+v_B^2),
&
Q_2&=v_{An}v_{Bn}(v_A^2+v_B^2),
&
Q_3&=v_A^2v_B^2,
\nonumber\\
Q_4&=v_{AB}^2,
&
Q_5&=v_A^2v_{Bn}^2+v_B^2v_{An}^2,
&
Q_6&=v_{AB}v_{An}v_{Bn},
&
Q_7&=v_{An}^2v_{Bn}^2,
\label{eq:app-adm-Q-basis}
\end{align}
while the quadratic \(G^2/r^2\) sector is spanned by
\begin{align}
T_1&=m_B v_A^2+m_A v_B^2,
&
T_2&=m_A v_A^2+m_B v_B^2,
&
T_3&=(m_A+m_B)v_{AB},
\nonumber\\
T_4&=(m_A+m_B)v_{An}v_{Bn},
&
T_5&=m_B v_{An}^2+m_A v_{Bn}^2,
&
T_6&=m_A v_{An}^2+m_B v_{Bn}^2,
\label{eq:app-adm-T-basis}
\end{align}
and the static comparable-mass slot is
\begin{equation}
S_1\equiv \frac{G^3m_A^2m_B^2}{r^3}.
\label{eq:app-adm-S-basis}
\end{equation}
The trial added block is therefore
\begin{equation}
L_{2,\mathrm{cross}}^{\mathrm{ans}}
=\frac{Gm_A m_B}{r}\sum_{i=1}^{7}q_iQ_i
+\frac{G^2m_A m_B}{r^2}\sum_{i=1}^{6}t_iT_i
+s_1S_1.
\label{eq:app-adm-cross-ansatz}
\end{equation}

On the Hamiltonian side one simply substitutes \(v\to v_0=p/m\).  It is therefore
useful to define the induced momentum-basis monomials
\begin{equation}
\widehat Q_i\equiv Q_i(v_0),
\qquad
\widehat T_i\equiv T_i(v_0),
\qquad
\widehat S_1\equiv S_1.
\label{eq:app-adm-hatted-basis-def}
\end{equation}
Explicitly,
\begin{align}
\widehat Q_1
&=
\frac{\pvec_A\!\cdot\!\pvec_B}{m_A m_B}
\left(\frac{p_A^2}{m_A^2}+\frac{p_B^2}{m_B^2}\right),
&
\widehat Q_2
&=
\frac{(\pvec_A\!\cdot\!\nvec)(\pvec_B\!\cdot\!\nvec)}{m_A m_B}
\left(\frac{p_A^2}{m_A^2}+\frac{p_B^2}{m_B^2}\right),
\nonumber\\
\widehat Q_3
&=
\frac{p_A^2p_B^2}{m_A^2m_B^2},
&
\widehat Q_4
&=
\frac{(\pvec_A\!\cdot\!\pvec_B)^2}{m_A^2m_B^2},
\nonumber\\
\widehat Q_5
&=
\frac{p_A^2(\pvec_B\!\cdot\!\nvec)^2+p_B^2(\pvec_A\!\cdot\!\nvec)^2}{m_A^2m_B^2},
&
\widehat Q_6
&=
\frac{(\pvec_A\!\cdot\!\pvec_B)(\pvec_A\!\cdot\!\nvec)(\pvec_B\!\cdot\!\nvec)}{m_A^2m_B^2},
\nonumber\\
\widehat Q_7
&=
\frac{(\pvec_A\!\cdot\!\nvec)^2(\pvec_B\!\cdot\!\nvec)^2}{m_A^2m_B^2},
\label{eq:app-adm-hatted-Qs}
\end{align}
\begin{align}
\widehat T_1
&=
\frac{m_B}{m_A^2}p_A^2+\frac{m_A}{m_B^2}p_B^2,
&
\widehat T_2
&=
\frac{p_A^2}{m_A}+\frac{p_B^2}{m_B},
&
\widehat T_3
&=
\frac{m_A+m_B}{m_A m_B}\,\pvec_A\!\cdot\!\pvec_B,
\nonumber\\
\widehat T_4
&=
\frac{m_A+m_B}{m_A m_B}
(\pvec_A\!\cdot\!\nvec)(\pvec_B\!\cdot\!\nvec),
&
\widehat T_5
&=
\frac{m_B}{m_A^2}(\pvec_A\!\cdot\!\nvec)^2
+\frac{m_A}{m_B^2}(\pvec_B\!\cdot\!\nvec)^2,
\nonumber\\
\widehat T_6
&=
\frac{(\pvec_A\!\cdot\!\nvec)^2}{m_A}
+\frac{(\pvec_B\!\cdot\!\nvec)^2}{m_B}.
\label{eq:app-adm-hatted-Ts}
\end{align}
By \cref{eq:app-adm-linear-new-block}, the added Hamiltonian contribution is simply
\begin{equation}
H_{2,\mathrm{cross}}^{\mathrm{ans}}
=-\frac{Gm_A m_B}{r}\sum_{i=1}^{7}q_i\widehat Q_i
-\frac{G^2m_A m_B}{r^2}\sum_{i=1}^{6}t_i\widehat T_i
-s_1\widehat S_1.
\label{eq:app-adm-Hcross-ans}
\end{equation}

Now let \(H_{2,\mathrm{seed}}\) denote the \(\TwoPN\) Hamiltonian obtained by Legendre-
transforming the frozen lower-order ledger plus
\cref{eq:app-adm-self-static-seed}.  The exact generic-frame \(\ADM\) residual is
\begin{equation}
\Delta H_2
\equiv
\bigl(H_{2,\mathrm{seed}}+H_{2,\mathrm{cross}}^{\mathrm{ans}}\bigr)
-H_{2\mathrm{PN}}^{\ADM}.
\label{eq:app-adm-residual-def}
\end{equation}
Because the basis
\begin{equation}
\mathfrak B
=\{\widehat Q_1,\ldots,\widehat Q_7,\widehat T_1,\ldots,\widehat T_6,\widehat S_1\}
\label{eq:app-adm-ordered-basis}
\end{equation}
is linearly independent in a generic frame, \(\Delta H_2=0\) is equivalent to the
vanishing of each coefficient separately.  In other words, the lift reduces to a
\(14\times 14\) linear system for
\begin{equation}
c=(q_1,\ldots,q_7,t_1,\ldots,t_6,s_1)^T.
\label{eq:app-adm-c-vector}
\end{equation}
The symbolic derivation ledger solves this system exactly and uniquely, giving
\begin{align}
q_1&=-\frac74,
& q_2&=-\frac14,
& q_3&=\frac{11}{8},
& q_4&=\frac14,
\nonumber\\
q_5&=-\frac58,
& q_6&=\frac32,
& q_7&=\frac38,
\label{eq:app-adm-q-solution}
\\[0.5ex]
t_1&=0,
& t_2&=\frac{11}{8},
& t_3&=-\frac{15}{4},
& t_4&=0,
& t_5&=0,
& t_6&=\frac{15}{8},
\label{eq:app-adm-t-solution}
\\[0.5ex]
s_1&=\frac54.
\label{eq:app-adm-s-solution}
\end{align}
So three quadratic basis coefficients vanish automatically,
\begin{equation}
t_1=t_4=t_5=0.
\label{eq:app-adm-vanishing-ts}
\end{equation}
That is a useful structural output: the exact lift does not need every symmetry-allowed
quadratic monomial.

Substituting \cref{eq:app-adm-q-solution,eq:app-adm-t-solution,eq:app-adm-s-solution}
into \cref{eq:app-adm-cross-ansatz} reproduces exactly the main-text added block
\cref{eq:adm-added-cross-block}.  The solution is therefore unique not merely in the
space of all possible \(\TwoPN\) corrections, but already in the smallest compact basis
that one would naturally try first.

\subsection{Exact equality to the standard generic-frame \texorpdfstring{$\ADM$}{ADM} target}
\label{app:adm-match-details}

With the solved coefficients inserted, the trial \(\TwoPN\) Hamiltonian is
\begin{equation}
H_{2,\mathrm{trial}}
\equiv
H_{2,\mathrm{seed}}+H_{2,\mathrm{cross}}^{\mathrm{ans}}\big|_{\text{solved }q,t,s}.
\label{eq:app-adm-Htrial}
\end{equation}
By construction,
\begin{equation}
H_{2,\mathrm{trial}}-H_{2\mathrm{PN}}^{\ADM}=0.
\label{eq:app-adm-zero-residual}
\end{equation}
This is the appendix form of the main theorem used in
\cref{sec:full-assembly}: once the repaired one-body/self/static seed is frozen, the
exact comparable-mass lift closes the remaining generic-frame conservative \(\TwoPN\)
residual identically.

It is worth recording explicitly what is being cancelled.  The seed
\cref{eq:app-adm-self-static-seed} already reproduces the full strict test-mass sector,
so \cref{eq:app-adm-zero-residual} does \emph{not} rely on any retuning of the one-body
closure.  All of the new information lives in comparable-mass structures that vanish when
one body is held fixed and the other is taken test-mass light.  This is why the lift can
be interpreted as a genuine comparable-mass completion rather than as a disguised
one-body repair.

The static cross upgrade is a particularly transparent check.  The seed contributes
\begin{equation}
\frac{G^3m_A m_B}{4r^3}(m_A^2+m_B^2),
\label{eq:app-adm-static-seed}
\end{equation}
while the added cross block contributes
\begin{equation}
\frac{5G^3m_A^2m_B^2}{4r^3}
=
\frac{G^3m_A m_B}{4r^3}(5m_A m_B).
\label{eq:app-adm-static-cross}
\end{equation}
Therefore the full assembled static mass polynomial is
\begin{equation}
\frac{G^3m_A m_B}{4r^3}(m_A^2+5m_A m_B+m_B^2),
\label{eq:app-adm-static-full}
\end{equation}
which is exactly the mass structure required by the standard conservative generic-frame
\(\ADM\) target.

A second useful check is the test-mass limit.  Setting body \(B\) heavy and at rest,
all quartic cross monomials vanish immediately, the \(G^2/r^2\) comparable-mass terms
vanish after division by the light test mass, and the static cross term is likewise
suppressed by one additional power of the light mass.  Thus the whole added block drops
out in the strict one-body limit, as it must.

Finally, note that no contact transformation is being used here to hide a mismatch.
The statement \cref{eq:app-adm-zero-residual} is literal equality of the perturbatively
Legendre-transformed Hamiltonian to the standard generic-frame conservative
\(\ADM\) target in the chosen representative.  The main-text full assembled Lagrangian
therefore satisfies
\begin{equation}
L_{2,\mathrm{full}}
\xrightarrow{\ \text{Legendre}\ }
H_{2\mathrm{PN}}^{\ADM}
\label{eq:app-adm-final-arrow}
\end{equation}
with no leftover residual and no unassigned comparable-mass coefficients.

\claimclass{\Exact{} for the derivation of
\cref{eq:app-adm-H0H1H2} and for the coefficient-by-coefficient linear solve in the
independent basis \cref{eq:app-adm-ordered-basis}.  \FullMatch{} for the exact equality
\cref{eq:app-adm-zero-residual} once the carried lower-order ledger, the one-body
closure, and the self/static seed are inserted.}

% ============================================================
\section{Tensor-wake decomposition of the solved added cross block}
\label{app:tensor-wake}

This appendix gives the first constructive rewrite of the uniquely solved added
comparable-mass conservative \(\TwoPN\) cross block.
The exact \(\ADM\) residual solve fixed that block algebraically.
The point here is different: to show that, once solved, the new coefficients already
organize into a small and intelligible wake/tensor module rather than an opaque
polynomial list.

The decomposition has four parts.
First, two quartic coefficients are forced by a universal kinetic dressing of the frozen
\(\OnePN\) vector wake.
Second, the remaining quartic residual closes exactly in a minimal five-channel rank-\(2\)
tensor basis.
Third, the solved \(G^2/r^2\) velocity block is reproduced by a compact local-potential
dressing of the vector/tensor channels.
Fourth, the static cross term is a single scalar overlap.
The later appendices sharpen this further into a body-local
\(P_0\oplus P_2\) mouth-port language and then into the support-minus-closure EFT.
The tensor-wake level is the intermediate constructive layer that makes those later
rewritings possible.

As in the rest of the paper, \(L_2\) denotes the coefficient of \(c^{-4}\) in the
conservative Lagrangian expansion
\(L_{\rm cons}=L_0+c^{-2}L_1+c^{-4}L_2+O(c^{-6})\).
Accordingly, the explicit factor \(c^{-4}\) is omitted throughout this appendix.

\subsection{Solved added cross block to be decomposed}
\label{app:tensor-wake-target}

The solved added conservative two-body \(\TwoPN\) cross block is
\begin{align}
L_{2,\mathrm{cross}}
={}&
\frac{Gm_A m_B}{r}
\Big[
-\frac74\,v_{AB}(v_A^2+v_B^2)
-\frac14\,v_{An}v_{Bn}(v_A^2+v_B^2)
+\frac{11}{8}\,v_A^2v_B^2
+\frac14\,v_{AB}^2
\notag\\
&\hspace{6.4em}
-\frac58\,(v_A^2v_{Bn}^2+v_B^2v_{An}^2)
+\frac32\,v_{AB}v_{An}v_{Bn}
+\frac38\,v_{An}^2v_{Bn}^2
\Big]
\notag\\
&+
\frac{G^2m_A m_B}{r^2}
\Big[
\frac{11}{8}(m_A v_A^2+m_B v_B^2)
-\frac{15}{4}(m_A+m_B)v_{AB}
+\frac{15}{8}(m_A v_{An}^2+m_B v_{Bn}^2)
\Big]
\notag\\
&+
\frac{5G^3m_A^2m_B^2}{4r^3}.
\label{eq:tensor-wake-target}
\end{align}
Here
\begin{equation}
\mathbf n\equiv \frac{\mathbf r}{r},
\qquad
v_{AB}\equiv \mathbf v_A\!\cdot\!\mathbf v_B,
\qquad
v_{An}\equiv \mathbf v_A\!\cdot\!\mathbf n,
\qquad
v_{Bn}\equiv \mathbf v_B\!\cdot\!\mathbf n,
\label{eq:tensor-wake-basic-defs}
\end{equation}
and the local pair potentials are
\begin{equation}
U_A\equiv \frac{Gm_B}{r},
\qquad
U_B\equiv \frac{Gm_A}{r}.
\label{eq:tensor-wake-UAUB}
\end{equation}
The rest of the appendix reconstructs
\cref{eq:tensor-wake-target} from a minimal tensor-wake channel system.

\subsection{Universal kinetic dressing of the frozen \texorpdfstring{$1$}{1}PN vector wake}
\label{app:tensor-wake-vector}

The carried \(\OnePN\) vector-wake kernel is
\begin{equation}
\mathcal W^{(1)}_{AB}
\equiv
C_{\parallel}v_{AB}+C_Lv_{An}v_{Bn},
\qquad
C_{\parallel}=-\frac72,
\qquad
C_L=-\frac12.
\label{eq:tensor-wake-1pn-kernel}
\end{equation}
The first two quartic coefficients in \cref{eq:tensor-wake-target} are then exactly
\begin{equation}
-\frac74\,v_{AB}(v_A^2+v_B^2)
-\frac14\,v_{An}v_{Bn}(v_A^2+v_B^2)
=
\frac12(v_A^2+v_B^2)\,\mathcal W^{(1)}_{AB}.
\label{eq:tensor-wake-leg-dressing}
\end{equation}
So these terms are not genuinely new \(\TwoPN\) fits at all; they are a universal
leg-dressing of the frozen \(\OnePN\) wake with coefficient
\begin{equation}
\sigma=\frac12.
\label{eq:tensor-wake-sigma}
\end{equation}
It is therefore natural to define
\begin{equation}
L_{2,\mathrm{vec,kin}}
\equiv
\frac{Gm_A m_B}{r}\,\frac12(v_A^2+v_B^2)\,\mathcal W^{(1)}_{AB}.
\label{eq:tensor-wake-Lveckin}
\end{equation}
Subtracting \cref{eq:tensor-wake-Lveckin} from the quartic \(G/r\) part of
\cref{eq:tensor-wake-target} isolates the genuinely new quartic residual.

\subsection{Minimal rank-\texorpdfstring{$2$}{2} tensor basis for the quartic residual}
\label{app:tensor-wake-rank2}

Define, relative to the pair direction \(\mathbf n\), the transverse and longitudinal
scalar channels
\begin{equation}
T_A\equiv v_A^2-v_{An}^2,
\qquad
L_A\equiv v_{An}^2,
\qquad
T_B\equiv v_B^2-v_{Bn}^2,
\qquad
L_B\equiv v_{Bn}^2,
\label{eq:tensor-wake-TL}
\end{equation}
and the two genuinely rank-\(2\) cross channels
\begin{equation}
S_{AB}
\equiv
\bigl(v_{AB}-v_{An}v_{Bn}\bigr)^2-\frac12 T_A T_B,
\label{eq:tensor-wake-S}
\end{equation}
\begin{equation}
M_{AB}
\equiv
2\bigl(v_{AB}-v_{An}v_{Bn}\bigr)v_{An}v_{Bn}.
\label{eq:tensor-wake-M}
\end{equation}
The interpretations are
\begin{equation}
T:\ \text{transverse-trace scalar},
\qquad
L:\ \text{longitudinal scalar},
\qquad
S:\ \text{transverse shear},
\qquad
M:\ \text{mixed transverse--longitudinal tensor}.
\label{eq:tensor-wake-channel-meanings}
\end{equation}

With these definitions, the quartic residual after removing the universal vector-wake
leg dressing is exactly
\begin{equation}
L_{2,\mathrm{tensor,quartic}}
=
\frac{Gm_A m_B}{r}
\Big[
\frac32\,T_A T_B
+\frac14\,S_{AB}
+M_{AB}
+\frac34\,(T_A L_B+L_A T_B)
+\frac94\,L_A L_B
\Big].
\label{eq:tensor-wake-Ltensorquartic}
\end{equation}
Expanding \cref{eq:tensor-wake-Ltensorquartic} back into the velocity invariants gives
\begin{align}
L_{2,\mathrm{tensor,quartic}}
={}&
\frac{Gm_A m_B}{r}
\Big[
\frac{11}{8}v_A^2v_B^2
+\frac14 v_{AB}^2
-\frac58(v_A^2v_{Bn}^2+v_B^2v_{An}^2)
\notag\\
&\hspace{6.4em}
+\frac32 v_{AB}v_{An}v_{Bn}
+\frac38 v_{An}^2v_{Bn}^2
\Big],
\label{eq:tensor-wake-quartic-expanded}
\end{align}
which is precisely the remaining quartic part of \cref{eq:tensor-wake-target}.
Thus the genuinely new quartic \(\TwoPN\) content is an exact minimal five-channel
rank-\(2\) tensor overlap with coefficients
\begin{equation}
k_{TT}=\frac32,
\qquad
k_S=\frac14,
\qquad
k_M=1,
\qquad
k_{TL}=\frac34,
\qquad
k_{LL}=\frac94.
\label{eq:tensor-wake-kcoeffs}
\end{equation}

\subsection{Positive scalar subsector}
\label{app:tensor-wake-positive}

Inside the tensor decomposition, the scalar \((T,L)\) subsector is governed by
\begin{equation}
K_{\mathrm{scalar}}
=
\begin{pmatrix}
\frac32 & \frac34\\[4pt]
\frac34 & \frac94
\end{pmatrix}.
\label{eq:tensor-wake-Kscalar}
\end{equation}
This matrix is positive:
\begin{equation}
\det K_{\mathrm{scalar}}=\frac{45}{16}>0,
\qquad
\operatorname{tr}K_{\mathrm{scalar}}=\frac{15}{4}>0,
\label{eq:tensor-wake-Kscalar-pos}
\end{equation}
with eigenvalues
\begin{equation}
\lambda_{\pm}=\frac{15\pm 3\sqrt5}{8}.
\label{eq:tensor-wake-Kscalar-eigs}
\end{equation}
So the scalar part of the new quartic tensor sector is constructive rather than
sign-indefinite.  This is why the later diagonalization into the body-local
\(P_0\oplus P_2\) mouth-port basis is so clean.

\subsection{Local-potential dressing of the \texorpdfstring{$G^2/r^2$}{G^2/r^2} block}
\label{app:tensor-wake-local-potential}

The solved quadratic-velocity sector of \cref{eq:tensor-wake-target} is
\begin{equation}
L_{2,\mathrm{quad,vel}}
=
\frac{G^2m_A m_B}{r^2}
\Big[
\frac{11}{8}(m_Av_A^2+m_Bv_B^2)
-\frac{15}{4}(m_A+m_B)v_{AB}
+\frac{15}{8}(m_Av_{An}^2+m_Bv_{Bn}^2)
\Big].
\label{eq:tensor-wake-quad-vel}
\end{equation}
This block is reproduced exactly by the local-potential dressing
\begin{equation}
L_{2,\mathrm{quad,pot}}
=
\frac{Gm_A m_B}{r}
\Big[
\tau_{\parallel}(U_A+U_B)C_{\parallel}v_{AB}
+\beta_T(U_AT_B+U_BT_A)
+\beta_L(U_AL_B+U_BL_A)
\Big].
\label{eq:tensor-wake-quad-pot}
\end{equation}
Matching to \cref{eq:tensor-wake-quad-vel} gives the unique coefficients
\begin{equation}
\tau_{\parallel}=\frac{15}{14},
\qquad
\beta_T=\frac{11}{8},
\qquad
\beta_L=\frac{13}{4}.
\label{eq:tensor-wake-tau-beta}
\end{equation}
Inserting \cref{eq:tensor-wake-UAUB,eq:tensor-wake-1pn-kernel,eq:tensor-wake-TL,eq:tensor-wake-tau-beta}
into \cref{eq:tensor-wake-quad-pot} gives
\begin{align}
L_{2,\mathrm{quad,pot}}
={}&
\frac{G^2m_A m_B}{r^2}
\Big[
-\frac{15}{4}(m_A+m_B)v_{AB}
+\frac{11}{8}(m_Av_A^2+m_Bv_B^2)
\notag\\
&\hspace{6.8em}
+\Bigl(\frac{13}{4}-\frac{11}{8}\Bigr)(m_Av_{An}^2+m_Bv_{Bn}^2)
\Big]
\notag\\
={}&
\frac{G^2m_A m_B}{r^2}
\Big[
\frac{11}{8}(m_Av_A^2+m_Bv_B^2)
-\frac{15}{4}(m_A+m_B)v_{AB}
+\frac{15}{8}(m_Av_{An}^2+m_Bv_{Bn}^2)
\Big],
\label{eq:tensor-wake-quad-pot-expanded}
\end{align}
which is exactly \cref{eq:tensor-wake-quad-vel}.

Structurally, the solved answer is smaller than a generic ansatz would suggest:
only the parallel vector dressing is needed, and the tensor-potential dressing is diagonal
in the transverse and longitudinal projector channels.

\subsection{Static scalar overlap and full constructive reconstruction}
\label{app:tensor-wake-full}

The remaining purely static comparable-mass piece is
\begin{equation}
L_{2,\mathrm{static,cross}}
=
\frac{5G^3m_A^2m_B^2}{4r^3}
=
\frac{Gm_A m_B}{r}\,\frac54 U_AU_B.
\label{eq:tensor-wake-static-cross}
\end{equation}
Therefore the full solved added cross block can be reconstructed exactly as
\begin{equation}
L_{2,\mathrm{cross}}
=
L_{2,\mathrm{vec,kin}}
+L_{2,\mathrm{tensor,quartic}}
+L_{2,\mathrm{quad,pot}}
+L_{2,\mathrm{static,cross}}.
\label{eq:tensor-wake-full-decomp}
\end{equation}
By direct expansion, the right-hand side of
\cref{eq:tensor-wake-full-decomp} reproduces \cref{eq:tensor-wake-target} term by term.

This gives the tensor-wake interpretation of the solved added cross block:
\begin{enumerate}[leftmargin=1.5em]
  \item a carried-forward \(\OnePN\) vector wake with universal kinetic dressing,
  \item a genuinely new five-channel rank-\(2\) tensor residual,
  \item a compact local-potential dressing of the vector/tensor projector channels, and
  \item a single scalar static overlap.
\end{enumerate}
In this language the algebraically solved \(\ADM\) residual is already telling us what a
constructive inner-throat derivation should reproduce.

\subsection{Bridge to the mouth-port language}
\label{app:tensor-wake-bridge}

The tensor-wake basis is not yet the final throat-response basis.
It is the most economical intermediate basis for seeing how the solved quartic residual
organizes before one commits to explicit mouth ports.
The next appendix shows that the same scalar tensor block diagonalizes into
\(P_0\oplus P_2\) channels.
At the scalar level one writes
\begin{equation}
P_0=T+L=v^2,
\qquad
P_{20}=-T+2L=3(v\cdot n)^2-v^2,
\label{eq:tensor-wake-P0P20}
\end{equation}
and the scalar kernel \cref{eq:tensor-wake-Kscalar} diagonalizes exactly in that basis.
So the tensor-wake decomposition should be read as the constructive bridge between
``the coefficients solved by the \(\ADM\) lift'' and ``the small
\(P_0\oplus P_2\) mouth-port hierarchy'' used in the next appendices.

\claimclass{\Exact{} for the algebraic decomposition of the solved added cross block into
\cref{eq:tensor-wake-full-decomp} once the target coefficients are inserted.
\Controlled{} for the interpretation of these exact overlaps as a constructive wake/tensor
response module that the inner-throat derivation should reproduce.
This appendix is explanatory rather than independent: it does not alter the solved
coefficients, but reorganizes them into a minimal carried-vector plus rank-\(2\) tensor
channel structure.}

% ============================================================
\section{\texorpdfstring{$P_0\oplus P_2$}{P0⊕P2} mouth-port rebuild and support-minus-closure EFT}
\label{app:mouth-port}

This appendix continues the constructive program started in
\cref{app:tensor-wake}.  There the solved added conservative two-body \(\TwoPN\)
comparable-mass cross block was rewritten as a carried \(\OnePN\) vector wake plus a
minimal rank-\(2\) tensor residual and a compact local-potential dressing.
The purpose of the present appendix is to sharpen that constructive rewrite into a
body-local mouth-port language and then into a small signed auxiliary-channel EFT.

Two statements are proved here.
First, after the universal leg-dressing piece inherited from the frozen
\(\OnePN\) wake is removed, the genuinely new quartic residual diagonalizes exactly
into one real monopole-like port and the full real quadrupole multiplet,
\(P_0\oplus P_2\).
Second, once the solved \(U\times \Pi\) terms and the static cross term are restored,
the new even block factorizes as a positive support bundle minus a single pure-potential
closure channel.
That is the sense in which the solved added \(\TwoPN\) cross sector becomes a genuine
\emph{support-minus-closure throat response theory} rather than an opaque coefficient list.

As in the preceding appendix, \(L_2\) denotes the coefficient of \(c^{-4}\) in the
conservative expansion
\(L_{\rm cons}=L_0+c^{-2}L_1+c^{-4}L_2+O(c^{-6})\),
so the explicit overall factor \(c^{-4}\) is suppressed throughout.

\subsection{Body-local \texorpdfstring{$P_0\oplus P_2$}{P0⊕P2} mouth-port basis}
\label{app:mouth-port-basis}

To display the real mouth-port channels it is convenient to work in the pair frame with
\begin{equation}
\mathbf n\equiv \frac{\mathbf r}{r}=\hat{\mathbf z}.
\label{eq:mouth-port-nz}
\end{equation}
Write each velocity as
\begin{equation}
\mathbf v_A=\mathbf u_A+d_A\mathbf n,
\qquad
\mathbf v_B=\mathbf u_B+d_B\mathbf n,
\qquad
\mathbf u_A\!\cdot\!\mathbf n=\mathbf u_B\!\cdot\!\mathbf n=0,
\label{eq:mouth-port-velocity-split}
\end{equation}
so that \(d_A=v_{An}\) and \(d_B=v_{Bn}\).
The real transverse components are denoted
\begin{equation}
\mathbf u_A=(u_{Ax},u_{Ay}),
\qquad
\mathbf u_B=(u_{Bx},u_{By}).
\label{eq:mouth-port-transverse-components}
\end{equation}
Rotational covariance guarantees that the final reconstruction is independent of the
choice of pair frame; the \(\hat z\) presentation is only a convenient way to write the
real channels explicitly.

The frozen \(\OnePN\) wake is already a dipole kernel,
\begin{equation}
L^{\rm wake}_{1\mathrm{PN}}
=-\frac72\,\mathbf v_A\!\cdot\!\mathbf v_B-\frac12 d_A d_B
=-\Bigl[\Pi_A^{(1\pm)}\!\cdot\!\Pi_B^{(1\pm)}+\Pi_A^{(10)}\Pi_B^{(10)}\Bigr],
\label{eq:mouth-port-1pn-wake}
\end{equation}
with the body-local odd dipole ports
\begin{equation}
\Pi_A^{(1\pm)}=\sqrt{\frac72}\,\mathbf u_A,
\qquad
\Pi_A^{(10)}=2d_A,
\label{eq:mouth-port-odd-ports-A}
\end{equation}
(and similarly for body \(B\)).
So the \(\OnePN\) sector already carries a local \(\ell=1\) mouth-response language.
The new question is how the genuinely new quartic residual organizes after the universal
kinetic leg dressing
\begin{equation}
\frac12(v_A^2+v_B^2)L^{\rm wake}_{1\mathrm{PN}}
\label{eq:mouth-port-leg-dressing}
\end{equation}
is removed.

The scalar tensor block derived in \cref{app:tensor-wake} used the variables
\begin{equation}
T=u^2,
\qquad
L=d^2,
\label{eq:mouth-port-TL}
\end{equation}
with scalar kernel
\begin{equation}
K_{TL}=
\begin{pmatrix}
\frac32 & \frac34\\[3pt]
\frac34 & \frac94
\end{pmatrix}.
\label{eq:mouth-port-KTL}
\end{equation}
Introducing the scalar combinations
\begin{equation}
P_0=T+L=v^2,
\qquad
P_{20}=-T+2L=3(v\!\cdot\!n)^2-v^2,
\label{eq:mouth-port-P0P20}
\end{equation}
this kernel diagonalizes exactly to
\begin{equation}
K_{P_0P_{20}}=
\begin{pmatrix}
\frac54 & 0\\[3pt]
0 & \frac14
\end{pmatrix}.
\label{eq:mouth-port-Kdiag}
\end{equation}
This immediately suggests the canonically normalized scalar ports
\begin{equation}
\Pi^{(0)}=\sqrt{\frac54}\,P_0=\frac{\sqrt5}{2}v^2,
\qquad
\Pi^{(20)}=\sqrt{\frac14}\,P_{20}=\frac12\bigl(3(v\!\cdot\!n)^2-v^2\bigr).
\label{eq:mouth-port-scalar-ports}
\end{equation}

The full real quadrupole multiplet is then completed by
\begin{equation}
\Pi^{(21c)}=\sqrt2\,(v\!\cdot\!n)u_x,
\qquad
\Pi^{(21s)}=\sqrt2\,(v\!\cdot\!n)u_y,
\label{eq:mouth-port-21-ports}
\end{equation}
\begin{equation}
\Pi^{(22c)}=\frac{u_x^2-u_y^2}{2\sqrt2},
\qquad
\Pi^{(22s)}=\frac{2u_xu_y}{2\sqrt2}.
\label{eq:mouth-port-22-ports}
\end{equation}
Using these definitions for each body, the quartic residual left after subtracting the
universal carried-forward piece \cref{eq:mouth-port-leg-dressing} is exactly
\begin{align}
Q^{\rm new}_{2\mathrm{PN}}
={}&
\Pi_A^{(0)}\Pi_B^{(0)}
+\Pi_A^{(20)}\Pi_B^{(20)}
+\Pi_A^{(21c)}\Pi_B^{(21c)}
+\Pi_A^{(21s)}\Pi_B^{(21s)}
\notag\\
&
+\Pi_A^{(22c)}\Pi_B^{(22c)}
+\Pi_A^{(22s)}\Pi_B^{(22s)}.
\label{eq:mouth-port-quartic-residual}
\end{align}
So the genuinely new quartic \(\TwoPN\) tensor sector is a positive Gram overlap of one
real \(P_0\) channel and the full real \(P_2\) multiplet.  This is the exact
body-local mouth-port version of the tensor-wake residual from the previous appendix.

The solved \(G^2/r^2\) velocity block also becomes simpler in this basis.  Using the
local pair potentials
\begin{equation}
U_A\equiv \frac{Gm_B}{r},
\qquad
U_B\equiv \frac{Gm_A}{r},
\label{eq:mouth-port-UAUB}
\end{equation}
it can be written as
\begin{equation}
-\frac{15}{4}(U_A+U_B)(\mathbf v_A\!\cdot\!\mathbf v_B)
+U_A\left(\frac{4}{\sqrt5}\Pi_B^{(0)}+\frac54\Pi_B^{(20)}\right)
+U_B\left(\frac{4}{\sqrt5}\Pi_A^{(0)}+\frac54\Pi_A^{(20)}\right).
\label{eq:mouth-port-upi-block}
\end{equation}
So only the scalar axisymmetric ports \(P_0\) and \(P_{20}\) are driven directly by the
local-potential source terms.

Finally, the static added cross piece is simply
\begin{equation}
\frac54 U_AU_B.
\label{eq:mouth-port-static-cross}
\end{equation}
Equations
\cref{eq:mouth-port-leg-dressing,eq:mouth-port-quartic-residual,eq:mouth-port-upi-block,eq:mouth-port-static-cross}
therefore reconstruct the full added conservative \(\TwoPN\) cross block in an explicit
body-local mouth-port basis.

\subsection{Support-minus-closure factorization}
\label{app:mouth-port-support}

The new even sector becomes especially transparent once the solved
\(U\times\Pi\) coefficients are absorbed into shifted body-local support channels.
Define
\begin{equation}
\mathcal S_A=
\Bigl(
\Pi_A^{(0)}+\frac{4}{\sqrt5}U_A,
\Pi_A^{(20)}+\frac54 U_A,
\Pi_A^{(21c)},
\Pi_A^{(21s)},
\Pi_A^{(22c)},
\Pi_A^{(22s)}
\Bigr),
\label{eq:mouth-port-support-A}
\end{equation}
with the obvious analogous definition for \(\mathcal S_B\).
Then the positive support channels alone would generate a pure static coefficient
\begin{equation}
\left(\frac{4}{\sqrt5}\right)^2+\left(\frac54\right)^2
=\frac{381}{80}.
\label{eq:mouth-port-support-static}
\end{equation}
But the solved static coefficient in \cref{eq:mouth-port-static-cross} is only
\begin{equation}
\frac54=\frac{100}{80}.
\label{eq:mouth-port-target-static}
\end{equation}
So the exact closure deficit is
\begin{equation}
\Delta_{\rm geom}=\frac{381}{80}-\frac{100}{80}=\frac{281}{80}.
\label{eq:mouth-port-delta-geom}
\end{equation}
This deficit is removed by one negative closure channel,
\begin{equation}
\mathcal C_A=\sqrt{\frac{281}{80}}\,U_A,
\qquad
\mathcal C_B=\sqrt{\frac{281}{80}}\,U_B.
\label{eq:mouth-port-closure-channel}
\end{equation}

The entire even block is therefore exactly
\begin{equation}
L^{\rm even}_{2\mathrm{PN}}=\mathcal S_A\!\cdot\!\mathcal S_B-\mathcal C_A\mathcal C_B.
\label{eq:mouth-port-even-factorized}
\end{equation}
The important structural fact is that the negative direction lives entirely in the pure
potential \(U\)-direction.  It does \emph{not} mix with the genuine mouth ports.
That is why it is natural to interpret the indefinite part as a geometry-closure term
rather than as a pathology of the mouth-response sector itself.

In the ordered source basis
\begin{equation}
X=(\Pi^{(0)},\Pi^{(20)},\Pi^{(21c)},\Pi^{(21s)},\Pi^{(22c)},\Pi^{(22s)},U),
\label{eq:mouth-port-Xbasis}
\end{equation}
the exact even response matrix is
\begin{equation}
M_{\rm even}=
\begin{pmatrix}
1&0&0&0&0&0&4/\sqrt5\\
0&1&0&0&0&0&5/4\\
0&0&1&0&0&0&0\\
0&0&0&1&0&0&0\\
0&0&0&0&1&0&0\\
0&0&0&0&0&1&0\\
4/\sqrt5&5/4&0&0&0&0&5/4
\end{pmatrix}.
\label{eq:mouth-port-Meven}
\end{equation}
It factorizes exactly as
\begin{equation}
M_{\rm even}=R_{\rm support}^{T}R_{\rm support}-R_{\rm geom}^{T}R_{\rm geom},
\label{eq:mouth-port-matrix-factorization}
\end{equation}
with
\begin{equation}
R_{\rm support}=
\begin{pmatrix}
1&0&0&0&0&0&4/\sqrt5\\
0&1&0&0&0&0&5/4\\
0&0&1&0&0&0&0\\
0&0&0&1&0&0&0\\
0&0&0&0&1&0&0\\
0&0&0&0&0&1&0
\end{pmatrix},
\qquad
R_{\rm geom}=\begin{pmatrix}0&0&0&0&0&0&\sqrt{281/80}\end{pmatrix}.
\label{eq:mouth-port-Rsupport-Rgeom}
\end{equation}
So the even block is literally a positive-semidefinite mouth-response block minus a
rank-\(1\) geometry-closure block.

\subsection{Local auxiliary-channel EFT and exact on-shell reconstruction}
\label{app:mouth-port-eft}

The factorization above has an immediate local auxiliary-field interpretation.
For a positive support channel \(q\) with body sources \(S_A,S_B\), use
\begin{equation}
L_q=-\frac12 q^2+q(S_A+S_B).
\label{eq:mouth-port-aux-positive}
\end{equation}
Here the auxiliary variable \(q\) is a local EFT support channel and is unrelated
to electric charge or to the gravity-side coefficient \(\krho\).
Eliminating \(q\) gives the cross term \(+S_AS_B\).
For the negative closure channel \(g\) with sources \(C_A,C_B\), use
\begin{equation}
L_g=+\frac12 g^2+g(C_A+C_B).
\label{eq:mouth-port-aux-negative}
\end{equation}
Eliminating \(g\) gives the cross term \(-C_AC_B\).
Finally, for the odd dipole wake channel \(d\) with body sources \(D_A,D_B\), use
\begin{equation}
L_d=-\frac12 d^2+d(D_A-D_B),
\label{eq:mouth-port-aux-odd}
\end{equation}
which yields the odd cross term \(-D_AD_B\) on shell.
Thus the solved \(\TwoPN\) sector is exactly the on-shell result of a small signed
auxiliary-channel EFT.

It is useful to keep the odd channels visible explicitly.  Define the odd dipole bundle
\begin{equation}
\mathcal D_A=
\Bigl(\sqrt{\tfrac72}v_{Ax},\sqrt{\tfrac72}v_{Ay},2v_{Az}\Bigr),
\qquad
\mathcal D_B=
\Bigl(\sqrt{\tfrac72}v_{Bx},\sqrt{\tfrac72}v_{By},2v_{Bz}\Bigr),
\label{eq:mouth-port-Dodd}
\end{equation}
so that the frozen \(\OnePN\) wake is
\begin{equation}
L^{\rm wake}_{1\mathrm{PN}}=-\mathcal D_A\!\cdot\!\mathcal D_B.
\label{eq:mouth-port-Dodd-wake}
\end{equation}
The added odd \(\TwoPN\) block is then simply the universal kinetic dressing plus the
solved local-potential dressing from \cref{app:tensor-wake}, namely
\begin{equation}
L^{\rm add}_{\rm odd,dressed}
=\frac12(v_A^2+v_B^2)L^{\rm wake}_{1\mathrm{PN}}
-\frac{15}{4}(U_A+U_B)(\mathbf v_A\!\cdot\!\mathbf v_B).
\label{eq:mouth-port-Lodd-dressed}
\end{equation}
Putting everything together, the entire solved added conservative \(\TwoPN\) cross block
is reconstructed exactly as
\begin{equation}
L^{\rm add}_{2\mathrm{PN},\,\mathrm{cross}}
=\frac{Gm_A m_B}{r}
\Bigl[
L^{\rm add}_{\rm odd,dressed}
+\mathcal S_A\!\cdot\!\mathcal S_B
-\mathcal C_A\mathcal C_B
\Bigr].
\label{eq:mouth-port-full-factorized}
\end{equation}
By direct expansion,
\cref{eq:mouth-port-full-factorized} reproduces the full solved added \(\ADM\)-lifted
cross target term by term.

\subsection{Why this is the right constructive split}
\label{app:mouth-port-interpretation}

This appendix matters because it is the first point where the new \(\TwoPN\) cross
sector stops looking like ``a polynomial found by the \(\ADM\) lift'' and starts
looking like a real throat-response theory.
The constructive split is now very sharp:
\begin{enumerate}[leftmargin=1.5em]
  \item the frozen \(\OnePN\) wake survives as an odd dipole backbone,
  \item the genuinely new even \(\TwoPN\) sector is a positive
  \(P_0\oplus P_2\) mouth-response bundle, and
  \item the only negative piece is a single pure-potential geometry-closure channel.
\end{enumerate}
That last point is especially important.  Because the negative direction lives entirely
in the \(U\)-direction, the present factorization does \emph{not} say that the mouth
operator itself is indefinite.  It says something more controlled: the support sector is
positive and minimal, while a separate geometry-energy closure term subtracts the excess
static overlap needed to reach the exact solved coefficient.  This is why the next
appendices can cleanly separate ``mouth response'' from ``geometry closure'' and then
package the result into the minimal irreducible-channel throat operator.

In this language, the remaining constructive task is no longer vague.  The inner-throat
program must explain two specific pieces:
\begin{enumerate}[leftmargin=1.5em]
  \item why the positive support bundle closes exactly on the
  \(P_0\oplus P_2\) channels with the solved source shifts
  \(4/\sqrt5\) and \(5/4\), and
  \item why geometry contributes the single rank-\(1\) pure-\(U\) closure deficit
  \(\Delta_{\rm geom}=281/80\).
\end{enumerate}
So the support-minus-closure EFT is not merely another algebraic repackaging.  It is the
first form in which the solved added \(\TwoPN\) cross block becomes a small,
body-local throat-response theory with a clear division between mouth support and
geometry closure.

\claimclass{\Exact{} for the algebraic equivalence of
\cref{eq:mouth-port-full-factorized} with the solved added conservative two-body
\(\TwoPN\) cross block once the coefficients are inserted.  \Controlled{} for the
interpretation of the positive block as a minimal mouth-response sector and of the
rank-\(1\) pure-\(U\) subtraction as a geometry-closure channel.  This appendix is
constructive and explanatory: it does not introduce new fitted coefficients, but reveals
that the solved \(\ADM\) cross block is already the on-shell output of a small signed
auxiliary-channel EFT.}

% ============================================================
\section{Minimal irreducible-channel throat operator and dynamic \texorpdfstring{$\DtN$}{DtN} scaffolding}
\label{app:irred-dtn}

This appendix turns the constructive rewrites of
\cref{app:tensor-wake,app:mouth-port}
into the smallest actual low-frequency throat operator that reproduces the solved added
conservative two-body \(\TwoPN\) cross target.  The point is no longer only that the
solved coefficients admit a useful decomposition.  The stronger statement is that they
already determine a minimal channel content for the throat-response problem.

At zero frequency, the full added conservative \(\TwoPN\) cross block is reproduced by
three ingredients:
\begin{enumerate}[leftmargin=1.5em]
  \item a carried odd \(\ell=1\) dipole sector with distinct transverse and
  longitudinal residues,
  \item a new even \(\ell=0\oplus 2\) support bundle with canonical unit support
  residues, and
  \item one direct pure-potential geometry-closure channel.
\end{enumerate}
Once that static operator is identified, the remaining dynamic problem becomes much
sharper.  One must explain why the zero-frequency residues take the values already solved
by the conservative \(\TwoPN\) algebra and then derive the corresponding low-frequency
pole structure from the inner-throat PDE / \(\DtN\) problem.

The appendix is organized accordingly.  First, we record the minimal irreducible-channel
operator fixed by the solved static target.  Next, we show that the same operator already
predicts the first local-potential three-body lift.  Finally, we build the minimal
frequency-dependent \(\DtN\) scaffolding: an isotropic 4D-ball unit test, which fails in
exactly the ways expected from the algebra, and a one-pole-per-channel axisymmetric
completion, which reduces to the solved static operator as \(\omega\to 0\).

As elsewhere in the paper,
\begin{equation}
L_{\rm cons}=L_0+\frac{1}{c^2}L_1+\frac{1}{c^4}L_2+O(c^{-6}),
\label{eq:irred-dtn-pn-expansion}
\end{equation}
so the explicit overall factor \(c^{-4}\) is suppressed throughout this appendix.

\subsection{Minimal irreducible channel content at zero frequency}
\label{app:irred-dtn-static}

Work in the pair frame
\begin{equation}
\mathbf n\equiv \frac{\mathbf r}{r}=\hat{\mathbf z},
\qquad
\mathbf v_A=\mathbf u_A+d_A\mathbf n,
\qquad
\mathbf v_B=\mathbf u_B+d_B\mathbf n,
\label{eq:irred-dtn-pair-frame}
\end{equation}
with \(\mathbf u_A\!\cdot\!\mathbf n=\mathbf u_B\!\cdot\!\mathbf n=0\) and
\begin{equation}
U_A\equiv \frac{Gm_B}{r},
\qquad
U_B\equiv \frac{Gm_A}{r}.
\label{eq:irred-dtn-UAUB}
\end{equation}
The frozen odd dipole wake may be written as
\begin{equation}
L^{\rm wake}_{1\mathrm{PN}}
=-\frac72\,\mathbf u_A\!\cdot\!\mathbf u_B-4d_A d_B,
\label{eq:irred-dtn-1pn-odd}
\end{equation}
which already displays two inequivalent odd channels:
\begin{equation}
1\perp:\ \mathbf u,
\qquad
10:\ d.
\label{eq:irred-dtn-odd-channel-labels}
\end{equation}
The solved added odd conservative \(\TwoPN\) block is then
\begin{equation}
L^{\rm add}_{\rm odd}
=
\frac12(v_A^2+v_B^2)L^{\rm wake}_{1\mathrm{PN}}
-(U_A+U_B)
\left[
\frac{15}{4}\,\mathbf u_A\!\cdot\!\mathbf u_B
+\frac{15}{4}\,d_A d_B
\right].
\label{eq:irred-dtn-added-odd}
\end{equation}
Dividing the solved potential dressings by the frozen static odd residues gives the body
source dressings
\begin{equation}
\eta_{\perp}=\frac{15}{14},
\qquad
\eta_{\parallel}=\frac{15}{16},
\label{eq:irred-dtn-odd-eta}
\end{equation}
so the odd sector is fully fixed at conservative zero frequency.

The even sector comes from the body-local
\(P_0\oplus P_2\) mouth-port basis already constructed in
\cref{app:mouth-port}.  Using the six real ports
\begin{equation}
\Pi
\equiv
\bigl(
\Pi^{(0)},\Pi^{(20)},\Pi^{(21c)},\Pi^{(21s)},\Pi^{(22c)},\Pi^{(22s)}
\bigr),
\label{eq:irred-dtn-pi-basis}
\end{equation}
the solved even block can be written as
\begin{equation}
L^{\rm even}_{2\mathrm{PN}}
=(\Pi_A+JU_A)\!\cdot\!(\Pi_B+JU_B)-\Delta_{\rm geom}U_AU_B,
\label{eq:irred-dtn-even-zero}
\end{equation}
with source vector
\begin{equation}
J=\left(\frac{4}{\sqrt5},\frac54,0,0,0,0\right)
\label{eq:irred-dtn-J}
\end{equation}
and geometry-closure deficit
\begin{equation}
\Delta_{\rm geom}=\frac{281}{80}.
\label{eq:irred-dtn-deltageom}
\end{equation}
The support part alone contributes the static coefficient
\begin{equation}
J\cdot J
=
\left(\frac{4}{\sqrt5}\right)^2
+\left(\frac54\right)^2
=\frac{381}{80},
\label{eq:irred-dtn-JdotJ}
\end{equation}
so the direct geometry subtraction in
\cref{eq:irred-dtn-deltageom} is exactly what lowers the pure-static coefficient to the
solved target value
\begin{equation}
\frac{381}{80}-\frac{281}{80}=\frac54.
\label{eq:irred-dtn-static-target}
\end{equation}

This yields the canonical zero-frequency channel data:
\begin{equation}
R_{1\perp}=\frac72,
\qquad
R_{10}=4,
\label{eq:irred-dtn-odd-residues}
\end{equation}
\begin{equation}
R_0=R_{20}=R_{21}=R_{22}=1,
\label{eq:irred-dtn-even-residues}
\end{equation}

together with the scalar source vector \cref{eq:irred-dtn-J} and the closure deficit
\cref{eq:irred-dtn-deltageom}.  In other words, the smallest irreducible channel set
required by the solved conservative added \(\TwoPN\) cross block is
\begin{equation}
\{1\perp,10,0,20,21c,21s,22c,22s,g\},
\label{eq:irred-dtn-minimal-channel-set}
\end{equation}
where \(g\) denotes the single pure-potential geometry-closure channel.

A compact way to summarize the full added conservative cross block is therefore
\begin{equation}
L^{\rm add}_{2\mathrm{PN},\,\mathrm{cross}}
=
\frac{Gm_A m_B}{r}
\left[
L^{\rm add}_{\rm odd}
+(\Pi_A+JU_A)\!\cdot\!(\Pi_B+JU_B)
-\Delta_{\rm geom}U_AU_B
\right].
\label{eq:irred-dtn-full-static-operator}
\end{equation}
Equation~\eqref{eq:irred-dtn-full-static-operator} is the minimal irreducible-channel
throat operator at zero frequency.

\subsection{Immediate local-potential \texorpdfstring{$3$}{3}-body lift}
\label{app:irred-dtn-3body}

The operator above already predicts the first local-potential three-body lift.  Promote
body-local pair potentials to
\begin{equation}
U_A^{\rm loc}=G\left(\frac{m_B}{r_{AB}}+\frac{m_C}{r_{AC}}\right),
\qquad
U_B^{\rm loc}=G\left(\frac{m_A}{r_{AB}}+\frac{m_C}{r_{BC}}\right).
\label{eq:irred-dtn-local-potentials-3body}
\end{equation}
Then the pair-\(AB\) operator immediately induces a background-\(C\) three-body lift.
The velocity-dependent piece is
\begin{align}
L^{(3\text{-body, vel})}_{AB|C}
={}&
\frac{G^2m_A m_B m_C}{8r_{AB}}
\left[
\frac{11v_B^2+15d_B^2-30(\mathbf v_A\!\cdot\!\mathbf v_B)}{r_{AC}}
+
\frac{11v_A^2+15d_A^2-30(\mathbf v_A\!\cdot\!\mathbf v_B)}{r_{BC}}
\right],
\label{eq:irred-dtn-3body-vel}
\end{align}
with
\begin{equation}
d_A\equiv \mathbf v_A\!\cdot\!\mathbf n_{AB},
\qquad
d_B\equiv \mathbf v_B\!\cdot\!\mathbf n_{AB},
\label{eq:irred-dtn-dA-dB-3body}
\end{equation}
and the static piece is
\begin{equation}
L^{(3\text{-body, stat})}_{AB|C}
=
\frac{5G^3m_A m_B m_C}{4}
\left[
\frac{m_A}{r_{AB}^2r_{AC}}
+
\frac{m_B}{r_{AB}^2r_{BC}}
+
\frac{m_C}{r_{AB}r_{AC}r_{BC}}
\right].
\label{eq:irred-dtn-3body-static}
\end{equation}
These terms are not used in the main two-body \(\TwoPN\) theorem, but they matter
constructively: once the minimal throat operator is fixed at zero frequency, it already
carries nontrivial \(N\)-body content.

\subsection{Isotropic \texorpdfstring{$4$}{4}D-ball \texorpdfstring{$\DtN$}{DtN} unit test}
\label{app:irred-dtn-ball}

The cleanest PDE-side unit test is the regular interior Helmholtz problem in a 4D ball.
For harmonic degree \(\ell\), the regular radial mode is
\begin{equation}
u_{\ell}(r)\propto r^{-1}J_{\ell+1}(kr),
\qquad
z\equiv ka,
\label{eq:irred-dtn-ball-radial}
\end{equation}
so the boundary Dirichlet-to-Neumann eigenvalue at radius \(a\) is
\begin{equation}
\Lambda_{\ell}(z)
=
-\frac1a
+\frac{z}{a}\frac{J'_{\ell+1}(z)}{J_{\ell+1}(z)}.
\label{eq:irred-dtn-ball-Lambda}
\end{equation}
Its small-\(z\) expansions for the first three channels are
\begin{equation}
\Lambda_0(z)
=-\frac{z^2}{4a}-\frac{z^4}{96a}+O(z^6),
\label{eq:irred-dtn-ball-Lambda0}
\end{equation}
\begin{equation}
\Lambda_1(z)
=\frac1a-\frac{z^2}{6a}-\frac{z^4}{288a}+O(z^6),
\label{eq:irred-dtn-ball-Lambda1}
\end{equation}
\begin{equation}
\Lambda_2(z)
=\frac2a-\frac{z^2}{8a}-\frac{z^4}{640a}+O(z^6).
\label{eq:irred-dtn-ball-Lambda2}
\end{equation}
Accordingly, the first admittances are
\begin{equation}
Y_1(z)=\frac1{\Lambda_1(z)}=a+\frac{az^2}{6}+\frac{az^4}{32}+O(z^6),
\label{eq:irred-dtn-ball-Y1}
\end{equation}
\begin{equation}
Y_2(z)=\frac1{\Lambda_2(z)}=\frac a2+\frac{az^2}{32}+\frac{3az^4}{1280}+O(z^6).
\label{eq:irred-dtn-ball-Y2}
\end{equation}

This isotropic branch is useful precisely because it fails in the right way.
It gets several structural features right: it is diagonal in harmonic channels, it
supports positive \(\ell=1\) and \(\ell=2\) towers, and after canonical normalization it is
compatible with an identity-like support metric within a fixed \(\ell\)-sector.  But it
also fails two requirements of the solved conservative operator immediately.

First, there is no finite static monopole stiffness:
\begin{equation}
\Lambda_0(0)=0.
\label{eq:irred-dtn-ball-no-monopole}
\end{equation}
So a pure isotropic cavity cannot generate the solved finite static monopole support or
source structure.  Second, a spherical branch cannot split the odd dipole residues.  It
contains one degenerate \(\ell=1\) multiplet rather than the two distinct static values
\cref{eq:irred-dtn-odd-residues}.  Therefore the isotropic 4D-ball branch is a useful
no-go unit test, but not the full throat.

\subsection{Minimal axisymmetric pole-completed \texorpdfstring{$\DtN$}{DtN} scaffolding}
\label{app:irred-dtn-axisymmetric}

The static operator above therefore forces the minimal PDE-side symmetry reduction
\begin{equation}
O(3)\longrightarrow O(2),
\label{eq:irred-dtn-O3O2}
\end{equation}
which splits the dipole and quadrupole multiplets exactly as required by the solved
conservative channel structure.
The smallest low-frequency completion consistent with the zero-frequency data is then a
one-pole-per-channel model with admittances
\begin{equation}
Y_{1\perp}(\omega)=\frac{7/2}{1-\omega^2/\Omega_{1\perp}^2},
\qquad
Y_{10}(\omega)=\frac{4}{1-\omega^2/\Omega_{10}^2},
\label{eq:irred-dtn-Yodd}
\end{equation}
\begin{equation}
Y_0(\omega)=\frac{1}{1-\omega^2/\Omega_0^2},
\qquad
Y_{20}(\omega)=\frac{1}{1-\omega^2/\Omega_{20}^2},
\label{eq:irred-dtn-Y0Y20}
\end{equation}
\begin{equation}
Y_{21}(\omega)=\frac{1}{1-\omega^2/\Omega_{21}^2},
\qquad
Y_{22}(\omega)=\frac{1}{1-\omega^2/\Omega_{22}^2},
\qquad
Y_g(\omega)=\frac{1}{1-\omega^2/\Omega_g^2}.
\label{eq:irred-dtn-Y21Y22Yg}
\end{equation}
Here \(Y_{21}\) stands for both real \(|m|=1\) quadrupole components and
\(Y_{22}\) for both real \(|m|=2\) components.  The corresponding bare \(\DtN\) kernels are
simply the inverses,
\begin{equation}
Z_{\alpha}(\omega)=Y_{\alpha}(\omega)^{-1},
\label{eq:irred-dtn-Zinverse}
\end{equation}
for \(\alpha\in\{1\perp,10,0,20,21,22,g\}\).  Explicitly,
\begin{equation}
Z_{1\perp}(\omega)=\frac{1-\omega^2/\Omega_{1\perp}^2}{7/2},
\qquad
Z_{10}(\omega)=\frac{1-\omega^2/\Omega_{10}^2}{4},
\label{eq:irred-dtn-Zodd}
\end{equation}
with analogous formulas for the even and geometry channels.

In the ordered source basis
\begin{equation}
X=(P_0,P_{20},P_{21c},P_{21s},P_{22c},P_{22s},U),
\label{eq:irred-dtn-Xbasis}
\end{equation}
the frequency-dependent even response matrix becomes
\begin{equation}
M_{\rm even}(\omega)=
\begin{pmatrix}
Y_0 & 0 & 0 & 0 & 0 & 0 & J_0Y_0\\
0 & Y_{20} & 0 & 0 & 0 & 0 & J_{20}Y_{20}\\
0 & 0 & Y_{21} & 0 & 0 & 0 & 0\\
0 & 0 & 0 & Y_{21} & 0 & 0 & 0\\
0 & 0 & 0 & 0 & Y_{22} & 0 & 0\\
0 & 0 & 0 & 0 & 0 & Y_{22} & 0\\
J_0Y_0 & J_{20}Y_{20} & 0 & 0 & 0 & 0 & J_0^2Y_0+J_{20}^2Y_{20}-\Delta_{\rm geom}Y_g
\end{pmatrix},
\label{eq:irred-dtn-Meven-omega}
\end{equation}
with
\begin{equation}
J_0=\frac{4}{\sqrt5},
\qquad
J_{20}=\frac54.
\label{eq:irred-dtn-J0J20}
\end{equation}
This is the dynamic continuation of the static support-minus-closure matrix from
\cref{app:mouth-port}.
The same signed factorization survives at all frequencies:
\begin{equation}
M_{\rm even}(\omega)
=
R_{\rm support}(\omega)^T R_{\rm support}(\omega)
-
R_{\rm geom}(\omega)^T R_{\rm geom}(\omega),
\label{eq:irred-dtn-Meven-factorized}
\end{equation}
with the support block dressed by the channel admittances and the geometry block dressed
by \(\sqrt{\Delta_{\rm geom}Y_g(\omega)}\).

The scalar drive vector is correspondingly promoted to
\begin{equation}
J_{\rm eff}(\omega)=\bigl(J_0Y_0(\omega),J_{20}Y_{20}(\omega),0,0,0,0\bigr),
\label{eq:irred-dtn-Jeff}
\end{equation}
and the effective pure-potential coefficient becomes
\begin{equation}
S_{\rm eff}(\omega)=J_0^2Y_0(\omega)+J_{20}^2Y_{20}(\omega)-\Delta_{\rm geom}Y_g(\omega).
\label{eq:irred-dtn-Seff}
\end{equation}
At zero frequency,
\begin{equation}
S_{\rm eff}(0)=\frac54,
\label{eq:irred-dtn-Seff0}
\end{equation}
exactly as required by the solved static cross term.  Its low-frequency expansion is
\begin{equation}
S_{\rm eff}(\omega)
=
\frac54
+\omega^2
\left(
\frac{16}{5\Omega_0^2}
+\frac{25}{16\Omega_{20}^2}
-\frac{281}{80\Omega_g^2}
\right)
+O(\omega^4).
\label{eq:irred-dtn-Seff-expansion}
\end{equation}
So the first nontrivial dynamic scalar data are now isolated sharply.

The static check is immediate:
\begin{equation}
\lim_{\omega\to 0}Y_{1\perp}=\frac72,
\quad
\lim_{\omega\to 0}Y_{10}=4,
\quad
\lim_{\omega\to 0}Y_0=
\lim_{\omega\to 0}Y_{20}=
\lim_{\omega\to 0}Y_{21}=
\lim_{\omega\to 0}Y_{22}=1,
\quad
\lim_{\omega\to 0}Y_g=1,
\label{eq:irred-dtn-static-limits}
\end{equation}
so the axisymmetric pole-completed model reduces exactly to the minimal
zero-frequency operator data already solved by the conservative \(\TwoPN\) algebra.

\subsection{What is fixed and what remains open}
\label{app:irred-dtn-status}

The constructive status after this step is unusually clean.
The conservative two-body \(\TwoPN\) program has already fixed the entire
zero-frequency operator ledger:
\begin{equation}
R_{1\perp}=\frac72,
\qquad
R_{10}=4,
\qquad
R_0=R_{20}=R_{21}=R_{22}=1,
\label{eq:irred-dtn-fixed-static}
\end{equation}
\begin{equation}
J=\left(\frac{4}{\sqrt5},\frac54,0,0,0,0\right),
\qquad
\Delta_{\rm geom}=\frac{281}{80}.
\label{eq:irred-dtn-fixed-Jdelta}
\end{equation}
What remains open is not the static conservative operator, but its dynamic PDE-side
completion.  In the one-pole-per-channel model, the still-open observables are exactly the
seven pole scales
\begin{equation}
\Omega_{1\perp},\quad
\Omega_{10},\quad
\Omega_0,\quad
\Omega_{20},\quad
\Omega_{21},\quad
\Omega_{22},\quad
\Omega_g.
\label{eq:irred-dtn-open-omegas}
\end{equation}
An optional near-spherical support simplification is
\begin{equation}
\Omega_{20}=\Omega_{21}=\Omega_{22},
\label{eq:irred-dtn-nearspherical}
\end{equation}
but it is not required anywhere in the main proof.

So the remaining PDE task is now very specific.
It is no longer to guess the conservative zero-frequency throat operator from scratch.
It is to derive dynamically, from the inner-throat wall and bulk physics, the pole
structure of a channel system whose static residues, source vector, and geometry-closure
deficit are already known.

\claimclass{\Exact{} for the zero-frequency operator data
\cref{eq:irred-dtn-odd-residues,eq:irred-dtn-even-residues,eq:irred-dtn-J,eq:irred-dtn-deltageom}
and for the statement that the axisymmetric pole-completed model reduces to those data in
its \(\omega\to 0\) limit.  \Controlled{} for the isotropic 4D-ball unit test as a
PDE-side no-go benchmark and for the one-pole-per-channel axisymmetric
\(\DtN\) completion as the minimal dynamic scaffolding consistent with the solved
conservative \(\TwoPN\) cross target.  This appendix does not claim a full moving-throat
PDE solution; it identifies the exact zero-frequency target and the first open dynamic
observables that the inner-throat program must still derive.}

% ============================================================
\section{Family-1 soft-wall and wall-stress program}
\label{app:family1}

This appendix records the constructive Family-1 chain that turns the solved static even
support sector of \cref{app:mouth-port-support} into an explicit local wall law on a
flared soft wall.  The main text does not need this material for the conservative
\(\TwoPN\) proof.  Its role is explanatory: it shows that the passed
\(P_0\oplus P_2\) support/source data are not merely a solved operator table, but can be
organized as the shadow of a definite Family-1 throat program.

The chain has five steps.
First, the solved support and source profiles already live in the same low-order
Family-1 flare basis \(\{1,\mu^2,\mu^4\}\).
Second, the passed static support data admit an exact reduced variational wall-Hessian.
Third, the strict isotropic steep-wall pullback fails, so a genuinely new wall degree of
freedom is required.
Fourth, one axisymmetric tangential wall-stress profile repairs the support sector
exactly and predicts the size of the remaining monopole add-on.
Finally, that monopole add-on can itself be reinterpreted as a reduced geometry-Hessian
compressibility rather than as a disconnected auxiliary insertion.

The result is a compact constructive reading of the static Family-1 wall sector:
\begin{equation}
\boxed{
\text{local support/source wall law}
\;+
\text{one reduced global monopole breathing mode}.
}
\label{eq:family1-static-summary}
\end{equation}
This is the static wall-side companion of the minimal dynamic channel scaffolding
assembled later in \cref{app:irred-dtn}.

\subsection{The solved support/source sector already sits in the Family-1 flare basis}
\label{app:family1-basis}

As in \cref{app:mouth-port-support}, the static even support targets are
\begin{equation}
K_{00}=\frac{4}{45},
\qquad
K_{1\perp}=\frac{2}{7},
\qquad
K_{10}=\frac{1}{4},
\qquad
K_{20}=\frac{4}{9},
\qquad
K_{21}=\frac{2}{3},
\qquad
K_{22}=\frac{8}{3}.
\label{eq:family1-target-support-values}
\end{equation}
The first constructive simplification is that the full passed \(\ell=1,2\) support data
can be written as one dipole-base profile plus one curvature correction:
\begin{equation}
Z_{\ell m}
=
\big\langle z_{\rm base}(\mu)\big\rangle_{\ell m}
+
\bigl(\lambda_\ell-2\bigr)
\big\langle z_{\rm curv}(\mu)\big\rangle_{\ell m},
\qquad
\lambda_\ell\equiv \ell(\ell+1),
\qquad
\mu\equiv \cos\theta,
\label{eq:family1-zbase-zcurv}
\end{equation}
with
\begin{equation}
z_{\rm base}(\mu)=\frac{17}{56}-\frac{5}{56}\mu^2,
\qquad
z_{\rm curv}(\mu)=\frac{593}{672}-\frac{1553}{672}\mu^2+\frac78\mu^4.
\label{eq:family1-zbase-zcurv-profiles}
\end{equation}
The entire quadrupole-only correction is therefore carried by the factor
\(\lambda_\ell-2\), so the dipole is special because \(\lambda_1-2=0\).

This is exactly the basis generated by the carried Family-1 flare profile
\begin{equation}
a(z)=a_0\Bigl(1+\beta e^{-z^2/z_m^2}\Bigr).
\label{eq:family1-flare-profile}
\end{equation}
Near the mouth center, with \(a_m=a(0)\) and suitable dimensionless
\(q,r\) that are purely flare-shape coefficients rather than charges, one has
\begin{equation}
\frac{a(z)}{a_m}=1-q\mu^2+r\mu^4+O(\mu^6).
\label{eq:family1-flare-mu-expansion}
\end{equation}
So the solved support operator already sits in the same low-order polynomial basis
produced by the actual Family-1 flare geometry:
\begin{equation}
\{1,\ \mu^2,\ \mu^4\}.
\label{eq:family1-flare-basis}
\end{equation}
The source sector collapses into the same basis:
\begin{equation}
S(\mu)=\frac{11}{8}+\frac{15}{8}P_2(\mu)
=\frac{7}{16}+\frac{45}{16}\mu^2.
\label{eq:family1-source-mu-basis}
\end{equation}
So both the support operator and the scalar source sector already live in the same simple
Family-1 flare basis before any dynamical completion is attempted.

A useful equivalent operator notation is
\begin{equation}
\mathcal B_{\rm eff}
=
\partial_n
+
z_{\rm base}(\mu)
+\frac12\Big\{-\Delta_S-2,\ z_{\rm curv}(\mu)\Big\},
\label{eq:family1-generalized-robin}
\end{equation}
restricted to the support sector.  On a spherical harmonic \(Y_{\ell m}\), the
anticommutator produces precisely the factor \(\lambda_\ell-2\).

\subsection{Exact reduced variational wall-Hessian}
\label{app:family1-hessian}

The next constructive step is stronger: the passed static support data admit an exact
reduced wall-Hessian on the truncated \(\ell=0,1,2\) mouth basis.  Write
\begin{equation}
K_{00}=K_{\rm mono},
\label{eq:family1-Kmono-def}
\end{equation}
and, for \(\ell=1,2\),
\begin{equation}
K_{\ell m}
=
K_{\rm mono}
+
T_0\,\ell(\ell+1)
+
\bigl(A_0+A_1\ell(\ell+1)\bigr)\langle P_2\rangle_{\ell m}
+
\bigl(B_0+B_1\ell(\ell+1)\bigr)\langle P_2^2\rangle_{\ell m}.
\label{eq:family1-reduced-wall-hessian}
\end{equation}
The required exact moments are
\begin{align}
\langle P_2\rangle_{1\perp}&=-\frac15,
&
\langle P_2\rangle_{10}&=\frac25,
&
\langle P_2\rangle_{20}&=\frac27,
&
\langle P_2\rangle_{21}&=\frac17,
&
\langle P_2\rangle_{22}&=-\frac27,
\label{eq:family1-P2-moments}
\\
\langle P_2^2\rangle_{1\perp}&=\frac17,
&
\langle P_2^2\rangle_{10}&=\frac{11}{35},
&
\langle P_2^2\rangle_{20}&=\frac37,
&
\langle P_2^2\rangle_{21}&=\frac17,
&
\langle P_2^2\rangle_{22}&=\frac17.
\label{eq:family1-P2sq-moments}
\end{align}
Solving \cref{eq:family1-reduced-wall-hessian} against the passed support data gives
\begin{equation}
K_{\rm mono}=\frac{4}{45},
\qquad
T_0=\frac{23}{135},
\label{eq:family1-T0}
\end{equation}
\begin{equation}
A_0=\frac{9095}{3528},
\qquad
A_1=-\frac{25559}{21168},
\qquad
B_0=-\frac{109}{56},
\qquad
B_1=\frac{1765}{3024}.
\label{eq:family1-A0A1B0B1}
\end{equation}
Substituting these values back reproduces all six channel stiffnesses with zero residual.
Thus the static axisymmetric support operator is not just a fitted Robin table; it is the
Hessian of a concrete reduced variational wall law.

This reduced wall-Hessian has a direct flare interpretation.  Define
\begin{equation}
\xi\equiv \frac{a_0^2}{z_m^2}.
\label{eq:family1-xi}
\end{equation}
Expanding the Family-1 flare profile through second order gives
\begin{equation}
\frac{\delta a(\theta)}{a_0}
=
\beta\Bigl[1-\frac{\xi}{3}+\frac{\xi^2}{18}\Bigr]
+
\beta\Bigl[-\frac{2\xi}{3}+\frac{2\xi^2}{9}\Bigr]P_2
+
\beta\Bigl[\frac{2\xi^2}{9}\Bigr]P_2^2
+O(\xi^3).
\label{eq:family1-flare-P2-P2sq}
\end{equation}
So the actual frozen Family-1 flare automatically generates the exact axisymmetric basis
\(\{1,P_2,P_2^2\}\) required by the reduced wall-Hessian.

There is also a useful sharp sub-result.  If the anisotropic gradient block of the
reduced wall law is assumed linear in the Family-1 flare profile, then
\begin{equation}
\frac{B_1}{A_1}=\frac{\xi}{\xi-3}.
\label{eq:family1-linear-gradient-ratio}
\end{equation}
Using \cref{eq:family1-A0A1B0B1} gives
\begin{equation}
\xi=\frac{12355}{12638}\approx 0.9776072163,
\qquad
\frac{z_m}{a_0}\approx 1.0113880097,
\label{eq:family1-zm-over-a0}
\end{equation}
so the minimal linear-gradient reading points to a flare width of order the throat
radius.

\subsection{Strict isotropic steep-wall no-go}
\label{app:family1-strict}

The next question is whether the solved support operator could already come from the
strictest possible steep-wall interpretation of the Family-1 geometry: one isotropic
normal penetration moment, one isotropic tangential-gradient moment, and the full
geometric pullback of the flare.  Write
\begin{equation}
\sigma(\mu)=1-q\mu^2+r\mu^4,
\label{eq:family1-sigma}
\end{equation}
and let \(J(\mu)\), \(F_\theta(\mu)\), and \(F_\phi(\mu)\) denote the resulting
quadratic-flare pullback factors.  The strict boundary-layer surface energy is
\begin{equation}
E^{\rm strict}_{\rm BL}
=
\frac12\int d\Omega
\left[
A\,J(\mu)\,\Psi^2
+
B\left(
F_{\theta}(\mu)(\partial_{\theta}\Psi)^2
+
F_{\phi}(\mu)\frac{(\partial_{\phi}\Psi)^2}{\sin^2\theta}
\right)
\right].
\label{eq:family1-strict-EBL}
\end{equation}
Projecting \cref{eq:family1-strict-EBL} onto the \(\ell=0,1,2\) mouth harmonics gives
exact channel formulas polynomial in \((A,B,q,r)\).  But a multi-seed least-squares scan
shows that the best-fit residual is not small:
\begin{equation}
\sum_i \Delta_i^2 \approx 0.5536733194,
\qquad
\max_i|\Delta_i|\approx 0.4390744081.
\label{eq:family1-strict-no-go-residual}
\end{equation}
So the passed static support sector is \emph{not} just the strict isotropic steep-wall
layer pulled back through Family-1 flare geometry.

This is a useful negative result.  It means that pure isotropic penetration plus pure
geometry pullback is too small a model, and that at least one extra local wall-stress,
traction, or profile degree of freedom is genuinely required.

\subsection{Exact traction-balance completion and local wall-stress law}
\label{app:family1-traction}

The smallest exact repair is to keep the monopole wall channel separate and solve the
\(\ell=1,2\) support sector with one base pressure profile and one tangential
wall-stress / curvature-compliance profile.  On the support sector, write
\begin{equation}
K_{\ell m}
=
\big\langle z_{\rm base}\big\rangle_{\ell m}
+
\bigl(\ell(\ell+1)-2\bigr)\langle t\rangle_{\ell m},
\label{eq:family1-traction-Klm}
\end{equation}
with low-order Family-1 profiles
\begin{equation}
z_{\rm base}(\mu)=b_0+b_2\mu^2,
\qquad
t(\mu)=t_0+t_2\mu^2+t_4\mu^4.
\label{eq:family1-t-profile-ansatz}
\end{equation}
Matching the five carried support targets uniquely fixes
\begin{equation}
b_0=\frac{17}{56},
\qquad
b_2=-\frac{5}{56},
\qquad
t_0=\frac{593}{672},
\qquad
t_2=-\frac{1553}{672},
\qquad
t_4=\frac78,
\label{eq:family1-traction-coeffs}
\end{equation}
so that
\begin{equation}
z_{\rm base}(\mu)=\frac{17}{56}-\frac{5}{56}\mu^2,
\qquad
t(\mu)=\frac{593}{672}-\frac{1553}{672}\mu^2+\frac78\mu^4.
\label{eq:family1-zbase-t-final}
\end{equation}
In Legendre form,
\begin{equation}
z_{\rm base}(\mu)=\frac{23}{84}-\frac{5}{84}P_2(\mu),
\qquad
t(\mu)=\frac{211}{1008}-\frac{129}{112}P_2(\mu)+\frac{7}{18}P_2(\mu)^2.
\label{eq:family1-zbase-t-legendre}
\end{equation}
This is exactly the same Family-1 low-order basis that kept appearing in the flare and
reduced-Hessian steps.

The equivalent local quadratic form is
\begin{equation}
E_{\rm supp}[\Psi]
=
\frac12\int d\Omega\,
\Big[p(\mu)\Psi^2+t(\mu)|\nabla_S\Psi|^2\Big],
\label{eq:family1-Esupp}
\end{equation}
with
\begin{equation}
p(\mu)=z_{\rm base}(\mu)-2t(\mu)-\frac12\Delta_S t(\mu).
\label{eq:family1-p-from-t}
\end{equation}
Using
\begin{equation}
\Delta_S t(\mu)
=-\frac{35}{2}\mu^4+\frac{2729}{112}\mu^2-\frac{1553}{336},
\label{eq:family1-DeltaS-t}
\end{equation}
one finds the exact pressure profile
\begin{equation}
p(\mu)=\frac{571}{672}-\frac{5141}{672}\mu^2+7\mu^4.
\label{eq:family1-pressure-profile}
\end{equation}
Thus the solved support sector is an explicit local wall model: one base pressure
profile, one tangential wall-stress / curvature-compliance profile, and the induced
pressure profile that follows from the anticommutator structure.  All five
\(\ell=1,2\) support residuals then vanish identically.

\subsection{Minimal anisotropic PDE-side completion and universal monopole prediction}
\label{app:family1-aniso}

The traction-balance rewrite above still leaves open whether a genuine PDE-side
boundary-layer completion could generate the needed tangential wall law.  The minimal new
test is to keep the normal penetration moment isotropic, keep the same strict flare
pullback, and promote only the tangential wall moment from a constant to the smallest
axisymmetric profile
\begin{equation}
B_{\rm tan}(\mu)=B_0+B_2\mu^2.
\label{eq:family1-Btan}
\end{equation}
The boundary-layer energy becomes
\begin{equation}
E_{\rm BL}^{\rm aniso}
=
\frac12\int d\Omega\,
\Big[
A\,J(\mu)\,\Psi^2
+
B_{\rm tan}(\mu)
\Big(
F_{\theta}(\mu)(\partial_{\theta}\Psi)^2
+
F_{\phi}(\mu)\frac{(\partial_{\phi}\Psi)^2}{\sin^2\theta}
\Big)
\Big].
\label{eq:family1-aniso-EBL}
\end{equation}
Projecting onto the \(\ell=1,2\) channels gives exact formulas linear in
\((A,B_0,B_2)\) and polynomial in \((q,r)\).  The resulting system has exact real
solutions, so the strict isotropic no-go is repaired by a single new local degree of
freedom: an axisymmetric tangential wall-stress profile.

A natural Family-1-like branch is the one with \(q>0\), \(r>0\), and
\(\sigma(\mu)>0\) on \([-1,1]\).  On that branch,
\begin{equation}
A\approx -0.281923219302567,
\qquad
B_0\approx 0.648914709228733,
\qquad
B_2\approx -1.085434603876673,
\label{eq:family1-familylike-branch-AB}
\end{equation}
\begin{equation}
q\approx 2.370915717168574,
\qquad
r\approx 2.758343474052255,
\label{eq:family1-familylike-branch-qr}
\end{equation}
and the support residuals vanish at machine precision.  In Legendre form,
\begin{equation}
B_{\rm tan}(\mu)=\beta_0+\beta_2 P_2(\mu),
\qquad
\beta_0=B_0+\frac13 B_2\approx 0.287103174603175,
\qquad
\beta_2=\frac23 B_2\approx -0.723623069251115.
\label{eq:family1-beta0-beta2}
\end{equation}
So the minimal successful PDE-side completion is an axisymmetric
\(P_2\)-modulated tangential traction profile.

The most useful structural consequence is a universal monopole identity:
\begin{equation}
K_{00}
=
K_{1\perp}+\frac12K_{10}-\frac1{10}K_{20}-\frac15K_{21}-\frac15K_{22}.
\label{eq:family1-universal-monopole-identity}
\end{equation}
Once the \(\ell=1,2\) support targets are matched exactly, the raw boundary-layer
monopole is fixed automatically:
\begin{equation}
K_{00}^{\rm BL}
=
\frac27+\frac12\cdot\frac14-\frac1{10}\cdot\frac49-\frac15\cdot\frac23
-\frac15\cdot\frac83
=-\frac{757}{2520}.
\label{eq:family1-K00BL}
\end{equation}
Since the carried monopole target is \(K_{00}^{\rm target}=4/45\), the separate
monopole wall add-on is fixed to
\begin{equation}
\Delta K_{00}^{\rm mono}
=
\frac{4}{45}-\left(-\frac{757}{2520}\right)
=
\frac{109}{280}.
\label{eq:family1-monopole-deficit}
\end{equation}
So the support-sector repair does not eliminate the separate monopole channel; it predicts
its exact required size.

\subsection{Support/source locking and the canonical scalar source vector}
\label{app:family1-source-lock}

The Family-1 wall program becomes even tighter once the scalar source sector is included.
With the exact carried base profile from \cref{eq:family1-zbase-t-final}, the exact
source profile obeys
\begin{equation}
S(\mu)=10-\frac{63}{2}\,z_{\rm base}(\mu).
\label{eq:family1-source-from-zbase}
\end{equation}
So the source sector is not an independent fit.  Once the base support profile is fixed,
the source profile is fixed automatically.

Equivalently, write the axisymmetric quadratic base profile as
\begin{equation}
z_{\rm base}(\mu)=a_0+a_2P_2(\mu).
\label{eq:family1-a0a2}
\end{equation}
Then the dipole support residues reconstruct it exactly:
\begin{equation}
a_0=\frac{K_{10}+2K_{1\perp}}{3},
\qquad
a_2=\frac{5}{3}(K_{10}-K_{1\perp}).
\label{eq:family1-a0a2-from-dipole}
\end{equation}
If the scalar source is written as
\begin{equation}
S(\mu)=p_{\rm iso}+q_{\rm ax}P_2(\mu),
\label{eq:family1-piso-qax}
\end{equation}
then the same constitutive law forces
\begin{equation}
p_{\rm iso}=10-\frac{21}{2}(K_{10}+2K_{1\perp}),
\qquad
q_{\rm ax}=-\frac{105}{2}(K_{10}-K_{1\perp}).
\label{eq:family1-piso-qax-from-dipole}
\end{equation}
At the passed dipole support values
\begin{equation}
K_{1\perp}=\frac27,
\qquad
K_{10}=\frac14,
\label{eq:family1-passed-dipole-values}
\end{equation}
this gives
\begin{equation}
p_{\rm iso}=\frac{11}{8},
\qquad
q_{\rm ax}=\frac{15}{8},
\label{eq:family1-piso-qax-values}
\end{equation}
exactly.

Using the canonical normalization convention of \cref{app:mouth-port,app:irred-dtn}, the
scalar source vector is
\begin{equation}
J_0=\frac{2}{\sqrt5}\left(p_{\rm iso}+\frac{q_{\rm ax}}{3}\right),
\qquad
J_{20}=\frac23 q_{\rm ax}.
\label{eq:family1-J0J20-def}
\end{equation}
Substituting \cref{eq:family1-piso-qax-from-dipole} gives
\begin{equation}
J_0=\frac{\sqrt5}{5}\bigl(20-56K_{10}-7K_{1\perp}\bigr),
\qquad
J_{20}=35\,(K_{1\perp}-K_{10}).
\label{eq:family1-J-from-dipole}
\end{equation}
so at the passed dipole support values,
\begin{equation}
J=\left(\frac{4}{\sqrt5},\frac54\right),
\label{eq:family1-J-final}
\end{equation}
exactly.  This is the same canonical scalar source vector that appears in
\cref{app:mouth-port,app:irred-dtn-static}.  The source sector used in the solved
conservative \(\TwoPN\) cross block is therefore already fixed once the dipole support
sector is fixed.

The full static local wall operator is correspondingly
\begin{equation}
\mathcal O_{\rm loc}
=
z_{\rm base}(\mu)+\frac12\{-\Delta_S-2,\ t(\mu)\},
\label{eq:family1-Oloc}
\end{equation}
with exact channel values
\begin{equation}
K_{00}^{\rm raw}=-\frac{757}{2520},
\qquad
K_{1\perp}=\frac27,
\qquad
K_{10}=\frac14,
\qquad
K_{20}=\frac49,
\qquad
K_{21}=\frac23,
\qquad
K_{22}=\frac83.
\label{eq:family1-Oloc-channels}
\end{equation}
So the minimal full static even-wall completion is simply
\begin{equation}
\mathcal O_{\rm full}
=
\mathcal O_{\rm loc}+\frac{109}{280}\,\mathbb P_{00},
\label{eq:family1-Ofull}
\end{equation}
with \(\mathbb P_{00}\) the monopole projector.

\subsection{Monopole closure from the reduced geometry Hessian}
\label{app:family1-geom}

The remaining static question is whether the monopole add-on
\(109/280\) can itself be given a geometry-side origin.  The answer is yes: it can be
reinterpreted as a low-frequency geometry compressibility of the throat variables
\((a,L)\).

Use the cylinder-like throat volume and area
\begin{equation}
V(a,L)=\frac{4\pi}{3}a^3L,
\qquad
A(a,L)=4\pi a^2L+\frac{8\pi}{3}a^3,
\label{eq:family1-VA}
\end{equation}
and the minimal curvature-completed geometry energy
\begin{equation}
E_{\rm geom}(a,L)
=
P_{\rm vac}V(a,L)+\sigma A(a,L)+\kappa_b\frac{a^2}{L}.
\label{eq:family1-Egeom}
\end{equation}
Treat the monopole port as a pressure-like source coupling through volume work,
\begin{equation}
\delta W=-p\,\delta V,
\qquad
g=\nabla_{(a,L)}V.
\label{eq:family1-volume-work}
\end{equation}
If \(H_0\) denotes the Hessian of \(E_{\rm geom}\) at the reference point
\((a_0,L_0)=(a_0,\Lambda a_0)\), then integrating out \((\delta a,\delta L)\) gives the
exact reduced monopole coefficient
\begin{equation}
\Delta K_{00}^{\rm geom}=\frac{g^T H_0^{-1}g}{V_0^2}.
\label{eq:family1-DeltaK00geom}
\end{equation}
So the old global monopole projector is reinterpreted as a genuine low-frequency geometry
compressibility.

Write
\begin{equation}
\sigma=\frac{\Sigma}{a_0^3},
\qquad
P_{\rm vac}=\rho\,\frac{\Sigma}{a_0^4},
\qquad
\kappa_b=\beta\,\frac{\Sigma}{a_0}.
\label{eq:family1-Sigma-rho-beta}
\end{equation}
For the baseline model with no curvature completion \((\beta=0)\), the reduced Hessian is
indefinite, so the literal \(PV+\sigma A\) geometry energy cannot by itself close the
monopole channel.  With the curvature completion, positivity requires
\begin{equation}
\beta>\beta_{\rm stab}(\Lambda,\rho)
\equiv
\frac{\pi\Lambda^3(\rho+2)^2}{2\Lambda\rho+3\Lambda+2},
\label{eq:family1-beta-stab}
\end{equation}
and positivity of the monopole coefficient also requires
\begin{equation}
\beta>\beta_{\Delta}(\Lambda,\rho)
\equiv
\frac{\pi\Lambda(2\Lambda\rho+5\Lambda-2)}{4}.
\label{eq:family1-beta-delta}
\end{equation}
Thus the geometry-side closure exists whenever
\begin{equation}
\beta>\max\bigl(\beta_{\rm stab},\beta_{\Delta}\bigr).
\label{eq:family1-beta-condition}
\end{equation}

On the EM branch,
\begin{equation}
\Lambda_{\rm EM}=\frac{\sqrt2\pi}{x_{01}}\approx 1.847486577120128,
\qquad
x_{01}=2.40482555769577\ldots,
\label{eq:family1-LambdaEM}
\end{equation}
a simple positive worked point is
\begin{equation}
\rho=\frac{1}{10},
\qquad
\beta=12,
\label{eq:family1-worked-rho-beta}
\end{equation}
which lies safely inside the passive/positive-support region.  Matching the exact target
\(\Delta K_{00}^{\rm geom}=109/280\) then fixes the overall geometry stiffness scale to
\begin{equation}
\Sigma_*\approx 0.2076143291835488854.
\label{eq:family1-Sigma-star}
\end{equation}
At this point,
\begin{equation}
\Delta K_{00}^{\rm geom}=\frac{109}{280}
\label{eq:family1-DeltaK00-match}
\end{equation}
at machine precision.

The resulting two-degree-of-freedom geometry sector is strongly dominated by one mode.
At the worked point above, the mode-resolved contributions to
\(\Delta K_{00}^{\rm geom}\) are approximately
\begin{equation}
(0.00878310,\ 0.38050262),
\label{eq:family1-mode-split}
\end{equation}
which sum to \(109/280\), with the larger mode carrying
\begin{equation}
97.743791\% 
\label{eq:family1-dominant-mode-percent}
\end{equation}
of the total static monopole response.  This explains why the earlier single global
monopole auxiliary was already the right effective description: it is a controlled
reduction of the full \((a,L)\) geometry sector, not a disconnected ad hoc insertion.

The cleanest static closure statement is therefore
\begin{equation}
\boxed{
\text{local Family-1 support/source wall law}
\;+
\text{dominant-coupling reduction of the reduced geometry Hessian}.
}
\label{eq:family1-final-static-closure}
\end{equation}
This closes the static Family-1 wall-stress program.  The subsequent dynamic tasks -- bulk
inertia, sidewall and endcap soft-wall corrections, and the final coupled low-frequency
response module -- belong to the dynamic \(\DtN\) scaffolding discussed in
\cref{app:irred-dtn}.

% ============================================================
\section{Geometry-breathing and dynamic monopole closures}
\label{app:geom-breathing}

This appendix records the scalar channel that remained after the static support/source
program of \cref{app:family1}.  The logic is now fairly sharp.
The local Family-1 wall law closes the full static \(\ell=1,2\) support/source sector,
but it leaves one scalar remainder in the monopole channel.
That remainder is not another arbitrary support coefficient.
It is the low-frequency shadow of a genuine reduced geometry response in the throat
variables \((a,L)\), and once a minimal dynamic reduction is written down it becomes an
exact positive two-pole breathing susceptibility.

A notational warning is worth stating first.
Two different scalar objects appear in the appendix program.
One is the pure-potential geometry-closure deficit
\(\Delta_{\mathrm{geom}}=281/80\) in the support-minus-closure EFT of
\cref{app:mouth-port,app:irred-dtn}; that quantity belongs to the
comparable-mass even cross operator.
The other is the monopole wall add-on
\begin{equation}
\Delta K_{00}=\frac{109}{280},
\label{eq:geom-breathing-deltaK00}
\end{equation}
which is the missing scalar support needed to lift the local static wall law to the full
passed monopole coefficient.
The present appendix concerns \emph{this second object}.
It is the breathing/compressibility closure of the scalar monopole wall channel,
not the pure-\(U\) closure subtraction of the even cross EFT.

The appendix has four layers.
First, we isolate the scalar static remainder left by the local wall law.
Second, we show that the missing coefficient is generated exactly by a reduced geometry
Hessian in \((a,L)\).
Third, we derive the exact dynamic monopole susceptibility of that reduced system and its
controlled one-pole low-frequency reduction.
Fourth, we fold in the Family-1 bulk, sidewall, and endcap inertia completions to obtain
an explicit wall-completed breathing branch.

\subsection{The scalar remainder left by the local static wall law}
\label{app:geom-breathing-static-gap}

The static support/source program of \cref{app:family1} compresses the full even wall
sector into
\begin{equation}
\text{[exact local support/source constitutive law]}
\;+
\Delta K_{00}\,\mathbb P_{00},
\label{eq:geom-breathing-local-plus-mono}
\end{equation}
where \(\mathbb P_{00}\) is the monopole projector.
The local part reproduces the non-monopole support and source data exactly, and the only
remaining scalar gap is
\begin{equation}
\Delta K_{00}=\frac{109}{280}.
\label{eq:geom-breathing-109-280}
\end{equation}
Equivalently, the carried local wall baseline in the monopole channel is
\begin{equation}
K_{00}^{\rm loc}=-\frac{757}{2520},
\label{eq:geom-breathing-K00loc}
\end{equation}
so the full passed monopole coefficient
\begin{equation}
K_{00}=\frac{4}{45}
\label{eq:geom-breathing-K00target}
\end{equation}
is recovered as
\begin{equation}
K_{00}=K_{00}^{\rm loc}+\Delta K_{00}
=-\frac{757}{2520}+\frac{109}{280}
=\frac{4}{45}.
\label{eq:geom-breathing-K00-split}
\end{equation}
Thus the static Family-1 wall program leaves exactly one unresolved scalar breathing
coefficient and no unresolved \(\ell=1,2\) support coefficient.

This scalar remainder is what the older notes had treated as a ``global monopole
auxiliary.''
The point of the present appendix is that the auxiliary is no longer merely a bookkeeping
insertion.
It can be derived as a reduced compressibility of the throat geometry itself.

\subsection{Static geometry-Hessian closure in \texorpdfstring{$(a,L)$}{(a,L)}}
\label{app:geom-breathing-static-hessian}

Use the cylinder-like throat geometry
\begin{equation}
V(a,L)=\frac{4\pi}{3}a^3L,
\qquad
A(a,L)=4\pi a^2L+\frac{8\pi}{3}a^3,
\label{eq:geom-breathing-VA}
\end{equation}
and the minimal curvature-completed geometry energy
\begin{equation}
E_{\rm geom}(a,L)
=
P_{\rm vac}V(a,L)+\sigma A(a,L)+\kappa_b\frac{a^2}{L}.
\label{eq:geom-breathing-Egeom}
\end{equation}
The curvature term is the smallest explicit completion that repairs the known failure of
literal \(P_{\rm vac}V+\sigma A\) geometry in the monopole channel.
Write the reference configuration as
\begin{equation}
(a_0,L_0)=(a_0,\Lambda a_0),
\qquad
V_0=V(a_0,L_0)=\frac{4\pi}{3}\Lambda a_0^4,
\label{eq:geom-breathing-reference}
\end{equation}
and let \(H_0\) be the Hessian of \(E_{\rm geom}\) at that point.
The uniform monopole port couples through volume work,
\begin{equation}
\delta W=-p\,\delta V,
\qquad
g\equiv \nabla_{(a,L)}V\big|_{(a_0,L_0)}.
\label{eq:geom-breathing-volume-work}
\end{equation}
Integrating out the reduced geometry variables then gives the exact static compressibility
coefficient
\begin{equation}
\boxed{
\Delta K_{00}^{\rm geom}
=
\frac{g^T H_0^{-1}g}{V_0^2}.
}
\label{eq:geom-breathing-static-compressibility}
\end{equation}
So the old monopole auxiliary is reinterpreted as a genuine reduced geometry-side
compressibility.

It is convenient to introduce the scaled parameters
\begin{equation}
\sigma=\frac{\Sigma}{a_0^3},
\qquad
P_{\rm vac}=\rho\,\frac{\Sigma}{a_0^4},
\qquad
\kappa_b=\beta\,\frac{\Sigma}{a_0},
\label{eq:geom-breathing-scaled-params}
\end{equation}
for which the exact reduced coefficient becomes
\begin{equation}
\Delta K_{00}^{\rm geom}
=
\frac{2\pi\Lambda^2\rho+5\pi\Lambda^2-2\pi\Lambda-4\beta}
{2\pi\Sigma\left(\pi\Lambda^3\rho^2+4\pi\Lambda^3\rho+4\pi\Lambda^3-2\Lambda\beta\rho-3\Lambda\beta-2\beta\right)}.
\label{eq:geom-breathing-exact-static-ratio}
\end{equation}
If the curvature completion is removed, \(\beta=0\), then the baseline Hessian is
\begin{equation}
H_{\rm base}
=
\begin{pmatrix}
8\pi(\Lambda\rho+\Lambda+2) & 4\pi(\rho+2)\\[3pt]
4\pi(\rho+2) & 0
\end{pmatrix},
\label{eq:geom-breathing-Hbase}
\end{equation}
with
\begin{equation}
\det H_{\rm base}=-16\pi^2(\rho+2)^2<0.
\label{eq:geom-breathing-Hbase-det}
\end{equation}
So literal \(P_{\rm vac}V+\sigma A\) geometry is not a passive \(2\)-DOF Hessian and
cannot by itself close the monopole channel.
This no-go is the static reason a curvature or soft-wall completion is necessary.

For the curvature-completed model, positivity and positive monopole support require
\begin{equation}
\beta>0,
\qquad
\beta>\beta_{\rm stab}(\Lambda,\rho)
\equiv
\frac{\pi\Lambda^3(\rho+2)^2}{2\Lambda\rho+3\Lambda+2},
\label{eq:geom-breathing-beta-stab}
\end{equation}

and
\begin{equation}
\beta>\beta_{\Delta}(\Lambda,\rho)
\equiv
\frac{\pi\Lambda(2\Lambda\rho+5\Lambda-2)}{4}.
\label{eq:geom-breathing-beta-delta}
\end{equation}
A representative EM-branch worked point is
\begin{equation}
\Lambda_{\rm EM}=\frac{\sqrt2\pi}{x_{01}}\approx 1.847486577120128,
\qquad
\rho=\frac{1}{10},
\qquad
\beta=12,
\label{eq:geom-breathing-EM-point}
\end{equation}
with
\begin{equation}
\beta_{\rm stab}(\Lambda_{\rm EM},\rho)\approx 11.0420171,
\qquad
\beta_{\Delta}(\Lambda_{\rm EM},\rho)\approx 11.0377513,
\label{eq:geom-breathing-EM-thresholds}
\end{equation}
so the point lies inside the passive positive-support region.
Matching the exact target
\begin{equation}
\Delta K_{00}^{\rm geom}=\frac{109}{280}
\label{eq:geom-breathing-static-match}
\end{equation}
then fixes the overall geometry stiffness scale to
\begin{equation}
\Sigma_*\approx 0.2076143291835488854.
\label{eq:geom-breathing-Sigma-star}
\end{equation}
At this point the normalized static coupling vector
\begin{equation}
\bar g\equiv \frac{\nabla V_0}{V_0}=
\left(3,\frac{1}{\Lambda_{\rm EM}}\right)
\label{eq:geom-breathing-gbar}
\end{equation}
feeds dominantly into one Hessian eigenmode, which is why the historical single breathing
auxiliary was already a good static reduction.

\subsection{Exact dynamic monopole susceptibility and the two-pole breathing branch}
\label{app:geom-breathing-dynamic}

To dynamize the static compressibility, use the dimensionless geometry variables
\begin{equation}
\Qgeom=\left(\frac{\delta a}{a_0},\frac{\delta L}{a_0}\right)
\label{eq:geom-breathing-Qgeom-vars}
\end{equation}
and the minimal affine displacement field inside the cylinder-like throat,
\begin{equation}
\boldsymbol\xi_\perp=\frac{\delta a}{a_0}\,\mathbf r,
\qquad
\xi_w=\frac{\delta L}{L_0}\,w.
\label{eq:geom-breathing-affine-field}
\end{equation}
For a uniform bulk inertia density \(\rho_{\rm eff}\), the reduced kinetic energy takes
exactly the form
\begin{equation}
T^{(2)}
=
\frac12\rho_{\rm eff}V_0a_0^2\,\dot{\Qgeom}^{T}\widehat M\,\dot{\Qgeom},
\qquad
\widehat M=
\begin{pmatrix}
\frac35 & 0\\[3pt]
0 & \frac1{12}
\end{pmatrix}.
\label{eq:geom-breathing-Mhat-affine}
\end{equation}
The corresponding quadratic energy is
\begin{equation}
E^{(2)}=\frac12\Sigma\,\Qgeom^{T}\widehat H\,\Qgeom,
\qquad
\widehat H\equiv \frac{a_0^2}{\Sigma}H_0.
\label{eq:geom-breathing-E2-Hhat}
\end{equation}
Define the scaled frequency variable
\begin{equation}
s\equiv \omega^2\frac{\rho_{\rm eff}V_0a_0^2}{\Sigma}.
\label{eq:geom-breathing-svar}
\end{equation}
Then the exact reduced geometry response is
\begin{equation}
\boxed{
Y_{\rm geom}(s)
=
\frac{1}{\Sigma}\,
\bar g^T\bigl(\widehat H-s\widehat M\bigr)^{-1}\bar g.
}
\label{eq:geom-breathing-Ygeom}
\end{equation}
The raw monopole wall channel is therefore
\begin{equation}
\boxed{
K_{00}^{\rm raw}(s)
=
-\frac{757}{2520}+Y_{\rm geom}(s).
}
\label{eq:geom-breathing-K00raws}
\end{equation}
At zero frequency,
\cref{eq:geom-breathing-K00raws} reduces to the static closure already discussed above.

Because \(\widehat H\) and \(\widehat M\) are positive on the worked branch,
\(Y_{\rm geom}(s)\) is a two-pole Stieltjes function.
Equivalently, if
\begin{equation}
\widehat H v_i=\lambda_i\widehat M v_i,
\qquad i\in\{-,+\},
\label{eq:geom-breathing-generalized-eigs}
\end{equation}
then
\begin{equation}
Y_{\rm geom}(s)
=
\frac{R_-}{1-s/\lambda_-}
+
\frac{R_+}{1-s/\lambda_+},
\qquad
R_\pm>0,
\qquad
R_-+R_+=\frac{109}{280}.
\label{eq:geom-breathing-two-pole}
\end{equation}
So the monopole wall completion is not a generic ad hoc pole fit.
It is the exact dynamic response of a reduced \((a,L)\) geometry sector.

At the EM worked point of
\cref{eq:geom-breathing-EM-point,eq:geom-breathing-Sigma-star}, one finds
\begin{equation}
\widehat H\approx
\begin{pmatrix}
114.33174685 & 19.35786733\\[3pt]
19.35786733 & 3.80598758
\end{pmatrix},
\label{eq:geom-breathing-Hhat-EM}
\end{equation}
with generalized pole parameters
\begin{equation}
\lambda_-\approx 5.23115613,
\qquad
\lambda_+\approx 230.99360624,
\label{eq:geom-breathing-lambdas-affine}
\end{equation}
and positive residues
\begin{equation}
R_-\approx 0.00327376153,
\qquad
R_+\approx 0.38601195275,
\qquad
R_-+R_+=\frac{109}{280}.
\label{eq:geom-breathing-residues-affine}
\end{equation}
The physical pole frequencies are
\begin{equation}
\Omega_\pm^2
=
\frac{\Sigma_*}{\rho_{\rm eff}V_0a_0^2}\,\lambda_\pm.
\label{eq:geom-breathing-physical-poles-affine}
\end{equation}
Thus the exact geometry breathing response is an explicit positive two-pole dynamic
closure whose static limit reproduces the required monopole coefficient.

\subsection{Controlled one-pole reduction and the old breathing auxiliary}
\label{app:geom-breathing-pade}

The historical single breathing auxiliary is recovered by a low-frequency Pad\'e reduction
of the exact two-pole function.
Expanding
\cref{eq:geom-breathing-Ygeom} at small \(s\) gives
\begin{equation}
Y_{\rm geom}(s)=\Delta_0+\Delta_2 s+O(s^2),
\label{eq:geom-breathing-small-s}
\end{equation}
with
\begin{equation}
\Delta_0=\frac{1}{\Sigma}\bar g^T\widehat H^{-1}\bar g,
\qquad
\Delta_2=\frac{1}{\Sigma}\bar g^T\widehat H^{-1}\widehat M\widehat H^{-1}\bar g.
\label{eq:geom-breathing-Delta0-Delta2}
\end{equation}
The associated \([1/1]\) Pad\'e reduction is
\begin{equation}
\boxed{
Y_{\rm geom}(s)\approx
\frac{\Delta_0}{1-s/\lambda_{\rm eff}},
\qquad
\lambda_{\rm eff}=\frac{\Delta_0}{\Delta_2}.
}
\label{eq:geom-breathing-pade}
\end{equation}
Since \(\Delta_0=109/280\), this may be written equivalently as
\begin{equation}
Y_{\rm mono}^{\rm eff}(s)
=
\frac{109/280}{1-s/\lambda_{\rm eff}}.
\label{eq:geom-breathing-ymono-eff}
\end{equation}
At the EM affine worked point,
\begin{equation}
\lambda_{\rm eff}\approx 169.48205088.
\label{eq:geom-breathing-lambda-eff-affine}
\end{equation}
On the natural low-frequency band
\begin{equation}
0\le s\le 0.1\lambda_-,
\label{eq:geom-breathing-natural-band}
\end{equation}
the maximum relative error of the one-pole reduction against the exact two-pole response
is only
\begin{equation}
\max|\delta Y/Y|\approx 8.87\times 10^{-5}.
\label{eq:geom-breathing-pade-error-affine}
\end{equation}
So the old global breathing auxiliary is a controlled low-frequency reduction of the
exact reduced geometry dynamics, not a separate phenomenological invention.

\subsection{Family-1 bulk inertia and the \texorpdfstring{$n=5$}{n=5} Thomas--Fermi branch}
\label{app:geom-breathing-TF}

The affine model above leaves the overall inertia scale \(\rho_{\rm eff}\) open.
The next step is to derive that scale from the Family-1 parent PDE on the filled-to-endcap
Thomas--Fermi branch.
Use the harmonic interior trap
\begin{equation}
V_{\rm in}(w)=\frac12 m_\psi\Omega_{\rm in}^2 w^2,
\label{eq:geom-breathing-Vin}
\end{equation}
with frozen barotrope
\begin{equation}
P(\rho)=K_{\rm EOS}\rho^n,
\qquad
h(\rho)=\frac{nK_{\rm EOS}}{n-1}\rho^{n-1}.
\label{eq:geom-breathing-eos}
\end{equation}
In the Thomas--Fermi/hydrostatic limit,
\begin{equation}
h(\rho)+\frac12m_\psi\Omega_{\rm in}^2 w^2=\mu_{\rm TF}.
\label{eq:geom-breathing-TF-balance}
\end{equation}
On the filled-to-endcap branch,
\begin{equation}
\mu_{\rm TF}=\frac12m_\psi\Omega_{\rm in}^2\left(\frac{L_0}{2}\right)^2,
\label{eq:geom-breathing-muTF}
\end{equation}
so the exact axial density profile is
\begin{equation}
\rho_0(w)=\rho_c\left(1-\frac{4w^2}{L_0^2}\right)^{\!1/(n-1)},
\qquad
|w|\le \frac{L_0}{2},
\label{eq:geom-breathing-rho0w}
\end{equation}
with central density
\begin{equation}
\rho_c=
\left[
\frac{(n-1)m_\psi\Omega_{\rm in}^2L_0^2}{8nK_{\rm EOS}}
\right]^{\!1/(n-1)}.
\label{eq:geom-breathing-rhoc}
\end{equation}
Define
\begin{equation}
c_0(n)=\frac12\int_{-1}^{1}(1-u^2)^{1/(n-1)}\,du.
\label{eq:geom-breathing-c0n-def}
\end{equation}
Then the effective bulk inertia density entering the monopole breathing channel is
\begin{equation}
\boxed{
\rho_{\rm eff}^{\rm TF}(n)=c_0(n)\rho_c,
}
\label{eq:geom-breathing-rhoeffTFn}
\end{equation}
with exact closed form
\begin{equation}
\boxed{
c_0(n)=
\frac{\sqrt\pi\,\Gamma\!\left(1+\frac{1}{n-1}\right)}
{2\,\Gamma\!\left(\frac32+\frac{1}{n-1}\right)}.
}
\label{eq:geom-breathing-c0n}
\end{equation}
For the frozen \(n=5\) branch,
\begin{equation}
\rho_{\rm eff}^{\rm TF}(5)=
\frac{\sqrt\pi\,\Gamma(1/4)}{6\Gamma(3/4)}
\left(
\frac{m_\psi\Omega_{\rm in}^2L_0^2}{10K_{\rm EOS}}
\right)^{1/4},
\label{eq:geom-breathing-rhoeffTF5}
\end{equation}
so the previously free bulk inertia scale is now tied directly to the same Family-1
microphysics that defines the parent throat branch.

The same reduction also renormalizes the axial inertia moment.
Define
\begin{equation}
c_2(n)=\frac12\int_{-1}^{1}u^2(1-u^2)^{1/(n-1)}\,du,
\label{eq:geom-breathing-c2n-def}
\end{equation}
for which
\begin{equation}
\frac{c_2(n)}{c_0(n)}=\frac{n-1}{3n-1}.
\label{eq:geom-breathing-c2-over-c0}
\end{equation}
Therefore
\begin{equation}
\widehat M_{LL}^{\rm TF}(n)=\frac{c_2}{4c_0}=\frac{n-1}{4(3n-1)}.
\label{eq:geom-breathing-MLL-TFn}
\end{equation}
At \(n=5\),
\begin{equation}
\boxed{
\widehat M_{\rm TF}(5)=
\begin{pmatrix}
\frac35 & 0\\[3pt]
0 & \frac1{14}
\end{pmatrix}.
}
\label{eq:geom-breathing-Mhat-TF5}
\end{equation}
So the Family-1 \(n=5\) branch renormalizes the axial breathing moment from the
uniform-slice value \(1/12\) to the EOS-sensitive value \(1/14\).

With this replacement, the exact reduced response becomes
\begin{equation}
Y_{\rm geom}^{\rm TF}(s)=
\frac{1}{\Sigma}\,
\bar g^T\bigl(\widehat H-s\widehat M_{\rm TF}(5)\bigr)^{-1}\bar g,
\qquad
s=\omega^2\frac{\rho_{\rm eff}^{\rm TF}(5)V_0a_0^2}{\Sigma}.
\label{eq:geom-breathing-Ygeom-TF}
\end{equation}
At the same EM worked point, one finds
\begin{equation}
\lambda_-\approx 5.92556258,
\qquad
\lambda_+\approx 237.91117494,
\label{eq:geom-breathing-lambdas-TF}
\end{equation}
with residues
\begin{equation}
R_-\approx 0.00262800,
\qquad
R_+\approx 0.38665771,
\qquad
R_-+R_+=\frac{109}{280},
\label{eq:geom-breathing-residues-TF}
\end{equation}
and effective Pad\'e pole
\begin{equation}
\lambda_{\rm eff}^{\rm TF}\approx 188.17695898.
\label{eq:geom-breathing-lambdaeff-TF}
\end{equation}
The one-pole reduction remains extremely accurate,
\begin{equation}
\max\text{ rel.err.}\approx 7.10\times 10^{-5}
\qquad
\text{for } 0\le s\le 0.1\lambda_-.
\label{eq:geom-breathing-TF-error}
\end{equation}
So the bulk part of the monopole dynamics is no longer phenomenological.

\subsection{Wall-completed dynamic monopole closure}
\label{app:geom-breathing-wall-completed}

Beyond the Thomas--Fermi bulk, the remaining dynamic ingredients are genuine sidewall and
endcap inertia corrections.
The static closure \(\Delta K_{00}=109/280\) is unaffected; only the inertia sector is
renormalized.
That is why the monopole wall response remains of the same structural form,
\begin{equation}
K_{00}(s)
=
-\frac{757}{2520}
+
\frac{R_-}{1-s/\lambda_-}
+
\frac{R_+}{1-s/\lambda_+},
\qquad
R_-+R_+=\frac{109}{280},
\label{eq:geom-breathing-K00-wall-completed-general}
\end{equation}
throughout the dynamic completion chain.

The separated-order Family-1 sidewall and endcap reductions of
\cref{app:family1} already gave a near-final composite wall completion.
But the most useful summary datum is the direct coupled full-profile branch, which feeds
both effects into the same reduced breathing operator.
On the balanced thin-layer reference branch one finds
\begin{equation}
R_{\rm mass}^{\rm full}\approx 0.886313972989725,
\qquad
\widehat M_{aa}^{\rm full}\approx 0.563114968953987,
\qquad
\widehat M_{LL}^{\rm full}\approx 0.065829228119349.
\label{eq:geom-breathing-full-profile-inertia}
\end{equation}
Using these values in the same geometry Hessian at the EM worked point gives
\begin{equation}
\lambda_-\approx 6.405572392138922,
\qquad
\lambda_+\approx 254.444968136936126,
\label{eq:geom-breathing-lambdas-full}
\end{equation}
with positive residues
\begin{equation}
R_-\approx 0.002552474771738,
\qquad
R_+\approx 0.386733239513976,
\qquad
R_-+R_+=\frac{109}{280}.
\label{eq:geom-breathing-residues-full}
\end{equation}
Therefore
\begin{equation}
K_{00}(0)=\frac{4}{45}
\label{eq:geom-breathing-K00zero-full}
\end{equation}
exactly, while the physical pole ratios relative to the sharp-wall Thomas--Fermi baseline
are shifted to
\begin{equation}
\frac{\Omega_-^2}{\Omega_{-,\rm TF}^2}\approx 1.219665554412172,
\qquad
\frac{\Omega_+^2}{\Omega_{+,\rm TF}^2}\approx 1.206678094538845.
\label{eq:geom-breathing-pole-ratios-full}
\end{equation}
The corresponding one-pole Pad\'e reduction is still excellent:
\begin{equation}
\lambda_{\rm eff}\approx 202.923516367519028,
\qquad
\max\text{ rel.err.}\approx 6.89\times 10^{-5}
\qquad
\text{on }0\le s\le 0.1\lambda_-.
\label{eq:geom-breathing-lambdaeff-full}
\end{equation}
So even after the wall completion is made explicit, the historical single breathing mode
remains a highly accurate low-frequency reduction of the exact two-pole geometry channel.

\subsection{Interpretation and status}
\label{app:geom-breathing-status}

The constructive status of the monopole channel is now much cleaner than it was at the
start of the \(\TwoPN\) program.
There are really three nested statements.

First, at strictly static level, the missing monopole coefficient is not an arbitrary
projector but the reduced compressibility of a curvature-completed throat geometry.
That is the content of
\cref{eq:geom-breathing-static-compressibility,eq:geom-breathing-static-match}.

Second, once the geometry variables \((a,L)\) are given a reduced inertia metric, the
same scalar channel becomes an exact positive two-pole breathing susceptibility.
That is the content of
\cref{eq:geom-breathing-two-pole,eq:geom-breathing-K00raws}.

Third, the Family-1 parent PDE and wall program now determine not only the static closure
but most of the dynamic normalization as well: the bulk inertia scale
\(\rho_{\rm eff}^{\rm TF}(5)\), the EOS-sensitive axial moment \(1/14\), and the first
explicit sidewall/endcap/full-profile corrections.
So the dynamic monopole branch is no longer a placeholder pole.
It is a concrete low-frequency channel with a known exact static residue, a known reduced
geometry origin, and a highly accurate one-pole reduction.

What still remains open is narrower than before.
This appendix does \emph{not} derive the non-monopole pole scales of the odd and even
\(\ell=1,2\) channels, and it does not claim a fully analytic microscopic derivation of
the local wall profiles beyond the reduced Family-1 completion.
But as far as the scalar breathing branch is concerned, the constructive program has
already reached a near-final form:
\begin{equation}
\boxed{
\text{exact static monopole closure}
\;+
\text{exact reduced two-pole breathing response}
\;+
\text{explicit Family-1 bulk and wall inertia completion}.
}
\label{eq:geom-breathing-final-summary}
\end{equation}
That is the sense in which the geometry-breathing and dynamic monopole closures are now
substantially in hand.

\claimclass{\Controlled{} for the reduced geometry-side model, affine breathing ansatz,
and Family-1 wall reductions.
\Exact{} for the static compressibility formula once the reduced model is declared, for
the exact two-pole response of the resulting \((a,L)\) system, and for the
Thomas--Fermi inertia integrals.
\Protocol{} for the specific thin-wall and coupled full-profile wall completion used to
package the final breathing branch.}

% ============================================================
\section{Final coupled Family-1 response module}
\label{app:family1-final}

This appendix is the constructive end-point of the throat-side program developed in
\cref{app:mouth-port,app:irred-dtn,app:family1,app:geom-breathing}.
Its role is not to add a new ingredient to the conservative \(\TwoPN\) proof, but to
show that the previously separated constructive pieces can now be packaged into one
coherent low-frequency Family-1 response module.
There are two parts to that statement.

First, the exact zero-frequency odd and even channel data already derived in the earlier
appendices remain unchanged: the odd dipole wake, the canonical
\(P_0\oplus P_2\) even support/source sector, and the pure-potential geometry-closure
channel still reproduce the solved added comparable-mass conservative \(\TwoPN\) cross
block exactly.
Second, the monopole breathing channel can now be fed not merely with a separated
sidewall-plus-endcap inertia estimate, but with a direct \emph{coupled} Family-1
full-profile wall completion.
The resulting monopole susceptibility still closes the same static scalar target,
but now does so with the best current dynamic normalization of the Family-1 program.

This is therefore the right place to summarize what the constructive appendix program has
actually achieved.
At zero frequency, the conservative added \(\TwoPN\) cross block is fixed exactly.
In the monopole wall channel, the scalar breathing response is fixed to an exact positive
low-frequency two-pole form whose inertia data now come from the direct coupled
Family-1 wall profile.
What remains open after this point is narrower: the genuinely non-monopole dynamic pole
scales and a more microscopic derivation of the wall profiles themselves.

\subsection{Static wall law and source profile remain unchanged}
\label{app:family1-final-static-law}

The direct coupled wall completion does \emph{not} change the static support/source law
already derived in \cref{app:family1}.
The carried local wall profiles remain
\begin{equation}
 z_{\rm base}(\mu)=\frac{17}{56}-\frac{5}{56}\mu^2,
 \qquad
 t(\mu)=\frac{593}{672}-\frac{1553}{672}\mu^2+\frac78\mu^4,
 \label{eq:family1-final-zbase-t}
\end{equation}
with exact scalar source profile
\begin{equation}
S(\mu)=10-\frac{63}{2}\,z_{\rm base}(\mu)
=\frac{7}{16}+\frac{45}{16}\mu^2.
\label{eq:family1-final-source-profile}
\end{equation}
Equivalently, the local static even-wall law is still the same operator already recorded
in \cref{app:family1,app:geom-breathing}:
\begin{equation}
\mathcal O_{\rm loc}
=
 z_{\rm base}(\mu)
+\frac12\Bigl\{-\Delta_S-2,\ t(\mu)\Bigr\}.
\label{eq:family1-final-Oloc}
\end{equation}
So the coupled full-profile step is purely a \emph{dynamic} refinement of the breathing
sector.
It does not alter the previously solved static support/source data.

\subsection{Exact zero-frequency odd and even channel packaging}
\label{app:family1-final-zero-freq}

The zero-frequency constructive module also remains unchanged.
Using the carried odd dipole wake from the \(\OnePN\) program and the canonical
\(P_0\oplus P_2\) even ports of \cref{app:mouth-port,app:irred-dtn}, the added
conservative comparable-mass \(\TwoPN\) cross block is still reproduced exactly by
\begin{equation}
L^{\rm add}_{\rm odd}
=
\frac12(v_A^2+v_B^2)L^{\rm wake}_{1\mathrm{PN}}
-\frac{15}{4}(U_A+U_B)(\mathbf v_A\!\cdot\!\mathbf v_B),
\label{eq:family1-final-Lodd}
\end{equation}
combined with the even block
\begin{equation}
L^{\rm add}_{\rm even}
=
\Pi_A\!\cdot\!\Pi_B
+U_A\,J\!\cdot\!\Pi_B
+U_B\,J\!\cdot\!\Pi_A
+\bigl(J\!\cdot\!J-\Delta_{\rm geom}\bigr)U_AU_B,
\label{eq:family1-final-Leven}
\end{equation}
where the six-component real even-port basis is the one already defined in
\cref{app:mouth-port,app:irred-dtn} and the source vector is
\begin{equation}
J=\left(\frac{4}{\sqrt5},\frac54,0,0,0,0\right),
\qquad
\Delta_{\rm geom}=\frac{281}{80}.
\label{eq:family1-final-J-deltageom}
\end{equation}
Again, \(J\) denotes the appendix-side throat-response source-weight vector, not a
Maxwell current.
Thus
\begin{equation}
J\!\cdot\!J
=\left(\frac{4}{\sqrt5}\right)^2+\left(\frac54\right)^2
=\frac{381}{80},
\label{eq:family1-final-JdotJ}
\end{equation}
so the support part alone would overproduce the static coefficient, and the direct
geometry-closure subtraction \(\Delta_{\rm geom}=281/80\) lowers it to the exact solved
pure-static added-cross value
\begin{equation}
\frac{381}{80}-\frac{281}{80}=\frac54.
\label{eq:family1-final-static-cross}
\end{equation}
In other words, the constructive residual against the exact solved added comparable-mass
\(\TwoPN\) target still vanishes identically at zero frequency.
Nothing in the coupled Family-1 wall completion changes that algebraic closure.

\subsection{Direct coupled full-profile Family-1 wall completion}
\label{app:family1-final-coupled-wall}

What \emph{does} change at the present stage is the monopole breathing normalization.
The earlier wall-completed breathing branch could be assembled in separated sidewall and
endcap order.
The present step goes beyond that approximation and evaluates the full coupled Family-1
Thomas--Fermi wall profile directly in the throat variables \((x,y)\).

On the balanced thin-layer-consistent reference branch,
\begin{equation}
\alpha_r=10,
\qquad
\varepsilon_r=0.05,
\qquad
p_r=2,
\label{eq:family1-final-radial-branch}
\end{equation}
\begin{equation}
\chi_{\rm cap}=4,
\qquad
\varepsilon_z=0.05,
\qquad
\alpha_{\rm cap}=4\varepsilon_z\chi_{\rm cap}=0.8,
\qquad
p_z=2,
\label{eq:family1-final-cap-branch}
\end{equation}
use the full coupled profile
\begin{equation}
\widetilde\rho(x,y)
=
\Biggl[
1-y^2
-\alpha_r\,S\!\left(\frac{x-1}{\varepsilon_r}\right)^{p_r}
-\alpha_{\rm cap}\,S\!\left(\frac{y-1}{2\varepsilon_z}\right)^{p_z}
\Biggr]_+^{1/4},
\label{eq:family1-final-rhotilde}
\end{equation}
for \(0\le x\le 1\), \(0\le y\le 1\), with the usual symmetry factor \(2\) in the
axial variable \(y\) and with the soft-wall switch function
\begin{equation}
S(\xi)=\frac{1+\tanh\xi}{2}.
\label{eq:family1-final-switch}
\end{equation}
The exact coupled mass and inertia integrals are then
\begin{equation}
I_2 = 2\int_0^1\!\!\int_0^1 x^2\widetilde\rho\,dx\,dy,
\qquad
I_4 = 2\int_0^1\!\!\int_0^1 x^4\widetilde\rho\,dx\,dy,
\qquad
I_w = 2\int_0^1\!\!\int_0^1 \frac{y^2}{4}x^2\widetilde\rho\,dx\,dy.
\label{eq:family1-final-I2I4Iw}
\end{equation}
Normalizing against the sharp-wall filled-to-endcap Thomas--Fermi baseline gives the
final coupled wall data
\begin{equation}
R_{\rm mass}^{\rm full}\approx 0.886313972989725,
\qquad
\widehat M_{aa}^{\rm full}\approx 0.563114968953987,
\qquad
\widehat M_{LL}^{\rm full}\approx 0.065829228119349.
\label{eq:family1-final-coupled-inertia}
\end{equation}
These are exactly the values carried into the final monopole breathing closure.

\subsection{Separated-order approximation was already very accurate}
\label{app:family1-final-separated-check}

It is useful to record one quantitative closure check.
The earlier separated-order completion had already given
\begin{equation}
R_{\rm mass}^{\rm sep}\approx 0.8846236634,
\qquad
\widehat M_{aa}^{\rm sep}\approx 0.5623810783,
\qquad
\widehat M_{LL}^{\rm sep}\approx 0.0671965962.
\label{eq:family1-final-separated-values}
\end{equation}
Relative to these values, the direct coupled full-profile corrections are only
\begin{equation}
\frac{R_{\rm mass}^{\rm full}-R_{\rm mass}^{\rm sep}}{R_{\rm mass}^{\rm sep}}
\approx +1.91\times 10^{-3},
\qquad
\frac{\widehat M_{aa}^{\rm full}-\widehat M_{aa}^{\rm sep}}{\widehat M_{aa}^{\rm sep}}
\approx +1.30\times 10^{-3},
\label{eq:family1-final-separated-mass-aa}
\end{equation}
\begin{equation}
\frac{\widehat M_{LL}^{\rm full}-\widehat M_{LL}^{\rm sep}}{\widehat M_{LL}^{\rm sep}}
\approx -2.03\times 10^{-2}.
\label{eq:family1-final-separated-LL}
\end{equation}
So the earlier separated reduction was already excellent for the overall mass scale and
radial inertia, and still reasonably good for the axial inertia.
That matters because it shows that the constructive program had essentially converged
before the final coupled check was performed.
The present coupled step is therefore a genuine tightening of the result, not a rescue of
an unstable earlier approximation.

\subsection{Final coupled monopole breathing response}
\label{app:family1-final-monopole}

Feeding the direct coupled wall data
\cref{eq:family1-final-coupled-inertia} into the same reduced geometry Hessian used in
\cref{app:geom-breathing} gives the final coupled monopole poles
\begin{equation}
\lambda_-\approx 6.405572392138922,
\qquad
\lambda_+\approx 254.444968136936126,
\label{eq:family1-final-lambdas}
\end{equation}
with positive residues
\begin{equation}
R_-\approx 0.002552474771738,
\qquad
R_+\approx 0.386733239513976,
\qquad
R_-+R_+=\frac{109}{280}.
\label{eq:family1-final-residues}
\end{equation}
So the final monopole wall channel is
\begin{equation}
K_{00}(s)
=
-\frac{757}{2520}
+\frac{R_-}{1-s/\lambda_-}
+\frac{R_+}{1-s/\lambda_+}.
\label{eq:family1-final-K00s}
\end{equation}
At zero frequency this reproduces the required static target exactly:
\begin{equation}
K_{00}(0)=\frac{4}{45}.
\label{eq:family1-final-K00zero}
\end{equation}
Relative to the sharp-wall Thomas--Fermi baseline, the physical monopole pole ratios are
\begin{equation}
\frac{\Omega_-^2}{\Omega_{-,\rm TF}^2}\approx 1.219665554412172,
\qquad
\frac{\Omega_+^2}{\Omega_{+,\rm TF}^2}\approx 1.206678094538845.
\label{eq:family1-final-pole-ratios}
\end{equation}
The corresponding one-pole Pad\'e reduction remains extremely accurate:
\begin{equation}
\lambda_{\rm eff}\approx 202.923516367519028,
\qquad
\max\,\mathrm{rel.err.}\approx 6.89\times 10^{-5}
\qquad
\text{on }0\le s\le 0.1\lambda_-.
\label{eq:family1-final-pade}
\end{equation}
So the historical single breathing auxiliary still survives as a very controlled
low-frequency reduction of the exact coupled two-pole monopole channel.

\subsection{Single-module summary and remaining open items}
\label{app:family1-final-summary}

The cleanest summary is now channel-theoretic rather than coefficient-theoretic.
At conservative \(\TwoPN\) order the constructive throat program has effectively reduced to
three coupled pieces:
\begin{equation}
\boxed{
\text{carried odd dipole wake}
\;\oplus\;
\text{exact }P_0\oplus P_2\text{ even support/source sector}
\;\oplus\;
\text{coupled two-pole monopole breathing channel }K_{00}(s).
}
\label{eq:family1-final-boxed-summary}
\end{equation}
The first two pieces reproduce the solved added comparable-mass conservative \(\TwoPN\)
cross block exactly at zero frequency.
The third piece closes the scalar monopole wall response with the best current dynamic
Family-1 normalization.
Together they form the final coupled Family-1 response module in the practical sense
relevant for the present paper.

It is also useful to state explicitly what remains open after this packaging.
Two unresolved items survive, and both are narrower than the original \(\TwoPN\) problem:
\begin{enumerate}[leftmargin=1.5em]
  \item the non-monopole dynamic pole scales of the odd and even
  \(\ell=1,2\) channels, which are true inner-throat PDE / \(\DtN\) observables rather
  than conservative zero-frequency necessities; and
  \item a fully analytic derivation of the local wall profiles from the microscopic
  soft-wall traction PDE, rather than from the reduced matching and full-profile
  evaluation used here.
\end{enumerate}
Neither of those open items blocks the conservative \(\TwoPN\) derivation itself.
What they block is only the next stage: a fully dynamical first-principles throat PDE
realization of the already-solved low-frequency operator.

So the practical finish-line statement of the constructive appendix program is now simply
this:
\begin{equation}
\boxed{
\text{the conservative }\TwoPN\text{ Family-1 throat-response module is effectively in hand.}
}
\label{eq:family1-final-practical-finish-line}
\end{equation}
That statement should be read in the same disciplined sense as the rest of the paper: not
as a claim that every dynamic PDE detail has been solved, but as the statement that the
full zero-frequency cross operator is closed exactly and that the scalar breathing branch
now has a direct coupled Family-1 wall completion.

\claimclass{\Controlled{} for the use of the specific balanced thin-layer Family-1 wall
profile and the reduced geometry breathing model.
\Exact{} for the unchanged zero-frequency odd/even algebraic reconstruction and for the
coupled full-profile integrals once that wall branch is declared.
\Protocol{} for the particular soft-wall / Thomas--Fermi reference branch used to package
the final dynamic monopole normalization.}


\end{document}
